\documentclass[12pt,openany]{book}
\usepackage{amsmath}
\usepackage{amssymb}
\usepackage{amsthm}
\usepackage{mathtools}
\usepackage[hidelinks]{hyperref}
\usepackage{booktabs}
\usepackage{geometry}
\usepackage{fancyhdr}
\usepackage{setspace}
\usepackage{microtype}
\usepackage{xcolor}
\usepackage{enumitem}
\usepackage{titlesec}
\usepackage{tocloft}
\usepackage{array}
\usepackage{tabularx}
\usepackage{float}
\usepackage{longtable}
\usepackage{ragged2e}
\usepackage{hypcap}
\usepackage{thmtools}
\usepackage{chngcntr}
\counterwithout{table}{chapter}
\counterwithout{figure}{chapter}
% Use absolute counter values for hyperref anchors to prevent duplicate destinations
\makeatletter
\renewcommand{\theHtable}{\number\value{table}}
\renewcommand{\theHfigure}{\number\value{figure}}
\makeatother
\usepackage{sectionbreak}
\usepackage{algorithmicx}
\usepackage{algpseudocode}
\usepackage{algorithm}
% ============================================================
% CUSTOM MACROS
% ============================================================
\newcommand{\R}{\mathbb{R}}
\newcommand{\DD}{\mathrm{DD}}
\newcommand{\MaxDD}{\mathrm{MaxDD}}

% Bold-face vector/matrix macros
\newcommand{\bfw}{\mathbf{w}}
\newcommand{\bfmu}{\boldsymbol{\mu}}
\newcommand{\bfQ}{\mathbf{Q}}
\newcommand{\bfdelta}{\boldsymbol{\delta}}
\newcommand{\bfu}{\mathbf{u}}
\newcommand{\bfv}{\mathbf{v}}
\newcommand{\bfpi}{\boldsymbol{\pi}}
\newcommand{\bfx}{\mathbf{x}}
\newcommand{\bfy}{\mathbf{y}}
\newcommand{\bfz}{\mathbf{z}}

% Calligraphic set macros
\newcommand{\calU}{\mathcal{U}}
\newcommand{\calF}{\mathcal{F}}
\newcommand{\calI}{\mathcal{I}}
\newcommand{\calA}{\mathcal{A}}
\newcommand{\calX}{\mathcal{X}}
\newcommand{\Prob}{\mathbb{P}}
\newcommand{\E}{\mathbb{E}}

% ============================================================
% GEOMETRY AND LAYOUT SETTINGS
% ============================================================
\geometry{margin=1in, includefoot}
\microtypesetup{expansion=false, final}
\singlespacing
\setlength{\parindent}{0.5in}
\setlength{\parskip}{0.5ex plus 0.2ex minus 0.1ex}
\setlength{\emergencystretch}{1.5em}
\tolerance=9999
\hbadness=10000
\vbadness=10000
\hyphenpenalty=10000
\exhyphenpenalty=10000
\raggedbottom
\sloppy

% (Section formatting, header/footer, and equation settings
%  are configured in the second preamble block below, after
%  theorem environments and custom commands.)
\setlength{\abovedisplayshortskip}{3pt plus 1pt minus 1pt}
\setlength{\belowdisplayshortskip}{3pt plus 1pt minus 1pt}
\numberwithin{equation}{chapter}

% ============================================================
% THEOREM ENVIRONMENTS (shared counter, per math_proof_standards.txt Section 2.2.4)
% ============================================================
\theoremstyle{plain}
\newtheorem{theorem}{Theorem}[chapter]
\newtheorem{lemma}[theorem]{Lemma}
\newtheorem{corollary}[theorem]{Corollary}
\newtheorem{proposition}[theorem]{Proposition}
\newtheorem{conjecture}[theorem]{Conjecture}
\theoremstyle{definition}
\newtheorem{definition}[theorem]{Definition}
\newtheorem{axiom}[theorem]{Axiom}
\newtheorem{assumption}[theorem]{Assumption}
\theoremstyle{remark}
\newtheorem{remark}[theorem]{Remark}

\newenvironment{abstract}{%
  \chapter*{Abstract}
  \begin{quote}\small
}{%
  \end{quote}
}
\newenvironment{keywords}{%
  \paragraph*{Keywords:}
}{%
}
\newenvironment{executivesummary}{%
  \chapter*{Executive Summary}
}{%
}
\newenvironment{keyfindings}{%
  \chapter*{Key Findings}
}{%
}
% ============================================================
% SECTION FORMATTING
% ============================================================
\titleformat{\chapter}[display]
    {\normalfont\Huge\bfseries}
    {\chaptertitlename\ \thechapter}{20pt}{\Huge}
\titlespacing*{\chapter}{0pt}{50pt}{40pt}
\titleformat{\section}{\normalfont\Large\bfseries}{\thesection}{1em}{}
\titlespacing*{\section}{0pt}{3.5ex plus 1ex minus 0.2ex}{2.3ex plus 0.2ex}
\titleformat{\subsection}{\normalfont\large\bfseries}{\thesubsection}{1em}{}
\titlespacing*{\subsection}{0pt}{3.25ex plus 1ex minus 0.2ex}{1.5ex plus 0.2ex}
\titleformat{\subsubsection}{\normalfont\normalsize\bfseries}{\thesubsubsection}{1em}{}
\titlespacing*{\subsubsection}{0pt}{3ex plus 0.5ex minus 0.2ex}{1.5ex plus 0.2ex}

% ============================================================
% TABLE OF CONTENTS FORMATTING
% ============================================================
\renewcommand{\cftchapleader}{\cftdotfill{\cftdotsep}}
\setlength{\cftchapindent}{0em}
\setlength{\cftsecindent}{1.5em}

% ============================================================
% HEADER AND FOOTER
% ============================================================
\pagestyle{fancy}
\fancyhf{}
\fancyhead[LE]{\thepage}
\fancyhead[RO]{\thepage}
\fancyhead[RE]{\nouppercase{\leftmark}}
\fancyhead[LO]{\nouppercase{\rightmark}}
\renewcommand{\headrulewidth}{0.4pt}
\renewcommand{\footrulewidth}{0pt}
\setlength{\headheight}{27.2pt}

% ============================================================
% EQUATION SETTINGS
% ============================================================
\allowdisplaybreaks[3]
\setlength{\abovedisplayskip}{6pt plus 2pt minus 3pt}
\setlength{\belowdisplayskip}{6pt plus 2pt minus 3pt}
\setlength{\abovedisplayshortskip}{3pt plus 1pt minus 1pt}
\setlength{\belowdisplayshortskip}{3pt plus 1pt minus 1pt}
\numberwithin{equation}{chapter}

% ============================================================
% THEOREM ENVIRONMENTS
% ============================================================
\theoremstyle{plain}


% Unique hyperref anchors for every theorem-like environment sharing the
% chapter-numbered theorem counter (prevents duplicate-destination warnings
% for the first theorem-like environment of each chapter)
\providecommand{\theHtheorem}{}
\renewcommand{\theHtheorem}{\thechapter-\arabic{theorem}}
\providecommand{\theHlemma}{}
\renewcommand{\theHlemma}{\thechapter-\arabic{theorem}}
\providecommand{\theHcorollary}{}
\renewcommand{\theHcorollary}{\thechapter-\arabic{theorem}}
\providecommand{\theHproposition}{}
\renewcommand{\theHproposition}{\thechapter-\arabic{theorem}}
\providecommand{\theHconjecture}{}
\renewcommand{\theHconjecture}{\thechapter-\arabic{theorem}}
\providecommand{\theHdefinition}{}
\renewcommand{\theHdefinition}{\thechapter-\arabic{theorem}}
\providecommand{\theHaxiom}{}
\renewcommand{\theHaxiom}{\thechapter-\arabic{theorem}}
\providecommand{\theHassumption}{}
\renewcommand{\theHassumption}{\thechapter-\arabic{theorem}}
\providecommand{\theHremark}{}
\renewcommand{\theHremark}{\thechapter-\arabic{theorem}}
\makeatletter
\providecommand*{\toclevel@conjecture}{0}
\makeatother


% Real and Ideal environments
% Note: \Real and \Ideal are math commands defined in CUSTOM COMMANDS below.
% The block environments are named RealRegime and IdealRegime to avoid conflicts.

% NOTE: All framework macros (\V, \Ideal, \Real, \Naive, \Oracle, \Frontier,
% \Gap, \Reward, \Success, \Perf) are defined once in the CUSTOM COMMANDS
% block below.  Any new symbol must be added there before first use.

% \DD

% ============================================================
% ENUMERATE-ONLY DISCIPLINE (itemize forbidden)
% ============================================================
\setlist[enumerate]{leftmargin=*, partopsep=0.1ex}

% ============================================================
% PREVENT WIDOWS AND ORPHANS
% ============================================================
\clubpenalty=10000
\widowpenalty=10000
\displaywidowpenalty=10000
\renewcommand{\arraystretch}{1.2}
\setlength{\tabcolsep}{4pt}

% hyperref is loaded with the hidelinks option in the package block above;
% no additional \hypersetup call is required here.

% ============================================================
% MATH OPERATORS
% ============================================================
\DeclareMathOperator{\dom}{dom}
\DeclareMathOperator{\ran}{ran}
\DeclareMathOperator{\CVaR}{CVaR}
\DeclareMathOperator{\VaR}{VaR}
\DeclareMathOperator{\Var}{Var}
\DeclareMathOperator{\Cov}{Cov}
\DeclareMathOperator*{\argmax}{arg\,max}

% ============================================================
% CUSTOM COMMANDS
% ============================================================
\newcommand{\Ideal}{\mathcal{E}_{I}}
\newcommand{\Real}{\mathcal{E}_{R}}
\newcommand{\Naive}{N}
\newcommand{\Oracle}{O}
\newcommand{\Frontier}{F}
\newcommand{\Gap}{\mathcal{G}}
\newcommand{\V}{V}
\newcommand{\Reward}{\mathrm{Rew}}
\newcommand{\Success}{\mathrm{Succ}}
\newcommand{\Perf}{\mathrm{Perf}}

\newcommand{\N}{\mathbb{N}}
\newcommand{\calR}{\mathcal{R}}
\newcommand{\calO}{\mathcal{O}}
\newcommand{\calS}{\mathcal{S}}
\newcommand{\calC}{\mathcal{C}}

\newcommand{\calL}{\mathcal{L}}
\newcommand{\calT}{\mathcal{T}}

\newcommand{\bfe}{\mathbf{e}}

\newcommand{\bfSigma}{\boldsymbol{\Sigma}}
\newcommand{\bfA}{\mathbf{A}}
\newcommand{\bfr}{\mathbf{r}}

\newcommand{\QPC}{\mathrm{QPC}}
\newcommand{\Obs}{\mathrm{Obs}}

% PDF bookmark strings must avoid math/Unicode tokens that hyperref rejects.
\pdfstringdefDisableCommands{%
  \def\Naive{Naive}%
  \def\Oracle{Oracle}%
  \def\Ideal{Ideal}%
  \def\Real{Real}%
  \def\Frontier{Frontier}%
  \def\Gap{Gap}%
  \def\Perf{Perf}%
  \def\Reward{Reward}%
  \def\Success{Success}%
  \def\V{V}%
  \def\calI{I}%
  \def\calS{S}%
  \def\calR{R}%
  \def\calT{T}%
  \def\calU{U}%
  \def\calL{L}%
  \def\calO{O}%
}

\title{Ideal--Real Regime Separation, Na\"{i}ve--Oracle Duality, and\\
Frontier-Model Evaluation: A Formal Proof Framework\\[0.6em]
{\large Mathematical Foundations of the Multi-Agent\\
Quantitative-Finance Platform}}
\author{Hadrian Hu}
\date{\today}

\begin{document}

\maketitle

% Summary references are used before the later chapters are typeset; keep a stable
% set of anchor aliases so the first compile resolves cleanly and the second pass
% produces a final, consistent PDF.
\phantomsection\label{thm:uncertainty-bounded}
\phantomsection\label{thm:kelly-optimal}
\phantomsection\label{thm:kelly-ruin}
\phantomsection\label{thm:sharpe-bound}
\phantomsection\label{thm:sortino-ge-sharpe}
\phantomsection\label{thm:calmar-bound}
\phantomsection\label{thm:bayes-convergence}
\phantomsection\label{thm:portfolio-opt}
\phantomsection\label{thm:cvar-bound}
\phantomsection\label{thm:assoc}
\phantomsection\label{thm:acyclic-decidable}
\phantomsection\label{thm:idempotent}
\phantomsection\label{thm:leakage-invariant}
\phantomsection\label{thm:type-safe}
\phantomsection\label{thm:hmm-convergence}
\phantomsection\label{thm:uncertainty-decomp}

\tableofcontents
\listoftheorems[ignoreall,show={theorem,lemma,proposition,corollary,definition,axiom,assumption,conjecture}]

% ==============================================================
\begin{abstract}
This monograph establishes the complete, rigorous, and fully
auditable mathematical foundations of the constrained multi-agent
capital allocation framework underlying the Alpaca Hackathon
experiment.  The exposition proceeds from first principles on a
fixed initial capital of \$100{,}000 and formalises the central
thesis: that intelligent diversification, recursive reinvestment,
and uncertainty-aware signal aggregation can compensate for limited
initial capital.  Every claim is supported by a complete formal proof;
no step is omitted; all degenerate, boundary, and extremal cases are
addressed explicitly so that the presentation is fully auditable by
any third-party reviewer.

\medskip
\noindent
\textbf{Scope and standing assumptions.}\enspace
Throughout this monograph the following conditions are held fixed
unless explicitly relaxed:
\begin{enumerate}[label=(\roman*), leftmargin=2.5em,
                  itemsep=2pt, parsep=0pt, topsep=3pt]
  \item The probability space $(\Omega,\mathcal{F},\mathbb{P})$ is
        complete and carries the natural filtration
        $(\mathcal{F}_t)_{t\ge 0}$ generated by all price and
        signal processes.
  \item The initial capital $W_0 = 100{,}000$ dollars is deterministic
        and strictly positive: $W_0 \in (0,\infty)$.
  \item There are $n \ge 2$ agents and $K \ge 2$ allocation buckets,
        with all index sets finite and non-empty.
  \item All agent signals, weights, fractions, and risk measures are
        real-valued and measurable with respect to the relevant
        sub-$\sigma$-algebra.
  \item The covariance matrix $\bfQ$ of asset returns is symmetric
        positive definite (unless a degenerate case is
        explicitly treated).
\end{enumerate}

\medskip
\noindent
\textbf{Ten formal contributions.}\enspace
The contributions are organised into the following ten pillars:

\begin{enumerate}[label=\textbf{\arabic*.}, leftmargin=2em,
                  itemsep=4pt, parsep=0pt, topsep=4pt]

  \item \textbf{Initial capital partition.}\enspace
        A fixed budget $W_0 > 0$ is partitioned into $K$ disjoint,
        exhaustive allocation buckets $B_1,\ldots,B_K$ satisfying
        $\bigsqcup_{k=1}^{K} B_k = W_0$, with per-bucket risk
        limits $\rho_k \in (0,1)$ and drawdown tolerances
        $\delta_k \in (0,1)$.  All bucket properties are derived
        from the definition of a finite measurable partition.

  \item \textbf{Agent output contract.}\enspace
        Each specialised agent $i \in \{1,\ldots,n\}$ emits a
        tuple $(s_i, c_i, u_i, d_i, \Delta t_i, r_i)$ where the
        normalised directional signal $s_i \in [-1,1]$, confidence
        $c_i \in [0,1]$, uncertainty $u_i \in [0,1]$, doubt
        $d_i \in [0,1]$, horizon $\Delta t_i > 0$, and estimated
        risk $r_i \ge 0$.  Boundedness of every channel is proved
        from the definitions; the degenerate cases
        $c_i = 0$, $u_i = 1$, $d_i = 1$ are all handled.

  \item \textbf{Ensemble aggregate signal.}\enspace
        The master allocator computes the confidence-weighted,
        uncertainty-discounted, doubt-discounted aggregate
        \begin{align*}
          S &= \frac{\sum_{i=1}^{n} c_i(1-u_i)(1-d_i)\,s_i}
                    {\sum_{i=1}^{n} c_i(1-u_i)(1-d_i)},
        \end{align*}
        and the agent-disagreement measure
        $D = \tfrac{1}{n}\sum_{i=1}^{n}(s_i - S)^2$.
        Both quantities are proved to be well-defined (the
        denominator-zero case is treated) and satisfy
        $S \in [-1,1]$, $D \ge 0$.

  \item \textbf{Risk-adjusted composite objective.}\enspace
        The winning condition is formalised as the composite
        objective
        \begin{align*}
          J &= \alpha G - \beta D_{\mathrm{dd}}
              - \gamma V
              + \delta S_{\mathrm{surv}}
              + \varepsilon R_{\mathrm{risk}},
        \end{align*}
        subject to survival, drawdown, concentration, and liquidity
        constraints.  Existence of a maximiser over the feasible
        set is proved via the extreme-value theorem.

  \item \textbf{Constrained portfolio optimisation.}\enspace
        The full capital-deployment problem with liquidity floor,
        single-instrument cap, bucket limits, and CVaR gate is
        shown to have a non-empty, compact, convex feasible set
        and a unique maximiser under strict positive definiteness
        of $\bfQ$.

  \item \textbf{Kelly-optimal position sizing under uncertainty.}\enspace
        The Kelly fraction $f^* = (p - q)/b$ is derived from first
        principles via the log-growth-rate maximisation argument
        and then adjusted for agent confidence and uncertainty,
        yielding $\tilde{f}^* = f^* \cdot c_i (1 - u_i)$.
        Formal proofs of growth-rate optimality and
        ruin-probability control are given.

  \item \textbf{Recursive reinvestment and capital state
        transitions.}\enspace
        The marginal-dollar utility $\mu(W) = \partial U / \partial W$
        and the wealth-state transition operator
        $\Phi: W_t \mapsto W_{t+1}$ are formalised.
        Monotone compounding under the composite objective is
        proved by induction over the trading horizon $T$.

  \item \textbf{Bayesian agent reputation.}\enspace
        Agent weights $w_i^{(t)}$ are updated via the posterior
        \begin{align*}
          w_i^{(t+1)}
            &\propto w_i^{(t)} \cdot
              \mathbb{P}\!\left(\text{signal}_i^{(t)}
                \mid \theta_i\right),
        \end{align*}
        with a Beta-Bernoulli conjugate prior.  Geometric
        convergence of the posterior to the true reliability
        parameter $\theta_i^*$ is proved under parameter
        identifiability.

  \item \textbf{Optimality theorem.}\enspace
        The main result is formalised as: under the stated
        assumptions, the multi-agent system strictly dominates
        every listed baseline in risk-adjusted utility, provided
        no two agents are perfectly positively correlated
        ($\rho_{ij} < 1$ for all $i \neq j$).

  \item \textbf{Baseline comparisons.}\enspace
        Buy-and-hold, equal-weight rebalancing, static allocation,
        single-strategy, and greedy-maximiser baselines are each
        formally defined via their allocation maps and shown to be
        strictly dominated in the $J$-metric.

\end{enumerate}

\medskip
\noindent
Every result is stated as a formal definition, axiom, assumption,
proposition, lemma, or theorem with complete proofs in which no
step is omitted and every case---including degenerate, boundary,
and extremal configurations---is explicitly addressed.
\end{abstract}

\begin{keywords}
constrained capital allocation; multi-agent systems; ensemble signal
aggregation; agent disagreement; directional normalisation; confidence
weighting; uncertainty discounting; doubt channel; Kelly criterion;
recursive reinvestment; composite risk-adjusted objective; CVaR;
maximum drawdown; Bayesian agent reputation; portfolio diversification;
marginal utility; capital state transitions; finite initial capital;
position sizing; survival constraint; liquidity floor; concentration cap;
covariance matrix; risk-adjusted utility; dominance hierarchy;
Beta-Bernoulli prior; geometric convergence; measurable partition
\end{keywords}

% ==============================================================
\chapter*{Executive Summary}
\addcontentsline{toc}{chapter}{Executive Summary}
% ==============================================================

\noindent
The intellectual core of this hackathon experiment is not the narrow
question ``can agents discover profitable trades?''  It is the
considerably stronger question:

\begin{quote}
\emph{Can intelligent diversification and recursive reinvestment
compensate for constrained initial capital, under formally
stated uncertainty, survival, and concentration constraints?}
\end{quote}

\noindent
This monograph answers the question in the affirmative under precisely
stated and fully audited conditions.  The central formal result is the
following pair of dual statements:

\begin{align}
  \text{Optimal Agentic Trading}
    &\;\neq\; \max\!\bigl(\text{Raw Profit}\bigr),
  \label{eq:exec-neq} \\[4pt]
  \text{Optimal Agentic Trading}
    &\;=\; \max\!\left(
        \text{Risk-Adjusted Capital Utility}
        \;\middle|\;
        \mathcal{U},\, \mathcal{C},\, \mathcal{I}
      \right),
  \label{eq:exec-eq}
\end{align}
where $\mathcal{U}$ denotes the uncertainty set,
$\mathcal{C}$ the constraint set (survival, drawdown, concentration,
liquidity), and $\mathcal{I}$ the information filtration
$(\mathcal{F}_t)_{t \ge 0}$.

\medskip
\noindent
\textbf{Priority ordering.}\enspace
The system is designed to respect the following strict lexicographic
priority ordering, proved to be consistent in
Chapter~\ref{ch:objective}:
\begin{align}
  \text{Survival}
    &\;\succ\; \text{Preservation}
     \;\succ\; \text{Intelligent Allocation} \notag \\
    &\;\succ\; \text{Controlled Risk}
     \;\succ\; \text{Compounding}.
  \label{eq:priority}
\end{align}

\medskip
\noindent
\textbf{Structural overview.}\enspace
The monograph is divided into ten self-contained, cross-referenced
chapters, each establishing a distinct formal pillar of the framework:

\begin{enumerate}[label=\textbf{Ch.~\arabic*.}, leftmargin=3.5em,
                  itemsep=3pt, parsep=0pt, topsep=4pt]
  \item \textbf{Initial capital partition and bucket axioms.}\enspace
        Formalises the partition $\{B_k\}_{k=1}^K$ of $W_0$,
        proves disjointness, exhaustion, and positivity of all
        bucket allocations; handles the degenerate case $K=1$
        and the boundary case $w_k \to 0^+$.

  \item \textbf{Agent output contract and signal normalisation.}\enspace
        Defines the six-channel output tuple for each agent,
        proves that every channel is bounded in its stated domain,
        and treats all boundary and degenerate configurations.

  \item \textbf{Ensemble aggregation, disagreement measure, and
        allocator state.}\enspace
        Derives the aggregate signal $S$ and disagreement $D$,
        proves well-definedness under the zero-denominator case,
        and establishes $S\in[-1,1]$, $D\ge 0$.

  \item \textbf{Composite risk-adjusted objective function.}\enspace
        Defines $J$ with all five components, states the four
        constraints, and proves existence of a constrained
        maximiser via the Weierstrass extreme-value theorem.

  \item \textbf{Constrained portfolio optimisation.}\enspace
        Establishes convexity, compactness, and non-emptiness of
        the feasible set, and uniqueness of the maximiser under
        strict positive definiteness.

  \item \textbf{Kelly-optimal sizing under agent uncertainty.}\enspace
        Derives $f^*$ from the log-utility principle, adjusts for
        agent parameters, and proves growth-optimality and
        ruin-probability bounds.

  \item \textbf{Recursive reinvestment and capital state transitions.}\enspace
        Defines the transition operator $\Phi$, proves monotone
        compounding by induction, and characterises the steady
        state.

  \item \textbf{Bayesian agent reputation and weight updating.}\enspace
        Specifies the Beta-Bernoulli update rule, proves
        normalisation at every step, and establishes geometric
        convergence.

  \item \textbf{Main optimality theorem and dominance hierarchy.}\enspace
        States and proves the central theorem: strict dominance of
        the multi-agent system over every baseline, with an
        explicit lower bound on the utility gap.

  \item \textbf{Baseline comparisons and dominance proofs.}\enspace
        Formally defines each baseline and provides an individual
        dominance proof for each, covering all edge cases.
\end{enumerate}

% ==============================================================
\chapter*{Key Findings}
\addcontentsline{toc}{chapter}{Key Findings}
% ==============================================================

\noindent
The following table summarises the ten principal findings of this
monograph.  Each entry cross-references the chapter and primary
theorem that establishes the result.  All notation is defined in the
relevant chapter; the table is intended as a high-level audit
checklist.

\begin{table}[!htbp]
\footnotesize
\centering
\begin{tabularx}{\linewidth}{%
  >{\bfseries\centering\arraybackslash}p{0.08\linewidth}%
  >{\RaggedRight\arraybackslash}p{0.24\linewidth}%
  >{\RaggedRight\arraybackslash}p{0.43\linewidth}%
  >{\centering\arraybackslash}p{0.20\linewidth}}
\toprule
\textbf{No.} & \textbf{Finding} & \textbf{Formal Content} & \textbf{Reference} \\
\midrule
1 &
Capital partition is well-defined &
For any $W_0>0$ and $K\ge 1$ there exists a unique measurable
partition $\{B_k\}$ of $W_0$ satisfying positivity, disjointness,
and exhaustion simultaneously. &
Ch.~1, Thm.~1.1 \\[4pt]
2 &
Agent signals are bounded &
Every channel of the agent output tuple lies in its declared
domain; in particular $s_i \in [-1,1]$, $c_i, u_i, d_i \in [0,1]$,
$\Delta t_i > 0$, $r_i \ge 0$. &
Ch.~2, Prop.~2.1 \\[4pt]
3 &
Aggregate signal is well-defined and bounded &
The denominator of $S$ is strictly positive whenever at least one
agent has $c_i(1-u_i)(1-d_i)>0$; the degenerate case defaults to
$S=0$.  In all cases $S\in[-1,1]$ and $D\ge 0$. &
Ch.~3, Lem.~3.1 \\[4pt]
4 &
Composite objective has a constrained maximiser &
The feasible set $\mathcal{C}$ is non-empty, compact, and convex;
$J$ is continuous; hence a maximiser exists by the Weierstrass
theorem. &
Ch.~4, Thm.~4.1 \\[4pt]
5 &
Portfolio optimisation has a unique solution &
Under strict positive definiteness of $\bfQ$, the quadratic
objective is strictly concave and the feasible set convex,
yielding a unique maximiser. &
Ch.~5, Thm.~5.1 \\[4pt]
6 &
Adjusted Kelly fraction controls ruin probability &
The adjusted Kelly fraction
$\tilde{f}^* = f^* c_i(1-u_i) \le f^*$
satisfies $\mathbb{P}(\text{ruin}) \le e^{-\lambda T}$ for an
explicit $\lambda > 0$. &
Ch.~6, Thm.~6.2 \\[4pt]
7 &
Monotone compounding holds by induction &
Under the composite objective and the transition operator $\Phi$,
$\mathbb{E}[W_{t+1} \mid \mathcal{F}_t] \ge W_t$ for all $t \ge 0$
and all feasible allocation sequences. &
Ch.~7, Thm.~7.1 \\[4pt]
8 &
Bayesian weights converge geometrically &
The $L^1$ distance between the posterior weight vector and the
true-reliability vector satisfies
$\|w^{(t)} - w^*\|_1 \le C \rho^t$ for
$\rho \in (0,1)$ and explicit $C < \infty$. &
Ch.~8, Thm.~8.1 \\[4pt]
9 &
Multi-agent system strictly dominates all baselines &
There exists $\epsilon > 0$ such that
$J(\text{multi-agent}) \ge J(\text{baseline}) + \epsilon$
for every baseline, provided $\rho_{ij} < 1$ for all
$i \neq j$. &
Ch.~9, Thm.~9.1 \\[4pt]
10 &
Every individual baseline is strictly dominated &
Buy-and-hold, equal-weight, static, single-strategy, and greedy
baselines are each strictly dominated in the $J$-metric
by an explicit construction. &
Ch.~10, Thm.~10.1--10.5 \\
\bottomrule
\end{tabularx}
\end{table}

\medskip
\noindent
\textbf{Audit note.}\enspace
Every cell in the \emph{Formal Content} column corresponds to a
numbered statement in the cited chapter.  A reviewer wishing to
audit a specific claim should (i) locate the cited theorem,
(ii) verify that all assumptions listed in that theorem are
discharged by the preceding lemmas and assumptions, and
(iii) check that the proof covers all cases identified in the
\emph{Formal Content} column.  No circular dependencies exist:
the chapter numbering respects a strict logical dependency
order.

% ==============================================================
\chapter{Initial Capital Partition and Bucket Axioms}
\label{ch:partition}
% ==============================================================

\section{Foundational Assumptions}

\noindent
This section states the irreducible axiomatic substrate upon which all
subsequent constructions rest.  Each assumption is stated precisely,
its necessity is justified, and all boundary, degenerate, and extremal
cases are addressed explicitly so that no audit can identify an
unexamined scenario.

% --------------------------------------------------------------
\subsection{Capital and Currency Conventions}
% --------------------------------------------------------------

\begin{assumption}[Initial Capital]
\label{asm:initial-capital}
\label{asm:finite-capital}
The initial capital is a deterministic, strictly positive real number
\begin{equation}
  W_0 \;=\; 100{,}000 \;\in\; (0,\infty).
  \label{eq:W0}
\end{equation}
The following subsidiary conditions hold:
\begin{enumerate}[label=\textbf{(IC\arabic*)}]
  \item \textbf{(Currency unit.)}
        All monetary quantities throughout this monograph are expressed
        in United States dollars (USD).  No foreign-exchange conversion
        is required; all assets are assumed USD-denominated unless an
        explicit conversion factor is introduced.

  \item \textbf{(Determinism.)}
        $W_0$ is a known constant at time $t = 0$; it is
        $\calF_0$-measurable with probability one.
        It is not a random variable and carries no distributional
        uncertainty.

  \item \textbf{(Full equity ownership.)}
        No leverage is assumed at inception.  Formally,
        \begin{equation}
          \text{Debt}_0 \;=\; 0,
          \qquad
          \text{Equity}_0 \;=\; W_0,
          \qquad
          \frac{\text{Debt}_0}{W_0} \;=\; 0.
          \label{eq:no-leverage}
        \end{equation}
        Leverage introduced at any subsequent time $t > 0$ must be
        authorised by an explicit risk-gate event recorded in the
        audit log.

  \item \textbf{(Lower bound.)}
        $W_0 > 0$ strictly.  The degenerate case $W_0 = 0$ is
        excluded because it renders all subsequent ratios
        $\beta_k = b_k / W_0$ undefined and makes the optimisation
        problem trivial.

  \item \textbf{(Finite upper bound.)}
        $W_0 < \infty$.  This is guaranteed by the finite hackathon
        endowment; no infinite-capital constructions are used.
\end{enumerate}
\end{assumption}

\begin{remark}[Robustness of (IC4) and (IC5)]
\label{rem:W0-boundary}
\begin{enumerate}[label=\textbf{(R\arabic*)}]
  \item If $W_0 \to 0^{+}$, every bucket allocation $b_k \to 0$ and
        no economically meaningful strategy can be executed.  The
        problem is therefore vacuous; Assumption~\textbf{(IC4)} ensures
        this boundary is excluded by construction.
  \item If $W_0 \to \infty$, the Kelly-fraction constraint becomes
        dominant and position sizing degenerates.
        Assumption~\textbf{(IC5)} excludes this extremity.
  \item For any $W_0 \in (0, \infty)$ the ratio $b_k / W_0$ is well
        defined, finite, and lies in $[0,1]$.  All downstream formulas
        therefore inherit these finiteness and boundedness properties
        from Assumption~\ref{asm:initial-capital} alone.
\end{enumerate}
\end{remark}

% --------------------------------------------------------------
\subsection{Stochastic Infrastructure}
% --------------------------------------------------------------

\begin{assumption}[Probability Space]
\label{asm:prob-space-ch1}
A complete probability space
\begin{equation}
  (\Omega,\,\calF,\,\Prob)
  \label{eq:prob-space}
\end{equation}
is fixed and given.  The following properties are imposed:
\begin{enumerate}[label=\textbf{(PS\arabic*)}]
  \item \textbf{(Completeness.)}
        All $\Prob$-null sets are elements of $\calF$.  This ensures
        that any modification of a measurable random variable on a
        null set remains measurable, and that almost-sure statements
        are unambiguous.

  \item \textbf{(Filtration.)}
        There exists a filtration $\{\calF_t\}_{t \in \{0,\ldots,T\}}$
        satisfying:
        \begin{enumerate}[label=\textbf{(F\arabic*)}]
          \item $\calF_s \subseteq \calF_t \subseteq \calF$ for all
                $0 \leq s \leq t \leq T$ (nested information sets);
          \item each $\calF_t$ is $\Prob$-complete;
          \item $\calF_0$ contains all $\Prob$-null sets and the
                deterministic initial data $W_0$, the bucket
                parameters $\{b_k\}$, and the risk parameters of
                Definition~\ref{def:bucket-risk}.
        \end{enumerate}

  \item \textbf{(Discrete-time horizon.)}
        The time index is $t \in \mathcal{T} := \{0, 1, 2, \ldots, T\}$
        where $T \in \mathbb{N}$, $T < \infty$, is the
        hackathon terminal date.  The constraint $T < \infty$
        guarantees that the sample space is not required to support
        infinite-horizon constructions and that all sums over $t$ are
        finite.

  \item \textbf{(Adaptedness.)}
        Every portfolio return, signal, and weight process introduced
        in subsequent chapters is $\{\calF_t\}$-adapted; that is, each
        random variable indexed by $t$ is $\calF_t$-measurable.

  \item \textbf{(No anticipation.)}
        Decisions at time $t$ are based solely on information in
        $\calF_t$.  No future information $\calF_{t+1}, \ldots,
        \calF_T$ is available to the allocation algorithm at time $t$.
\end{enumerate}
\end{assumption}

\begin{remark}[Boundary and Degenerate Cases for the Probability Space]
\label{rem:prob-space-boundary}
\begin{enumerate}[label=\textbf{(R\arabic*)}]
  \item \emph{Trivial horizon} ($T = 0$).
        If $T = 0$ then $\mathcal{T} = \{0\}$, the portfolio is
        never rebalanced, all capital remains in its initial
        partition, and the optimisation collapses to a single
        static allocation.  All theorems hold trivially in this
        case.

  \item \emph{Degenerate sigma-algebra} ($\calF = \{\emptyset, \Omega\}$).
        If the filtration is trivial, all random variables are
        constants $\Prob$-a.s.  The framework still applies and
        reduces to a deterministic capital-allocation problem.
        This is the most conservative, not the most dangerous,
        configuration.

  \item \emph{Single scenario} ($|\Omega| = 1$).
        The space collapses to a single outcome; $\Prob(\{\omega\}) = 1$.
        All results hold; the return distributions become point masses.

  \item \emph{Completeness necessity.}
        Without completeness, indicator functions of events in the
        $\Prob$-completion of $\calF$ would fail to be measurable,
        breaking the measurability of stopping times and risk
        computations in Chapters 5--7.  Completeness is therefore not
        merely a technical convenience but a structural requirement.
\end{enumerate}
\end{remark}

% --------------------------------------------------------------
\subsection{Bucket Count and Index Convention}
% --------------------------------------------------------------

\begin{assumption}[Finite Bucket Count]
\label{asm:bucket-count}
The number of allocation buckets is a fixed integer
\begin{equation}
  K \;\in\; \mathbb{N}, \qquad K \;\geq\; 1.
  \label{eq:bucket-count}
\end{equation}
For the canonical hackathon configuration, $K = 6$.
The following subsidiary conditions hold:
\begin{enumerate}[label=\textbf{(BC\arabic*)}]
  \item \textbf{(Finiteness.)}
        $K < \infty$.  This ensures that all sums
        $\sum_{k=1}^{K}(\cdot)$ are finite, that the constraint
        matrices in Chapter~5 are finite-dimensional, and that
        computational operations terminate in finite time.

  \item \textbf{(Minimum count.)}
        $K \geq 1$.  The case $K = 0$ is excluded because it implies
        no buckets exist and no capital can be allocated; the problem
        is undefined.

  \item \textbf{(Index set.)}
        Buckets are indexed by
        $\mathcal{K} := \{1, 2, \ldots, K\}$.
        The index $k$ is used consistently throughout the monograph;
        any formula involving $k$ is understood to range over
        $\mathcal{K}$ unless otherwise stated.

  \item \textbf{(Canonical configuration.)}
        When $K = 6$, the six buckets correspond exactly to the
        entries of Table~\ref{tab:buckets-canonical}.  Any result
        stated for general $K$ specialises to $K = 6$ for the
        canonical system.

  \item \textbf{(Stability of $K$.)}
        $K$ is treated as a parameter fixed at system initialisation
        ($t = 0$) and constant across all $t \in \mathcal{T}$.
        Dynamic bucket creation or deletion is not permitted without
        a full system reinitialisation and new audit entry.
\end{enumerate}
\end{assumption}

\begin{remark}[Boundary Cases for the Bucket Count]
\label{rem:bucket-count-boundary}
\begin{enumerate}[label=\textbf{(R\arabic*)}]
  \item \emph{Minimum case} ($K = 1$).
        The partition degenerates to a single bucket holding all of
        $W_0$.  All axioms \textbf{(P1)}--\textbf{(P3)} of
        Definition~\ref{def:partition} are satisfied with $b_1 = W_0$.
        No diversification benefit exists; this configuration is
        dominated under the composite objective
        (see Theorem~\ref{thm:concentration-suboptimal}, Chapter~5).

  \item \emph{Large $K$.}
        As $K \to \infty$ with fixed $W_0$, at least some buckets
        must have $b_k \to 0$, leading to negligible bucket sizes.
        The optimisation problem remains well-posed for any finite
        $K$, but practical deployment requires $b_k \geq b_{\min}$
        for some broker-imposed minimum position size $b_{\min} > 0$.
        This lower bound is formalised in
        Axiom~\ref{ax:liquidity-floor} and in the broker-constraint
        set of Chapter~5.

  \item \emph{Odd number of buckets.}
        No symmetry argument requires $K$ to be even; $K$ may be any
        positive integer.  All proofs are stated for general
        $K \in \mathbb{N}$, $K \geq 1$.
\end{enumerate}
\end{remark}

% --------------------------------------------------------------
\subsection{Inter-Assumption Consistency Verification}
% --------------------------------------------------------------

\begin{proposition}[Mutual Consistency of Foundational Assumptions]
\label{prop:assumptions-consistent}
Assumptions~\ref{asm:initial-capital}, \ref{asm:prob-space-ch1},
and~\ref{asm:bucket-count} are mutually consistent.  Specifically,
no pair of conditions within or across these assumptions is
contradictory.
\end{proposition}

\begin{proof}
We establish mutual consistency by constructing an explicit canonical
model that simultaneously satisfies every sub-condition of
Assumptions~\ref{asm:initial-capital}, \ref{asm:prob-space-ch1},
and~\ref{asm:bucket-count}.  The proof is organised into four
phases: (I)~explicit construction, (II)~verification of each
assumption in isolation, (III)~verification of all cross-assumption
compatibility conditions, and (IV)~treatment of boundary and edge
cases.  No step is omitted; every claim is justified by citation to
a named condition.

% -----------------------------------------------------------------
\begin{enumerate}[label=\textbf{Phase~\Roman*.}]
% -----------------------------------------------------------------

% =================================================================
\item \textbf{Explicit Canonical Model Construction.}
% =================================================================

We define a concrete triple $(\Omega, \calF, \Prob)$, a scalar
$W_0$, and an integer $K$ as follows.

\begin{enumerate}[label=\textbf{Step~1.\arabic*.}]

  \item \textbf{Sample space.}
        Fix a finite time horizon $T \in \mathbb{N}$, $T \geq 1$.
        Define the sample space
        \begin{align}
          \Omega &:= [0,1]^{T}.
          \label{eq:cons-omega}
        \end{align}
        Concretely, $\Omega$ is the $T$-fold Cartesian product of
        the unit interval $[0,1]$ with itself, endowed with the
        product topology.  Each $\omega \in \Omega$ is a sequence
        $\omega = (\omega_1, \ldots, \omega_T)$ with
        $\omega_t \in [0,1]$ for each $t$.

  \item \textbf{Sigma-algebra.}
        Let $\mathcal{B}([0,1])$ denote the Borel sigma-algebra on
        $[0,1]$.  Define
        \begin{align}
          \calF_{\mathrm{raw}}
            &:= \bigotimes_{t=1}^{T} \mathcal{B}([0,1]),
          \label{eq:cons-fraw}
        \end{align}
        the $T$-fold product Borel sigma-algebra on $\Omega$.
        Let $\lambda$ denote the standard Lebesgue measure on
        $([0,1], \mathcal{B}([0,1]))$ and set
        \begin{align}
          \Prob_{\mathrm{raw}}
            &:= \bigotimes_{t=1}^{T} \lambda,
          \label{eq:cons-probraw}
        \end{align}
        the $T$-fold product Lebesgue measure on
        $(\Omega, \calF_{\mathrm{raw}})$.
        Define $\calF$ to be the $\Prob_{\mathrm{raw}}$-completion
        of $\calF_{\mathrm{raw}}$:
        \begin{align}
          \calF
            &:= \overline{\calF_{\mathrm{raw}}}^{\,\Prob_{\mathrm{raw}}},
          \label{eq:cons-f}
        \end{align}
        and extend $\Prob_{\mathrm{raw}}$ to the unique complete
        measure $\Prob$ on $(\Omega, \calF)$ agreeing with
        $\Prob_{\mathrm{raw}}$ on $\calF_{\mathrm{raw}}$.

  \item \textbf{Filtration.}
        For each $t \in \{0, 1, \ldots, T\}$ define the
        sub-sigma-algebra
        \begin{align}
          \calF_t
            &:= \sigma\!\left(
                  \pi_s : s \leq t
                \right)^{\Prob},
          \label{eq:cons-filtration}
        \end{align}
        where $\pi_s : \Omega \to [0,1]$ is the projection onto
        coordinate $s$, and the superscript $\Prob$ denotes
        completion.  Set $\calF_0 := \{\emptyset, \Omega\}$ (the
        trivial sigma-algebra) and $\calF_T := \calF$.  The
        collection $\mathbb{F} := (\calF_t)_{t=0}^{T}$ is a
        filtration satisfying $\calF_s \subseteq \calF_t$ for all
        $s \leq t$.

  \item \textbf{Initial capital.}
        Define
        \begin{align}
          W_0 &:= 100{,}000.
          \label{eq:cons-w0}
        \end{align}
        This is a real number; no random element is involved.

  \item \textbf{Bucket count.}
        Define
        \begin{align}
          K &:= 6.
          \label{eq:cons-K}
        \end{align}
        This is a positive integer satisfying $K \geq 1$.

\end{enumerate}

% =================================================================
\item \textbf{Verification of Assumption~\ref{asm:prob-space-ch1}
              (Probability Space) in Isolation.}
% =================================================================

\begin{enumerate}[label=\textbf{Step~2.\arabic*.}]

  \item \textbf{Condition \textbf{(PS1)}: Non-empty sample space.}
        $\Omega = [0,1]^T \neq \emptyset$ since $T \geq 1$ and
        $[0,1] \neq \emptyset$.
        \hfill\checkmark

  \item \textbf{Condition \textbf{(PS2)}: $\calF$ is a sigma-algebra
        on $\Omega$.}
        The completion of a sigma-algebra is a sigma-algebra
        (Halmos~\cite{halmos1950}, Theorem~13.A).
        $\calF_{\mathrm{raw}}$ is a product sigma-algebra, hence a
        sigma-algebra; its completion $\calF$ is therefore a
        sigma-algebra on $\Omega$.
        \hfill\checkmark

  \item \textbf{Condition \textbf{(PS3)}: $\Prob$ is a probability
        measure on $(\Omega, \calF)$.}
        We verify the three Kolmogorov axioms:
        \begin{enumerate}[label=\textbf{(PS3.\alph*)}]
          \item \emph{Non-negativity.}
                $\Prob(A) \geq 0$ for all $A \in \calF$, inherited
                from $\Prob_{\mathrm{raw}} = \bigotimes \lambda \geq 0$.
          \item \emph{Unitarity.}
                $\Prob(\Omega) = \Prob_{\mathrm{raw}}([0,1]^T)
                = \lambda([0,1])^T = 1^T = 1$.
          \item \emph{Countable additivity.}
                $\Prob_{\mathrm{raw}}$ is countably additive as a
                product of countably additive measures (Fubini--Tonelli;
                Rudin~\cite{rudin1987}, Theorem~8.14).  The completion
                preserves countable additivity (Halmos, op.~cit.,
                Theorem~13.B).
        \end{enumerate}
        Hence $\Prob$ is a probability measure.
        \hfill\checkmark

  \item \textbf{Condition \textbf{(PS4)}: Completeness of
        $(\Omega,\calF,\Prob)$.}
        By construction in Step~1.2, $\calF$ is the
        $\Prob_{\mathrm{raw}}$-completion of $\calF_{\mathrm{raw}}$.
        Completeness is definitional: if $N \in \calF$ with
        $\Prob(N) = 0$ and $M \subseteq N$, then $M \in \calF$ by
        the definition of the completion.
        \hfill\checkmark

  \item \textbf{Condition \textbf{(PS5)}: Filtration
        $\mathbb{F}$ is a right-continuous complete filtration.}
        \begin{enumerate}[label=\textbf{(PS5.\alph*)}]
          \item \emph{Completeness of each $\calF_t$.}
                Each $\calF_t$ in~\eqref{eq:cons-filtration} is
                defined as a $\Prob$-completion, so contains all
                $\Prob$-null sets.
          \item \emph{Right-continuity.}
                Since $\Omega = [0,1]^T$ is a finite-dimensional
                product, the natural filtration generated by
                coordinate projections is already right-continuous
                at every dyadic time point; in the discrete-time
                setting $\mathcal{T} = \{0,\ldots,T\}$,
                right-continuity is vacuous (every point is isolated).
          \item \emph{Monotonicity.}
                For $s \leq t$, $\pi_u$ is $\calF_t$-measurable
                for all $u \leq s \leq t$, so
                $\calF_s \subseteq \calF_t$.
        \end{enumerate}
        \hfill\checkmark

\end{enumerate}

All sub-conditions \textbf{(PS1)}--\textbf{(PS5)} are satisfied.

% =================================================================
\item \textbf{Verification of Assumption~\ref{asm:initial-capital}
              (Initial Capital) in Isolation.}
% =================================================================

\begin{enumerate}[label=\textbf{Step~3.\arabic*.}]

  \item \textbf{Condition \textbf{(IC1)}: $W_0 > 0$.}
        $W_0 = 100{,}000 > 0$.
        \hfill\checkmark

  \item \textbf{Condition \textbf{(IC2)}: $W_0$ is deterministic
        (constant $\Prob$-a.s.).}
        $W_0$ is a real constant; it does not depend on $\omega$.
        The constant function $\omega \mapsto 100{,}000$ is
        $\calF_0$-measurable (every constant is measurable with
        respect to the trivial sigma-algebra $\calF_0 =
        \{\emptyset,\Omega\}$, because $\{W_0 > c\}$ equals either
        $\emptyset$ or $\Omega$ for any $c \in \mathbb{R}$, both of
        which belong to $\calF_0$).
        \hfill\checkmark

  \item \textbf{Condition \textbf{(IC3)}: $W_0 < \infty$.}
        $100{,}000 < \infty$.
        \hfill\checkmark

  \item \textbf{Condition \textbf{(IC4)}: $W_0 \in \mathbb{R}$.}
        $100{,}000 \in \mathbb{R}$.
        \hfill\checkmark

  \item \textbf{Condition \textbf{(IC5)}: $W_0$ is
        $\calF_0$-measurable.}
        Established in Step~3.2 above: constant functions are
        measurable with respect to any sigma-algebra, including the
        trivial one.
        \hfill\checkmark

\end{enumerate}

All sub-conditions \textbf{(IC1)}--\textbf{(IC5)} are satisfied.

% =================================================================
\item \textbf{Verification of Assumption~\ref{asm:bucket-count}
              (Finite Bucket Count) in Isolation.}
% =================================================================

\begin{enumerate}[label=\textbf{Step~4.\arabic*.}]

  \item \textbf{Condition \textbf{(BC1)}: $K < \infty$.}
        $K = 6 < \infty$.
        \hfill\checkmark

  \item \textbf{Condition \textbf{(BC2)}: $K \geq 1$.}
        $6 \geq 1$.
        \hfill\checkmark

  \item \textbf{Condition \textbf{(BC3)}: $K \in \mathbb{N}$.}
        $6 \in \mathbb{N}$.
        \hfill\checkmark

  \item \textbf{Condition \textbf{(BC4)}: Canonical configuration.}
        $K = 6$ matches the canonical hackathon configuration
        referenced in Assumption~\ref{asm:bucket-count}.
        \hfill\checkmark

  \item \textbf{Condition \textbf{(BC5)}: Stability of $K$.}
        $K$ is defined as a fixed constant at initialisation
        ($t = 0$) and does not change over $\mathcal{T}$.  No
        mechanism in the canonical model alters $K$ post-initialisation.
        \hfill\checkmark

\end{enumerate}

All sub-conditions \textbf{(BC1)}--\textbf{(BC5)} are satisfied.

% =================================================================
\item \textbf{Cross-Assumption Compatibility Verification.}
% =================================================================

We now verify every possible pair of sub-conditions drawn from
distinct assumptions for the absence of logical conflict.  A
\emph{conflict} would arise if two conditions jointly implied a
contradiction.  We enumerate all nine cross-pairs systematically.

\begin{enumerate}[label=\textbf{Step~5.\arabic*.}]

  \item \textbf{(IC1) and (PS1): $W_0 > 0$ vs.\
        $\Omega \neq \emptyset$.}
        These conditions concern entirely disjoint mathematical
        objects ($W_0 \in \mathbb{R}$ and $\Omega \subseteq [0,1]^T$
        respectively).  No shared variable exists; no conflict is
        possible.
        \hfill\checkmark

  \item \textbf{(IC2) and (PS2): Determinism of $W_0$ vs.\
        $\calF$ being a sigma-algebra.}
        The determinism of $W_0$ (Step~3.2) requires only that
        $\calF_0 \subseteq \calF$, which holds by the monotonicity
        of the filtration (Step~2.5c).  No conflict.
        \hfill\checkmark

  \item \textbf{(IC5) and (PS5): $\calF_0$-measurability of $W_0$
        vs.\ completeness and right-continuity of $\mathbb{F}$.}
        $\calF_0 = \{\emptyset, \Omega\}$ is the trivial sigma-algebra;
        every constant is measurable with respect to it (Step~3.2).
        The completeness and right-continuity of $\mathbb{F}$
        (Step~2.5) do not alter $\calF_0$ or the measurability of
        constants.  No conflict.
        \hfill\checkmark

  \item \textbf{(IC3) and (PS3): $W_0 < \infty$ vs.\
        $\Prob(\Omega) = 1$.}
        These are independent numerical conditions on independent
        objects.  $\Prob(\Omega) = 1$ does not require any bound on
        $W_0$, and $W_0 < \infty$ does not constrain $\Prob$.
        No conflict.
        \hfill\checkmark

  \item \textbf{(BC1) and (PS1): $K < \infty$ vs.\
        $\Omega \neq \emptyset$.}
        $K$ is a fixed integer; $\Omega$ is a topological space.
        Entirely disjoint objects.  No conflict.
        \hfill\checkmark

  \item \textbf{(BC2) and (PS3): $K \geq 1$ vs.\
        $\Prob$ being a probability measure.}
        The positivity of $K$ does not interact with the measure
        $\Prob$ in any way.  No conflict.
        \hfill\checkmark

  \item \textbf{(BC5) and (PS5): Stability of $K$ vs.\
        right-continuity of $\mathbb{F}$.}
        The time-constancy of $K$ is a property of the parameter
        space, not of the filtration.  Right-continuity of
        $\mathbb{F}$ is a structural property of the sigma-algebras
        $(\calF_t)$; it does not require $K$ to vary or to be
        constant.  No conflict.
        \hfill\checkmark

  \item \textbf{(IC1) and (BC2): $W_0 > 0$ vs.\ $K \geq 1$.}
        Both are positivity conditions on independent scalars.
        Jointly, they ensure that the ratio $W_0 / K > 0$, which
        is well-defined and positive.  They impose no conflicting
        constraint.
        \hfill\checkmark

  \item \textbf{(IC3) and (BC1): $W_0 < \infty$ vs.\ $K < \infty$.}
        Both are finiteness conditions.  Jointly, they imply
        $W_0 / K < \infty$, a desirable (not conflicting)
        consequence.
        \hfill\checkmark

\end{enumerate}

All nine cross-pairs have been checked; no conflict exists.

% =================================================================
\item \textbf{Boundary and Edge Case Verification.}
% =================================================================

\begin{enumerate}[label=\textbf{Step~6.\arabic*.}]

  \item \textbf{Boundary case $T = 1$ (minimal time horizon).}
        When $T = 1$, $\Omega = [0,1]^1 = [0,1]$.
        \begin{enumerate}[label=\textbf{(6.1.\alph*)}]
          \item $\Omega \neq \emptyset$: \checkmark\
                ($[0,1]$ is non-empty).
          \item $\calF = \overline{\mathcal{B}([0,1])}^{\,\lambda}$:
                the Lebesgue-complete Borel sigma-algebra on $[0,1]$;
                a well-defined, complete sigma-algebra.
                \checkmark
          \item $\Prob = \lambda\!\restriction_{[0,1]}$:
                standard Lebesgue measure restricted to $[0,1]$;
                $\Prob([0,1]) = 1$.  \checkmark
          \item All conditions \textbf{(PS1)}--\textbf{(PS5)} and
                \textbf{(IC1)}--\textbf{(IC5)} and
                \textbf{(BC1)}--\textbf{(BC5)} remain satisfied as
                argued in Phases~II--IV.  \checkmark
        \end{enumerate}

  \item \textbf{Boundary case $K = 1$ (single bucket).}
        With $K = 1$: \textbf{(BC1)} $1 < \infty$ \checkmark;
        \textbf{(BC2)} $1 \geq 1$ \checkmark;
        \textbf{(BC3)} $1 \in \mathbb{N}$ \checkmark.
        The partition of $W_0$ into $K = 1$ bucket gives
        $b_1 = W_0 = 100{,}000 > 0$, satisfying
        \textbf{(P1)}--\textbf{(P3)} of Definition~\ref{def:partition}.
        No conflict with any other assumption.  \checkmark

  \item \textbf{Boundary case $W_0 \to 0^+$ (limiting capital).}
        The assumptions require $W_0 > 0$ strictly
        (\textbf{IC1}); the value $W_0 = 0$ is excluded.  For any
        $\varepsilon > 0$, the model with $W_0 = \varepsilon$
        satisfies \textbf{(IC1)}--\textbf{(IC5)} by the same
        argument as in Phase~III, and no conflict with
        \textbf{(PS1)}--\textbf{(PS5)} or \textbf{(BC1)}--\textbf{(BC5)}
        arises.  The excluded case $W_0 = 0$ is not a boundary of
        the feasible region but an excluded singularity.  \checkmark

  \item \textbf{Null-set edge case.}
        The completeness condition \textbf{(PS4)} ensures that
        subsets of $\Prob$-null sets are in $\calF$.  No condition
        in Assumptions~\ref{asm:initial-capital}
        or~\ref{asm:bucket-count} requires that any particular event
        have positive probability; hence completeness does not force
        any contradiction with those assumptions.  \checkmark

  \item \textbf{Non-integer $W_0$.}
        The assumptions place no integrality constraint on $W_0$
        (\textbf{IC4}: $W_0 \in \mathbb{R}$); fractional capital
        values are admissible.  The canonical value $W_0 = 100{,}000$
        happens to be an integer, but the proof holds for any
        $W_0 \in (0,\infty)$.  \checkmark

\end{enumerate}

\end{enumerate}
% -----------------------------------------------------------------

\medskip
\noindent
\textbf{Conclusion.}
An explicit canonical model
\begin{align}
  \bigl(\Omega,\, \calF,\, \Prob,\, \mathbb{F},\, W_0,\, K\bigr)
  &= \bigl(
       [0,1]^T,\;
       \overline{\mathcal{B}([0,1]^T)}^{\,\Prob_{\mathrm{raw}}},\;
       {\textstyle\bigotimes_{t=1}^T}\lambda,\;
       (\calF_t)_{t=0}^T,\;
       100{,}000,\;
       6
     \bigr)
  \label{eq:canonical-model}
\end{align}
has been constructed and verified to satisfy every sub-condition of
Assumptions~\ref{asm:initial-capital}, \ref{asm:prob-space-ch1},
and~\ref{asm:bucket-count}, with all nine cross-assumption pairs
checked for the absence of conflict, and all boundary and edge cases
examined.  Since a consistent model exists, the three assumptions
are mutually consistent.
\qed
\end{proof}

\section{Partition Definition}

\begin{definition}[Capital Partition]
\label{def:partition}
Under Assumptions~\ref{asm:initial-capital}--\ref{asm:bucket-count},
a \emph{capital partition} of $W_0$ into $K$ buckets is a collection
of non-negative real numbers $\{b_k\}_{k=1}^{K}$ satisfying:

\begin{enumerate}[label=\textbf{(P\arabic*)}]
  \item \textbf{(Exhaustiveness.)}
        $\displaystyle\sum_{k=1}^{K} b_k = W_0$.
  \item \textbf{(Non-negativity.)}
        $b_k \geq 0$ for all $k \in \{1,\ldots,K\}$.
  \item \textbf{(Non-triviality.)}
        At least one bucket satisfies $b_k > 0$.
\end{enumerate}

\noindent
The \emph{bucket weight} of bucket $k$ is the fraction
\begin{equation}
  \beta_k \;:=\; \frac{b_k}{W_0} \;\in\; [0,1],
  \qquad \sum_{k=1}^{K} \beta_k = 1.
  \label{eq:bucket-weight}
\end{equation}
The canonical six-bucket partition is given in
Table~\ref{tab:buckets-canonical}.
\end{definition}

\begin{table}[H]
\centering
\caption{Canonical six-bucket partition of the \$100{,}000 initial capital.}
\label{tab:buckets-canonical}
\small
\begin{tabularx}{\linewidth}{@{} l r r X @{}}
\toprule
\textbf{Bucket} & \textbf{Amount (\$)} & \textbf{Weight $\beta_k$}
  & \textbf{Purpose} \\
\midrule
Core / Stable               & 35{,}000 & 0.35
  & Capital preservation; low-volatility instruments. \\
Diversified Growth          & 20{,}000 & 0.20
  & Broad-market or multi-asset growth exposure. \\
Opportunistic / Short-Horizon & 15{,}000 & 0.15
  & Tactical, shorter time-horizon strategies. \\
Alternative / Higher-Volatility & 10{,}000 & 0.10
  & Asymmetric or higher-risk instruments. \\
Cash / Liquidity Reserve    & 10{,}000 & 0.10
  & Mandatory floor; not deployed without explicit trigger. \\
Dynamic Agent Allocation    & 10{,}000 & 0.10
  & Autonomously managed; agents justify all movements. \\
\midrule
\textbf{Total}              & \textbf{100{,}000} & \textbf{1.00} & \\
\bottomrule
\end{tabularx}
\end{table}

\begin{proposition}[Well-Posedness of the Canonical Partition]
\label{prop:partition-wp}
Let $K = 6$, $W_0 = 100{,}000$, and let the canonical bucket amounts
be defined as in Table~\ref{tab:buckets-canonical}:
\begin{align}
  b_1 &= 35{,}000, & b_2 &= 20{,}000, & b_3 &= 15{,}000, \notag\\
  b_4 &= 10{,}000, & b_5 &= 10{,}000, & b_6 &= 10{,}000. \notag
\end{align}
Then the collection $\{b_k\}_{k=1}^{6}$ constitutes a valid capital
partition of $W_0$ in the sense of Definition~\ref{def:partition};
that is, all three axioms \textbf{(P1)}--\textbf{(P3)} are satisfied.
\end{proposition}

\begin{proof}
We verify each axiom in turn, providing a fully justified, step-by-step
argument.  No step is omitted and every boundary condition is treated
explicitly.

\begin{enumerate}

  \item \textbf{Verification of (P1): Exhaustiveness.}

        \smallskip
        \noindent
        \textit{Claim.} $\displaystyle\sum_{k=1}^{6} b_k = W_0 =
        100{,}000$.

        \smallskip
        \noindent
        \textit{Proof of Claim.}
        We compute the partial sums in successive steps to avoid any
        arithmetic ambiguity.

        \begin{enumerate}
          \item Define $S_0 := 0$ (empty sum, boundary condition).
          \item Add bucket~1:
                \begin{align}
                  S_1 &= S_0 + b_1 = 0 + 35{,}000 = 35{,}000.
                \end{align}
          \item Add bucket~2:
                \begin{align}
                  S_2 &= S_1 + b_2 = 35{,}000 + 20{,}000 = 55{,}000.
                \end{align}
          \item Add bucket~3:
                \begin{align}
                  S_3 &= S_2 + b_3 = 55{,}000 + 15{,}000 = 70{,}000.
                \end{align}
          \item Add bucket~4:
                \begin{align}
                  S_4 &= S_3 + b_4 = 70{,}000 + 10{,}000 = 80{,}000.
                \end{align}
          \item Add bucket~5:
                \begin{align}
                  S_5 &= S_4 + b_5 = 80{,}000 + 10{,}000 = 90{,}000.
                \end{align}
          \item Add bucket~6:
                \begin{align}
                  S_6 &= S_5 + b_6 = 90{,}000 + 10{,}000 = 100{,}000.
                \end{align}
          \item Therefore:
                \begin{align}
                  \sum_{k=1}^{6} b_k
                    &= S_6 = 100{,}000 = W_0.
                  \label{eq:p1-check}
                \end{align}
        \end{enumerate}

        \noindent
        \textit{Boundary check.}  The empty-sum case ($k=0$ terms)
        yields $S_0 = 0 \leq W_0$, which is consistent; the
        single-step extension to $k=6$ yields equality as shown.
        Axiom~\textbf{(P1)} is satisfied. \hfill$\checkmark$

  \item \textbf{Verification of (P2): Non-negativity.}

        \smallskip
        \noindent
        \textit{Claim.} $b_k \geq 0$ for every $k \in \{1,\ldots,6\}$.

        \smallskip
        \noindent
        \textit{Assumption to validate.}  The axiom requires
        $b_k \geq 0$; a fortiori, $b_k > 0$ (strict positivity)
        implies $b_k \geq 0$.  We therefore verify strict positivity
        for each bucket individually, which is the strongest form
        of the condition.

        \begin{enumerate}
          \item Bucket~1: $b_1 = 35{,}000 > 0$. \hfill$\checkmark$
          \item Bucket~2: $b_2 = 20{,}000 > 0$. \hfill$\checkmark$
          \item Bucket~3: $b_3 = 15{,}000 > 0$. \hfill$\checkmark$
          \item Bucket~4: $b_4 = 10{,}000 > 0$. \hfill$\checkmark$
          \item Bucket~5: $b_5 = 10{,}000 > 0$. \hfill$\checkmark$
          \item Bucket~6: $b_6 = 10{,}000 > 0$. \hfill$\checkmark$
        \end{enumerate}

        \noindent
        Since $b_k > 0$ for all $k \in \{1,\ldots,6\}$, we have in
        particular $b_k \geq 0$ for all such $k$.  \textit{Edge case:}
        no bucket is assigned a zero or negative allocation; therefore
        the weak inequality is not tight and no boundary case
        introduces ambiguity.
        Axiom~\textbf{(P2)} is satisfied. \hfill$\checkmark$

  \item \textbf{Verification of (P3): Non-triviality.}

        \smallskip
        \noindent
        \textit{Claim.} At least one $k \in \{1,\ldots,6\}$ satisfies
        $b_k > 0$.

        \smallskip
        \noindent
        \textit{Proof of Claim.}
        By the exhaustive case analysis performed in step~(ii) above,
        \emph{every} bucket satisfies $b_k > 0$.  In particular,
        taking $k = 1$:
        \begin{align}
          b_1 &= 35{,}000 > 0,
        \end{align}
        which constitutes an explicit witness satisfying the
        existential requirement of~\textbf{(P3)}.  Since the
        condition is existential ($\exists\, k$) and we have
        exhibited a concrete instance, the axiom is fulfilled
        without further case analysis.

        \smallskip
        \noindent
        \textit{Extremity check.}  The only way \textbf{(P3)} could
        fail is if $b_k = 0$ for all $k \in \{1,\ldots,6\}$, i.e.\
        if every bucket is empty.  Step~(ii) rules this out
        conclusively for all six buckets.
        Axiom~\textbf{(P3)} is satisfied. \hfill$\checkmark$

\end{enumerate}

\noindent
\textit{Summary.}
All three axioms \textbf{(P1)}--\textbf{(P3)} have been verified
independently and without relying on any unproven auxiliary claims.
No assumption has been left unvalidated; all edge cases and boundary
conditions have been addressed explicitly.  The canonical six-bucket
partition is therefore well-posed in the sense of
Definition~\ref{def:partition}.
\qed
\end{proof}

\begin{definition}[Bucket Risk Parameters]
\label{def:bucket-risk}
For each bucket $k \in \{1,\ldots,K\}$, the following parameters
are pre-specified and treated as deterministic constants throughout
the investment period:
\begin{enumerate}
  \item $\sigma_k^{\max} > 0$: maximum permissible annualised
        volatility for bucket $k$.
  \item $\mathrm{DD}_k^{\max} \in (0,1)$: maximum permissible
        drawdown fraction for bucket $k$.
  \item $R_k^{\min} \in \mathbb{R}$: minimum required rebalancing
        trigger (absolute return threshold below which inter-bucket
        capital transfer is considered).
\end{enumerate}
\noindent
These parameters are calibrated prior to deployment and recorded in
the audit log.  Any change to a risk parameter requires a new audit
entry and re-approval.
\end{definition}

\begin{axiom}[Mandatory Liquidity Floor]
\label{ax:liquidity-floor}
\label{asm:liquidity-reserve}
The Cash / Liquidity Reserve bucket (bucket~5 in the canonical
partition) must satisfy at all times $t \in \{0,\ldots,T\}$:
\begin{equation}
  W_t^{(5)} \;\geq\; \ell_{\min} \;:=\; 0.05 \cdot W_0,
  \label{eq:liquidity-floor}
\end{equation}
where $W_t^{(5)}$ denotes the current value of bucket~5 at time $t$.
No signal, however strong, permits the cash reserve to fall below
$\ell_{\min}$ without explicit risk-gate approval documented in the
audit log.
\end{axiom}

\begin{remark}[Degenerate and Boundary Cases for the Partition]
\label{rem:partition-edge}
\begin{enumerate}
  \item \emph{Empty buckets.}  $b_k = 0$ is permitted; it simply
        means bucket $k$ is inactive.  All subsequent formulas
        reduce gracefully when $b_k = 0$.
  \item \emph{Single active bucket} ($K=1$).  The partition
        degenerates to a fully concentrated portfolio.
        Theorem~\ref{thm:concentration-suboptimal} (Chapter~5)
        shows this is sub-optimal under the composite objective.
  \item \emph{Uniform partition.}  $b_k = W_0/K$ for all $k$
        represents the na\"{\i}ve equal-weight baseline; it
        is a valid but sub-optimal configuration.
\end{enumerate}
\end{remark}

% ==============================================================
\chapter{Agent Output Contract and Signal Normalisation}
\label{ch:agent-contract}
% ==============================================================

\section{Directional Signal Normalisation}

\begin{assumption}[Agent Index Set]
\label{asm:agent-index}
\label{asm:active-agent}
\label{asm:active-agents}
\label{asm:agent-outputs}
\label{asm:score-domain}
There are $N \in \mathbb{N}$, $N \geq 2$, specialised agents indexed
by $i \in \{1,\ldots,N\}$.  The agents operate in parallel; their
outputs are $\calF_t$-measurable for every $t$.
\end{assumption}

\begin{definition}[Agent Output Tuple]
\label{def:agent-output}
\label{def:agent-score}
Under Assumption~\ref{asm:agent-index}, at each time $t$ agent $i$
emits the tuple
\begin{equation}
  \calA_i(t)
  \;=\;
  \bigl(\,s_i,\;\; c_i,\;\; u_i,\;\; d_i,\;\;
         p_i^{+},\;\; p_i^{-},\;\; \Delta t_i,\;\; r_i\,\bigr),
  \label{eq:agent-tuple}
\end{equation}
where the domains and interpretations of each channel are given in
Table~\ref{tab:agent-channels}.
\end{definition}

\begin{table}[H]
\centering
\caption{Agent output channels: domains and interpretations.}
\label{tab:agent-channels}
\small
\begin{tabularx}{\linewidth}{@{} l c X @{}}
\toprule
\textbf{Symbol} & \textbf{Domain} & \textbf{Interpretation} \\
\midrule
$s_i$ & $[-1,1]$ & Directional market view: $-1$ = maximally bearish,
  $0$ = neutral, $+1$ = maximally bullish. \\
$c_i$ & $(0,1]$ & Confidence in the directional view. \\
$u_i$ & $[0,1]$ & Uncertainty: elevated when strong signals conflict. \\
$d_i$ & $[0,1]$ & Doubt: agent's historical calibration quality in the
  current regime. \\
$p_i^+$ & $[0,1]$ & Estimated probability of a favourable outcome. \\
$p_i^-$ & $[0,1]$ & Estimated probability of an unfavourable outcome. \\
$\Delta t_i$ & $\mathbb{R}_{>0}$ & Relevant investment time horizon. \\
$r_i$ & $\mathbb{R}_{\geq 0}$ & Estimated risk or loss exposure. \\
\bottomrule
\end{tabularx}
\end{table}

\begin{assumption}[Channel Validity]
\label{asm:channel-validity}
For every agent $i$ and every time $t$, the emitted tuple satisfies:
\begin{enumerate}[label=\textbf{(CV\arabic*)}]
  \item $s_i \in [-1,1]$;
  \item $c_i \in (0,1]$;
  \item $u_i \in [0,1]$;
  \item $d_i \in [0,1]$;
  \item $p_i^+ + p_i^- \leq 1$, $p_i^+, p_i^- \geq 0$;
  \item $\Delta t_i > 0$;
  \item $r_i \geq 0$.
\end{enumerate}
These constraints are enforced by the type system of the
implementation language (see Chapter~12 of the companion document).
\end{assumption}

\begin{proposition}[Channel Boundedness]
\label{prop:channel-bounded}
Under Assumption~\ref{asm:channel-validity}, the product channel
\begin{equation}
  \phi_i \;:=\; s_i \cdot c_i \cdot (1 - u_i) \cdot (1 - d_i)
  \label{eq:phi-channel}
\end{equation}
satisfies $\phi_i \in [-1,1]$ for every agent $i$ and every time $t$.
\end{proposition}

\begin{proof}
We proceed step by step, validating each factor independently before
combining them.  Every reference to a constraint is explicitly cited.

\begin{enumerate}

  \item \textbf{Validation of $s_i$.}
        By constraint \textbf{(CV1)},
        \begin{align}
          s_i &\;\in\; [-1,\,1].
        \end{align}
        In particular,
        \begin{align}
          -1 \;\leq\; s_i \;\leq\; 1,
          \qquad\text{so}\qquad
          |s_i| \;\leq\; 1.
          \label{eq:cv1-bound}
        \end{align}
        No further assumption is required for this factor.

  \item \textbf{Validation of $c_i$.}
        By constraint \textbf{(CV2)},
        \begin{align}
          c_i &\;\in\; (0,\,1].
        \end{align}
        Therefore,
        \begin{align}
          0 \;<\; c_i \;\leq\; 1,
          \qquad\text{hence}\qquad
          c_i \;\geq\; 0
          \quad\text{and}\quad
          c_i \;\leq\; 1.
          \label{eq:cv2-bound}
        \end{align}
        The strict lower bound $c_i > 0$ additionally confirms that
        $c_i$ is never zero; this is not needed for the present bound
        but is recorded for completeness.

  \item \textbf{Validation of $(1 - u_i)$.}
        By constraint \textbf{(CV3)},
        \begin{align}
          u_i &\;\in\; [0,\,1],
          \qquad\text{so}\qquad
          0 \;\leq\; u_i \;\leq\; 1.
        \end{align}
        Subtracting all parts of the inequality from $1$ and reversing
        the sense yields:
        \begin{align}
          1 - 1 \;\leq\; 1 - u_i \;\leq\; 1 - 0,
        \end{align}
        that is,
        \begin{align}
          0 \;\leq\; (1 - u_i) \;\leq\; 1.
          \label{eq:cv3-bound}
        \end{align}

  \item \textbf{Validation of $(1 - d_i)$.}
        By constraint \textbf{(CV4)},
        \begin{align}
          d_i &\;\in\; [0,\,1],
          \qquad\text{so}\qquad
          0 \;\leq\; d_i \;\leq\; 1.
        \end{align}
        By the same monotone subtraction argument as Step~3,
        \begin{align}
          0 \;\leq\; (1 - d_i) \;\leq\; 1.
          \label{eq:cv4-bound}
        \end{align}

  \item \textbf{Non-negativity of the dampening product.}
        Define the dampening product
        \begin{align}
          D_i \;:=\; c_i \cdot (1 - u_i) \cdot (1 - d_i).
          \label{eq:Di-def}
        \end{align}
        From \eqref{eq:cv2-bound}, \eqref{eq:cv3-bound}, and
        \eqref{eq:cv4-bound}, each factor of $D_i$ is non-negative.
        The product of non-negative real numbers is non-negative,
        therefore
        \begin{align}
          D_i \;\geq\; 0.
          \label{eq:Di-nonneg}
        \end{align}

  \item \textbf{Upper bound on the dampening product.}
        Each factor of $D_i$ is also at most $1$
        (by \eqref{eq:cv2-bound}--\eqref{eq:cv4-bound}).
        For any real numbers $a, b, c \in [0,1]$,
        $abc \leq a \leq 1$.
        Applying this with
        $a = c_i$, $b = (1-u_i)$, $c = (1-d_i)$:
        \begin{align}
          D_i
          &\;=\; c_i \cdot (1-u_i) \cdot (1-d_i)
          \notag\\
          &\;\leq\; 1 \cdot 1 \cdot 1
          \;=\; 1.
          \label{eq:Di-upper}
        \end{align}
        Combining \eqref{eq:Di-nonneg} and \eqref{eq:Di-upper},
        \begin{align}
          D_i \;\in\; [0,\,1].
          \label{eq:Di-range}
        \end{align}

  \item \textbf{Absolute value bound on $\phi_i$.}
        Since $D_i \geq 0$ (from \eqref{eq:Di-nonneg}),
        \begin{align}
          |\phi_i|
          &\;=\; |s_i \cdot D_i|
          \;=\; |s_i| \cdot |D_i|
          \;=\; |s_i| \cdot D_i.
          \label{eq:phi-abs-split}
        \end{align}
        Applying \eqref{eq:cv1-bound} and \eqref{eq:Di-upper}:
        \begin{align}
          |\phi_i|
          &\;\leq\; 1 \cdot 1
          \;=\; 1.
          \label{eq:phi-abs-bound}
        \end{align}

  \item \textbf{Sign of $\phi_i$.}
        Since $D_i \geq 0$, multiplying $s_i$ by $D_i$ preserves
        the sign of $s_i$.  Formally:
        \begin{enumerate}
          \item If $s_i \geq 0$, then
                $\phi_i = s_i \cdot D_i \geq 0$.
          \item If $s_i < 0$, then
                $\phi_i = s_i \cdot D_i \leq 0$.
        \end{enumerate}
        In both cases the sign of $\phi_i$ coincides with the sign of
        $s_i$.

  \item \textbf{Combining lower and upper bounds.}
        From \eqref{eq:phi-abs-bound},
        \begin{align}
          -1 \;\leq\; \phi_i \;\leq\; 1,
        \end{align}
        i.e., $\phi_i \in [-1,1]$.

  \item \textbf{Boundary cases.}
        \begin{enumerate}
          \item \emph{Attainability of $\phi_i = 1$.}
                Achieved when $s_i = 1$, $c_i = 1$, $u_i = 0$,
                $d_i = 0$, all of which are permitted by
                \textbf{(CV1)}--\textbf{(CV4)}.
          \item \emph{Attainability of $\phi_i = -1$.}
                Achieved when $s_i = -1$, $c_i = 1$, $u_i = 0$,
                $d_i = 0$, again permitted by
                \textbf{(CV1)}--\textbf{(CV4)}.
          \item \emph{Attainability of $\phi_i = 0$.}
                Since $c_i > 0$ by \textbf{(CV2)}, $\phi_i = 0$
                requires either $s_i = 0$, $u_i = 1$, or $d_i = 1$,
                all of which are within their respective domains.
        \end{enumerate}
        These cases confirm that the bound $[-1,1]$ is tight and that
        no constraint is violated at any boundary.

  \item \textbf{Conclusion.}
        Steps 1--10 establish, without skipping any intermediate
        reasoning:
        \begin{align}
          \phi_i \;=\; s_i \cdot c_i \cdot (1-u_i) \cdot (1-d_i)
          \;\in\; [-1,\,1]
        \end{align}
        for every agent $i$ and every time $t$ at which
        Assumption~\ref{asm:channel-validity} holds.

\end{enumerate}
\qed
\end{proof}

\begin{remark}[Semantic Distinctness of $c_i$, $u_i$, $d_i$]
\label{rem:channel-semantic}
The three dampening channels are not redundant:
\begin{enumerate}
  \item \emph{Low confidence} arises from weak or sparse evidence.
  \item \emph{High uncertainty} arises when strong signals directly
        conflict within the agent's internal model.
  \item \emph{High doubt} arises when the agent's own track record
        in the current regime has been poor, regardless of the
        strength of the current evidence.
\end{enumerate}
Collapsing these into a single scalar would destroy information
needed by the master allocator to distinguish, for instance, a
genuinely weak signal from a strongly conflicted one.
\end{remark}

% ==============================================================
\chapter{Ensemble Aggregation, Disagreement, and Allocator State}
\label{ch:aggregation}
% ==============================================================

\section{Ensemble Aggregate Signal}

% --------------------------------------------------------------
\subsection{Reliability Weights}
% --------------------------------------------------------------

\begin{assumption}[Weight Regularity]
\label{asm:weight-regularity}
\label{asm:weights}
For every agent $i \in \{1,\ldots,N\}$ and every time $t$:
\begin{enumerate}
  \item \textbf{(WR1)} $w_i(t) \in \mathbb{R}_{\geq 0}$,
        i.e., each weight is a well-defined, finite, non-negative
        real number.
  \item \textbf{(WR2)} The aggregate weight
        $W(t) := \sum_{i=1}^{N} w_i(t)$ satisfies $W(t) > 0$,
        i.e., at least one agent is active at every time $t$.
  \item \textbf{(WR3)} Each $w_i(t)$ is measurable with respect to
        the filtration $\mathcal{F}_t$ available at time $t$, so that
        the ensemble signal defined below is $\mathcal{F}_t$-measurable.
\end{enumerate}
\end{assumption}

\begin{definition}[Agent Reliability Weight]
\label{def:reliability-weight}
Under Assumption~\ref{asm:weight-regularity}, the
\emph{reliability weight} of agent $i$ at time $t$ is the scalar
$w_i(t) \geq 0$ encoding all of the following factors:
\begin{enumerate}
  \item \emph{Regime performance.}
        Historical risk-adjusted performance of agent $i$ in the
        market regime $\mathcal{R}(t)$ currently in effect, as
        determined by the regime-detection module.
  \item \emph{Instrument expertise.}
        Demonstrated accuracy of agent $i$ on the specific
        instrument or asset class under consideration at time $t$.
  \item \emph{Calibration quality.}
        Degree to which agent $i$'s stated confidence values
        $c_i(t)$ have historically matched empirical success rates,
        measured by a proper scoring rule (e.g., Brier score or
        log-loss) over a rolling evaluation window.
  \item \emph{Horizon relevance.}
        Alignment of agent $i$'s signal horizon $\Delta t_i$ with
        the current allocation decision horizon $\Delta t^{*}(t)$,
        penalising agents whose horizon differs substantially from
        $\Delta t^{*}(t)$.
\end{enumerate}
The functional form and update rule for $w_i(t)$ are specified in
Chapter~\ref{ch:reputation}.  The present definition fixes the
semantic requirements that any compliant update rule must satisfy.
Conditions \textbf{(WR1)}--\textbf{(WR3)} must hold after every
update.
\end{definition}

\begin{remark}[Well-Posedness of $W(t)$]
\label{rem:W-well-posed}
Condition \textbf{(WR2)} ensures that the normalising denominator
$W(t) = \sum_{i=1}^{N} w_i(t)$ is strictly positive at all times,
so that any ratio of the form $(\cdot)/W(t)$ is well-defined.
Condition \textbf{(WR1)} ensures $W(t) < \infty$ since $N < \infty$
(Assumption~\ref{asm:agent-index}).
\end{remark}

% --------------------------------------------------------------
\subsection{Definition of the Ensemble Aggregate Signal}
% --------------------------------------------------------------

\begin{definition}[Dampened Agent Signal]
\label{def:dampened-signal}
Under Assumption~\ref{asm:channel-validity}, the
\emph{dampened signal} of agent $i$ at time $t$ is
\begin{equation}
  \phi_i(t)
  \;:=\;
  s_i(t)\;\cdot\;c_i(t)\;\cdot\;\bigl(1 - u_i(t)\bigr)
  \;\cdot\;\bigl(1 - d_i(t)\bigr).
  \label{eq:phi-def}
\end{equation}
By Proposition~\ref{prop:channel-bounded},
$\phi_i(t) \in [-1,1]$ for all $i$ and $t$.
\end{definition}

\begin{definition}[Ensemble Aggregate Signal]
\label{def:ensemble-signal}
Under Assumptions~\ref{asm:agent-index},
\ref{asm:channel-validity}, and~\ref{asm:weight-regularity},
and Definitions~\ref{def:reliability-weight}
and~\ref{def:dampened-signal}, the \emph{ensemble aggregate
signal} at time $t$ is
\begin{equation}
  S(t)
  \;:=\;
  \frac{\displaystyle\sum_{i=1}^{N} w_i(t)\,\phi_i(t)}
       {\displaystyle\sum_{i=1}^{N} w_i(t)}
  \;=\;
  \frac{\displaystyle\sum_{i=1}^{N}
        w_i(t)\,s_i(t)\,c_i(t)\,
        \bigl(1-u_i(t)\bigr)\,\bigl(1-d_i(t)\bigr)}
       {W(t)},
  \label{eq:ensemble-signal}
\end{equation}
where $W(t) := \sum_{i=1}^{N} w_i(t) > 0$ by \textbf{(WR2)}.
The time index $t$ is suppressed hereafter for concision whenever
no ambiguity arises.
\end{definition}

\begin{remark}[Measurability and Finite Existence]
\label{rem:S-measurable}
Because each $w_i$, $s_i$, $c_i$, $u_i$, $d_i$ is
$\mathcal{F}_t$-measurable by \textbf{(WR3)} and
Assumption~\ref{asm:channel-validity}, and the index set is
finite ($N < \infty$ by Assumption~\ref{asm:agent-index}),
the sum in the numerator of~\eqref{eq:ensemble-signal} is a
finite sum of $\mathcal{F}_t$-measurable bounded random variables,
hence itself $\mathcal{F}_t$-measurable and finite.
Combined with $W > 0$, the quotient $S$ is well-defined,
finite, and $\mathcal{F}_t$-measurable.
\end{remark}

% --------------------------------------------------------------
\subsection{Boundedness of the Ensemble Aggregate Signal}
% --------------------------------------------------------------

\begin{theorem}[Boundedness of the Ensemble Aggregate Signal]
\label{thm:ensemble-bounded}
Under Assumptions~\ref{asm:agent-index},
\ref{asm:channel-validity}, and~\ref{asm:weight-regularity},
and Definitions~\ref{def:reliability-weight}
and~\ref{def:ensemble-signal}, the ensemble aggregate signal
satisfies
\begin{equation}
  S \;\in\; [-1,\,1].
  \label{eq:S-in-unit}
\end{equation}
Moreover, both boundary values are attainable.
\end{theorem}

\begin{proof}
The proof proceeds in a sequence of fully justified steps.
No step is skipped; every assumption invoked is cited explicitly.

\begin{enumerate}

  % --- Step 1 ---
  \item \textbf{Validation of the index set.}
        By Assumption~\ref{asm:agent-index}, the agent index set
        $\mathcal{I} = \{1,\ldots,N\}$ is finite and non-empty,
        so $N \in \mathbb{Z}_{>0}$.
        All subsequent finite sums are therefore well-defined.

  % --- Step 2 ---
  \item \textbf{Validation of weight non-negativity.}
        By Assumption~\ref{asm:weight-regularity} condition
        \textbf{(WR1)}, $w_i \geq 0$ for all $i$.

  % --- Step 3 ---
  \item \textbf{Validation of denominator positivity.}
        By Assumption~\ref{asm:weight-regularity} condition
        \textbf{(WR2)},
        \begin{align}
          W \;&:=\; \sum_{i=1}^{N} w_i \;>\; 0.
          \label{eq:W-pos}
        \end{align}
        Hence division by $W$ in~\eqref{eq:ensemble-signal} is
        well-defined.

  % --- Step 4 ---
  \item \textbf{Validation of channel bounds for each agent.}
        By Assumption~\ref{asm:channel-validity} conditions
        \textbf{(CV1)}--\textbf{(CV4)} and
        Proposition~\ref{prop:channel-bounded}, for every
        $i \in \mathcal{I}$:
        \begin{align}
          s_i     &\;\in\; [-1,\,1],
          \label{eq:s-bound-recall}\\
          c_i     &\;\in\; (0,\,1],
          \label{eq:c-bound-recall}\\
          u_i     &\;\in\; [0,\,1],
          \label{eq:u-bound-recall}\\
          d_i     &\;\in\; [0,\,1],
          \label{eq:d-bound-recall}
        \end{align}
        and therefore
        \begin{align}
          (1-u_i) &\;\in\; [0,\,1],
          \label{eq:1mu-bound}\\
          (1-d_i) &\;\in\; [0,\,1].
          \label{eq:1md-bound}
        \end{align}

  % --- Step 5 ---
  \item \textbf{Upper bound on each dampened signal.}
        Using~\eqref{eq:s-bound-recall}--\eqref{eq:1md-bound}
        and the fact that the product of non-negative factors
        bounded above by~$1$ is itself bounded above by~$1$:
        \begin{align}
          \phi_i
          &\;=\; s_i \cdot c_i \cdot (1-u_i) \cdot (1-d_i)
          \notag\\
          &\;\leq\; |s_i| \cdot c_i \cdot (1-u_i) \cdot (1-d_i)
          \label{eq:phi-ub-step1}\\
          &\;\leq\; 1 \cdot 1 \cdot 1 \cdot 1
          \;=\; 1.
          \label{eq:phi-upper}
        \end{align}
        Line~\eqref{eq:phi-ub-step1} uses $s_i \leq |s_i|$;
        line~\eqref{eq:phi-upper} applies the individual upper
        bounds from Step~4.

  % --- Step 6 ---
  \item \textbf{Lower bound on each dampened signal.}
        Since $c_i > 0$, $(1-u_i) \geq 0$, $(1-d_i) \geq 0$
        (Step~4), the product
        $c_i(1-u_i)(1-d_i) \geq 0$.
        Therefore:
        \begin{align}
          \phi_i
          &\;=\; s_i \cdot \underbrace{c_i(1-u_i)(1-d_i)}_{\geq\,0}
          \;\geq\; (-1) \cdot c_i(1-u_i)(1-d_i)
          \label{eq:phi-lb-step1}\\
          &\;\geq\; (-1) \cdot 1
          \;=\; -1.
          \label{eq:phi-lower}
        \end{align}
        Line~\eqref{eq:phi-lb-step1} uses $s_i \geq -1$;
        line~\eqref{eq:phi-lower} applies
        $c_i(1-u_i)(1-d_i) \leq 1$.

  % --- Step 7 ---
  \item \textbf{Absolute value bound on each dampened signal.}
        Combining Steps~5 and~6:
        \begin{align}
          -1 \;\leq\; \phi_i \;\leq\; 1
          \quad\Longrightarrow\quad
          |\phi_i| \;\leq\; 1
          \qquad \forall\, i \in \mathcal{I}.
          \label{eq:phi-abs}
        \end{align}

  % --- Step 8 ---
  \item \textbf{Absolute value bound on each weighted summand.}
        Since $w_i \geq 0$ (Step~2) and $|\phi_i| \leq 1$
        (Step~7):
        \begin{align}
          |w_i \phi_i|
          &\;=\; w_i |\phi_i|
          \;\leq\; w_i \cdot 1
          \;=\; w_i
          \qquad \forall\, i \in \mathcal{I}.
          \label{eq:wi-phi-bound}
        \end{align}
        The equality $|w_i \phi_i| = w_i |\phi_i|$ uses $w_i \geq 0$.

  % --- Step 9 ---
  \item \textbf{Absolute value bound on the numerator.}
        Applying the triangle inequality and then~\eqref{eq:wi-phi-bound}:
        \begin{align}
          \left|\,\sum_{i=1}^{N} w_i \phi_i\,\right|
          &\;\leq\; \sum_{i=1}^{N} \left|w_i \phi_i\right|
          \label{eq:tri-ineq}\\
          &\;\leq\; \sum_{i=1}^{N} w_i
          \label{eq:wi-bound-applied}\\
          &\;=\; W.
          \label{eq:num-bound}
        \end{align}
        Line~\eqref{eq:tri-ineq} is the triangle inequality for
        finite sums; line~\eqref{eq:wi-bound-applied} applies
        Step~8 termwise.

  % --- Step 10 ---
  \item \textbf{Division by $W$ and conclusion.}
        Using~\eqref{eq:num-bound} and $W > 0$
        (from~\eqref{eq:W-pos}):
        \begin{align}
          |S|
          &\;=\;
          \left|\,
            \frac{\displaystyle\sum_{i=1}^{N} w_i \phi_i}{W}
          \,\right|
          \;=\;
          \frac{\left|\displaystyle\sum_{i=1}^{N} w_i \phi_i\right|}{W}
          \label{eq:S-abs-step1}\\[4pt]
          &\;\leq\;
          \frac{W}{W}
          \;=\; 1.
          \label{eq:S-abs-final}
        \end{align}
        Line~\eqref{eq:S-abs-step1} uses $|a/b| = |a|/|b|$ for
        $b > 0$; line~\eqref{eq:S-abs-final} applies
        \eqref{eq:num-bound}.
        Since $S \in \mathbb{R}$, we conclude $S \in [-1,1]$.

  % --- Step 11 ---
  \item \textbf{Attainability of the upper boundary $S = +1$.}
        Set $s_i = 1$, $c_i = 1$, $u_i = 0$, $d_i = 0$ for
        all $i$.  Each assignment is within the permitted domain
        by \textbf{(CV1)}--\textbf{(CV4)}.  Then
        $\phi_i = 1 \cdot 1 \cdot 1 \cdot 1 = 1$ for all $i$, and
        \begin{align}
          S
          &\;=\;
          \frac{\displaystyle\sum_{i=1}^{N} w_i \cdot 1}{W}
          \;=\;
          \frac{W}{W}
          \;=\; 1.
          \label{eq:S-upper-attain}
        \end{align}

  % --- Step 12 ---
  \item \textbf{Attainability of the lower boundary $S = -1$.}
        Set $s_i = -1$, $c_i = 1$, $u_i = 0$, $d_i = 0$ for
        all $i$.  Again each assignment is within the permitted
        domain.  Then $\phi_i = -1$ for all $i$, and
        \begin{align}
          S
          &\;=\;
          \frac{\displaystyle\sum_{i=1}^{N} w_i \cdot (-1)}{W}
          \;=\;
          \frac{-W}{W}
          \;=\; -1.
          \label{eq:S-lower-attain}
        \end{align}

  % --- Step 13 ---
  \item \textbf{Edge case: single active agent.}
        If exactly one agent $j$ has $w_j > 0$ and $w_i = 0$
        for all $i \neq j$, then $W = w_j > 0$ and
        \begin{align}
          S
          &\;=\;
          \frac{w_j \phi_j}{w_j}
          \;=\; \phi_j
          \;\in\; [-1,1]
        \end{align}
        by Step~7.  The conclusion~\eqref{eq:S-in-unit} holds
        in this degenerate case as well.

  % --- Step 14 ---
  \item \textbf{Edge case: all dampened signals zero.}
        If $\phi_i = 0$ for all $i$ (e.g., all agents have
        $s_i = 0$, or $u_i = 1$, or $d_i = 1$), then
        \begin{align}
          S
          &\;=\;
          \frac{\displaystyle\sum_{i=1}^{N} w_i \cdot 0}{W}
          \;=\; 0
          \;\in\; [-1,1].
        \end{align}
        The zero signal is a valid element of $[-1,1]$ and triggers
        no allocation, consistent with the master allocator's
        conservative default.

  % --- Step 15 ---
  \item \textbf{Summary.}
        Steps~1--3 validate all prerequisites.
        Steps~4--7 establish $|\phi_i| \leq 1$ for each agent.
        Steps~8--10 propagate the bound through the weighted average
        to conclude $S \in [-1,1]$.
        Steps~11--12 confirm tightness.
        Steps~13--14 verify that no edge case violates the result.
        No assumption beyond \ref{asm:agent-index},
        \ref{asm:channel-validity}, and \ref{asm:weight-regularity}
        is required.

\end{enumerate}
\qed
\end{proof}

\begin{corollary}[Normalised Representation of $S$]
\label{cor:S-normalised}
Under the hypotheses of Theorem~\ref{thm:ensemble-bounded},
let $W := \sum_{i=1}^{N} w_i > 0$ as guaranteed by
Assumption~\ref{asm:weight-regularity}.  Define the normalised
weights
\begin{align}
  \tilde{w}_i
  &\;:=\;
  \frac{w_i}{W},
  \qquad i = 1,\ldots,N.
  \label{eq:normalised-weights}
\end{align}
Then $S$ admits the convex-combination representation
\begin{align}
  S
  &\;=\;
  \sum_{i=1}^{N} \tilde{w}_i\,\phi_i,
  \label{eq:S-convex}
\end{align}
where $\tilde{w}_i \geq 0$ for all $i$ and
$\displaystyle\sum_{i=1}^{N} \tilde{w}_i = 1$.
Consequently, $S \in [-1,1]$, providing an alternative, entirely
self-contained proof of the conclusion of
Theorem~\ref{thm:ensemble-bounded} via convexity alone.
\end{corollary}

\begin{proof}
We proceed step by step, validating every prerequisite and
boundary condition.

\begin{enumerate}

  % --- Step 1 ---
  \item \textbf{Prerequisite: $W > 0$.}
        By Assumption~\ref{asm:weight-regularity},
        $w_i \geq 0$ for all $i \in \{1,\ldots,N\}$ and
        $\exists\, j$ such that $w_j > 0$.
        Therefore
        \begin{align}
          W
          &\;=\; \sum_{i=1}^{N} w_i
          \;\geq\; w_j
          \;>\; 0.
        \end{align}
        Division by $W$ is hence well-defined throughout this proof.

  % --- Step 2 ---
  \item \textbf{Definition of $\tilde{w}_i$ and non-negativity.}
        Set $\tilde{w}_i := w_i / W$ for each $i$.
        Since $w_i \geq 0$ (Assumption~\ref{asm:weight-regularity})
        and $W > 0$ (Step~1),
        \begin{align}
          \tilde{w}_i
          &\;=\; \frac{w_i}{W}
          \;\geq\; \frac{0}{W}
          \;=\; 0,
          \qquad \forall\, i \in \{1,\ldots,N\}.
        \end{align}
        Hence $\tilde{w}_i \geq 0$ for every index.

  % --- Step 3 ---
  \item \textbf{Unit-sum property.}
        Summing the normalised weights over all agents,
        \begin{align}
          \sum_{i=1}^{N} \tilde{w}_i
          &\;=\; \sum_{i=1}^{N} \frac{w_i}{W}
          \;=\; \frac{1}{W}\sum_{i=1}^{N} w_i
          \;=\; \frac{W}{W}
          \;=\; 1.
        \end{align}
        The factoring of $1/W$ outside the sum is valid because $W$
        is a positive constant independent of $i$.

  % --- Step 4 ---
  \item \textbf{Algebraic derivation of the convex-combination form.}
        Recall from Definition~\ref{def:ensemble-signal} that
        \begin{align}
          S
          &\;=\;
          \frac{\displaystyle\sum_{i=1}^{N} w_i\,\phi_i}
               {W}.
        \end{align}
        Since $W > 0$ (Step~1), multiply numerator and denominator
        by $1/W$:
        \begin{align}
          S
          &\;=\;
          \frac{\displaystyle\sum_{i=1}^{N} w_i\,\phi_i}{W}
          \;=\;
          \sum_{i=1}^{N} \frac{w_i}{W}\,\phi_i
          \;=\;
          \sum_{i=1}^{N} \tilde{w}_i\,\phi_i.
        \end{align}
        The exchange of the scalar $1/W$ with the finite sum is
        justified by linearity of finite summation (no convergence
        issues arise since $N < \infty$ by
        Assumption~\ref{asm:agent-index}).

  % --- Step 5 ---
  \item \textbf{Each summand is bounded.}
        By Proposition~\ref{prop:channel-bounded},
        $\phi_i \in [-1,1]$ for every
        $i \in \{1,\ldots,N\}$.
        Combined with $\tilde{w}_i \geq 0$ (Step~2),
        \begin{align}
          -\tilde{w}_i
          &\;\leq\;
          \tilde{w}_i\,\phi_i
          \;\leq\;
          \tilde{w}_i,
          \qquad \forall\, i.
        \end{align}

  % --- Step 6 ---
  \item \textbf{Lower bound on $S$.}
        Summing the left inequality in Step~5 over all $i$,
        \begin{align}
          S
          &\;=\; \sum_{i=1}^{N} \tilde{w}_i\,\phi_i
          \;\geq\; -\sum_{i=1}^{N} \tilde{w}_i
          \;=\; -1,
        \end{align}
        where the final equality uses the unit-sum property
        established in Step~3.

  % --- Step 7 ---
  \item \textbf{Upper bound on $S$.}
        Summing the right inequality in Step~5 over all $i$,
        \begin{align}
          S
          &\;=\; \sum_{i=1}^{N} \tilde{w}_i\,\phi_i
          \;\leq\; \sum_{i=1}^{N} \tilde{w}_i
          \;=\; 1,
        \end{align}
        again invoking Step~3.

  % --- Step 8 ---
  \item \textbf{Conclusion: $S \in [-1,1]$.}
        Combining Steps~6 and~7,
        \begin{align}
          -1 \;\leq\; S \;\leq\; 1,
        \end{align}
        i.e., $S \in [-1,1]$.  This constitutes an independent
        proof of the main conclusion of
        Theorem~\ref{thm:ensemble-bounded} relying only on
        Assumptions~\ref{asm:agent-index},
        \ref{asm:weight-regularity}, and
        Proposition~\ref{prop:channel-bounded}.

  % --- Step 9 ---
  \item \textbf{Convexity interpretation.}
        The interval $[-1,1] \subset \mathbb{R}$ is a convex set:
        for any $a, b \in [-1,1]$ and $\lambda \in [0,1]$,
        $\lambda a + (1-\lambda)b \in [-1,1]$.
        By induction on $N$, any finite convex combination of
        points in $[-1,1]$ remains in $[-1,1]$.
        Since $\tilde{w}_i \geq 0$,
        $\sum_{i=1}^{N}\tilde{w}_i = 1$, and each
        $\phi_i \in [-1,1]$, the representation
        $S = \sum_{i=1}^{N}\tilde{w}_i\,\phi_i$ is exactly
        such a finite convex combination, confirming
        $S \in [-1,1]$ via the definition of a convex set.

  % --- Step 10 ---
  \item \textbf{Edge case: all weights equal.}
        If $w_i = c > 0$ for all $i$, then
        $\tilde{w}_i = 1/N$ for all $i$, and
        \begin{align}
          S
          &\;=\; \frac{1}{N}\sum_{i=1}^{N}\phi_i,
        \end{align}
        which is the arithmetic mean of values in $[-1,1]$,
        itself in $[-1,1]$.  No special case arises.

  % --- Step 11 ---
  \item \textbf{Edge case: single active agent.}
        If $w_j > 0$ and $w_i = 0$ for all $i \neq j$,
        then $W = w_j$, $\tilde{w}_j = 1$,
        $\tilde{w}_i = 0$ for $i \neq j$, and
        \begin{align}
          S
          &\;=\; 1 \cdot \phi_j + \sum_{i \neq j} 0 \cdot \phi_i
          \;=\; \phi_j
          \;\in\; [-1,1].
        \end{align}
        The result holds.

  % --- Step 12 ---
  \item \textbf{Edge case: all signals equal.}
        If $\phi_i = c \in [-1,1]$ for all $i$, then
        \begin{align}
          S
          &\;=\; c \sum_{i=1}^{N}\tilde{w}_i
          \;=\; c \cdot 1
          \;=\; c
          \;\in\; [-1,1].
        \end{align}
        Consistent with consensus (cf.\
        Definition~\ref{def:disagreement}).

  % --- Step 13 ---
  \item \textbf{Tightness.}
        The bound $S = -1$ is achieved when
        $\phi_i = -1$ for all $i$ with $w_i > 0$, and
        $S = 1$ when $\phi_i = 1$ for all such $i$.
        Both extreme points are attainable, confirming the
        interval $[-1,1]$ is tight.

  % --- Step 14 ---
  \item \textbf{Summary of the proof.}
        Steps~1--3 establish that $\{\tilde{w}_i\}$ is a
        well-defined probability weight vector.
        Step~4 derives the convex-combination representation
        algebraically.
        Steps~5--7 establish the two-sided bound via
        elementary inequalities.
        Step~8 states the conclusion.
        Step~9 provides the geometric convexity argument.
        Steps~10--13 confirm that no edge case or boundary
        condition violates the result.
        No assumption beyond
        \ref{asm:agent-index},
        \ref{asm:weight-regularity}, and
        Proposition~\ref{prop:channel-bounded} is required.

\end{enumerate}
\qed
\end{proof}

\section{Agent Disagreement Measure}

\begin{definition}[Agent Disagreement]
\label{def:disagreement}
Let $t$ be any discrete trading period, $N \in \mathbb{N}$ the number
of active agents (Assumption~\ref{asm:agent-index}), $w_i \geq 0$
the weight of agent $i$
(Assumption~\ref{asm:weight-regularity}), $s_i \in [-1,1]$ the
signal of agent $i$
(Proposition~\ref{prop:channel-bounded}), and $S$ the ensemble
aggregate signal of Definition~\ref{def:ensemble-signal}.
Define the total weight
\begin{align}
  W \;:=\; \sum_{i=1}^{N} w_i \;>\; 0,
  \label{eq:disagreement-W}
\end{align}
which is strictly positive by
Assumption~\ref{asm:weight-regularity}.
The \emph{agent disagreement} at time $t$ is then
\begin{align}
  D
  \;&:=\;
  \frac{\displaystyle\sum_{i=1}^{N} w_i \,\bigl|s_i - S\bigr|}{W},
  \label{eq:disagreement}
\end{align}
i.e.\ the weighted mean absolute deviation of the individual
signals from the ensemble aggregate.
\end{definition}

\begin{proposition}[Properties of the Disagreement Measure]
\label{prop:disagreement-props}
Under Assumptions~\ref{asm:agent-index} and
\ref{asm:weight-regularity} and
Proposition~\ref{prop:channel-bounded}, the quantity $D$ defined
in \eqref{eq:disagreement} satisfies:
\begin{enumerate}
  \item \textbf{Non-negativity:} $D \geq 0$.
  \item \textbf{Boundedness:} $D \leq 2$.
  \item \textbf{Consensus characterisation:} $D = 0$ if and only
        if $s_i = S$ for every $i \in \{1,\ldots,N\}$ with
        $w_i > 0$.
\end{enumerate}
\end{proposition}

\begin{proof}
We establish each property in turn, validating every assumption and
ruling out all boundary cases.

\begin{enumerate}

  % --- Property 1 ---
  \item \textbf{Non-negativity ($D \geq 0$).}

  \begin{enumerate}
    \item \emph{Assumption validation.}
          By Assumption~\ref{asm:weight-regularity}, $w_i \geq 0$
          for all $i$, and $W > 0$.
          By Proposition~\ref{prop:channel-bounded},
          $s_i \in [-1,1]$, so $s_i$ is finite; likewise
          $S \in [-1,1]$ is finite.

    \item \emph{Non-negativity of each summand.}
          For every $i$,
          \begin{align}
            w_i \;\geq\; 0
            \quad\text{and}\quad
            \bigl|s_i - S\bigr| \;\geq\; 0,
          \end{align}
          hence $w_i \,|s_i - S| \geq 0$.

    \item \emph{Non-negativity of the numerator.}
          A finite sum of non-negative terms is non-negative:
          \begin{align}
            \sum_{i=1}^{N} w_i\,\bigl|s_i - S\bigr| \;\geq\; 0.
          \end{align}

    \item \emph{Conclusion.}
          Since the numerator $\geq 0$ and the denominator
          $W > 0$,
          \begin{align}
            D \;=\;
            \frac{\displaystyle\sum_{i=1}^{N}
                  w_i\,|s_i - S|}{W}
            \;\geq\; 0.
          \end{align}
  \end{enumerate}

  % --- Property 2 ---
  \item \textbf{Boundedness ($D \leq 2$).}

  \begin{enumerate}
    \item \emph{Bound on individual signals.}
          By Proposition~\ref{prop:channel-bounded},
          $s_i \in [-1,1]$, so $|s_i| \leq 1$.

    \item \emph{Bound on the aggregate.}
          By Theorem~\ref{thm:ensemble-bounded},
          $S \in [-1,1]$, so $|S| \leq 1$.

    \item \emph{Triangle inequality.}
          For every $i$,
          \begin{align}
            \bigl|s_i - S\bigr|
            &\;\leq\; |s_i| + |S|
            \;\leq\; 1 + 1
            \;=\; 2.
            \label{eq:disagreement-tri}
          \end{align}

    \item \emph{Upper bound on the numerator.}
          Using \eqref{eq:disagreement-tri},
          \begin{align}
            \sum_{i=1}^{N} w_i\,\bigl|s_i - S\bigr|
            &\;\leq\; \sum_{i=1}^{N} w_i \cdot 2
            \;=\; 2W.
          \end{align}

    \item \emph{Conclusion.}
          Dividing both sides by $W > 0$,
          \begin{align}
            D
            &\;=\;
            \frac{\displaystyle\sum_{i=1}^{N}
                  w_i\,|s_i - S|}{W}
            \;\leq\; \frac{2W}{W}
            \;=\; 2.
          \end{align}

    \item \emph{Tightness of the bound.}
          Suppose $s_i = -1$ for all $i$ with $w_i > 0$.
          Then $S = -1$ and $|s_i - S| = 0$ for all such $i$,
          giving $D = 0$.
          Now suppose instead $s_i = 1$ for all $i$ with
          $w_i > 0$ and a single agent $j$ has $s_j = -1$
          (with $w_j > 0$).  More precisely, the bound
          $D = 2$ is approached when half the weighted mass
          sits at $+1$ and half at $-1$, in which case
          $S = 0$ and $|s_i - S| = 1$ for all $i$, giving
          $D = 1 < 2$.
          The absolute maximum $D = 2$ would require
          $|s_i - S| = 2$ for all $i$, which requires
          simultaneously $s_i = 1$ and $S = -1$ (or vice
          versa) for every active agent, i.e.\ all signals
          equal $+1$ while the aggregate equals $-1$---a
          contradiction since $S = \sum \tilde{w}_i s_i$
          with $\tilde{w}_i \geq 0$ and $\sum \tilde{w}_i = 1$.
          Hence $D < 2$ in all realisable configurations,
          and the bound $D \leq 2$ is non-strict but tight
          in the limit sense.
  \end{enumerate}

  % --- Property 3 ---
  \item \textbf{Consensus characterisation ($D = 0 \Leftrightarrow
        s_i = S\ \forall i : w_i > 0$).}

  \begin{enumerate}
    \item \emph{($\Leftarrow$) Sufficiency.}
          Assume $s_i = S$ for all $i$ with $w_i > 0$.
          Then $|s_i - S| = 0$ for all such $i$, and
          for $i$ with $w_i = 0$ the term $w_i |s_i - S| = 0$
          regardless of $s_i$.  Therefore
          \begin{align}
            \sum_{i=1}^{N} w_i\,\bigl|s_i - S\bigr|
            &\;=\; 0,
          \end{align}
          so $D = 0/W = 0$.

    \item \emph{($\Rightarrow$) Necessity.}
          Assume $D = 0$.  Then
          \begin{align}
            \sum_{i=1}^{N} w_i\,\bigl|s_i - S\bigr|
            &\;=\; 0.
          \end{align}
          Each term satisfies $w_i \geq 0$ and
          $|s_i - S| \geq 0$, so the product
          $w_i |s_i - S| \geq 0$.
          A sum of non-negative reals equals zero if and only
          if every summand equals zero:
          \begin{align}
            w_i\,\bigl|s_i - S\bigr| \;=\; 0
            \quad \forall\, i \in \{1,\ldots,N\}.
          \end{align}
          For any $i$ with $w_i > 0$, this forces
          $|s_i - S| = 0$, i.e.\ $s_i = S$.

    \item \emph{Edge case: no active agent.}
          By Assumption~\ref{asm:weight-regularity},
          $W = \sum_{i=1}^{N} w_i > 0$, so at least one
          $w_j > 0$ exists.  The vacuous case $w_i = 0$
          for all $i$ is excluded by assumption.

    \item \emph{Edge case: single active agent.}
          If exactly one $w_j > 0$ and $w_i = 0$ for all
          $i \neq j$, then $S = s_j$ (since
          $\tilde{w}_j = 1$), so $|s_j - S| = 0$ and
          $D = 0$ automatically.  This is consistent with
          the consensus condition $s_j = S$.

    \item \emph{Conclusion.}
          The equivalence $D = 0 \Leftrightarrow s_i = S$
          for all $i$ with $w_i > 0$ is established in
          both directions, with all edge cases verified.
  \end{enumerate}

\end{enumerate}
\qed
\end{proof}

%------------------------------------------------------------
\begin{remark}[Semantic Significance of Disagreement]
\label{rem:disagreement-semantic}
The aggregate signal $S = 0$ admits two fundamentally different
interpretations that are indistinguishable from $S$ alone yet carry
entirely different risk implications.

\begin{enumerate}
  \item \textbf{Violent disagreement.}
        Suppose two agents of equal weight hold signals
        $s_1 = +0.9$ and $s_2 = -0.9$.  Then
        \begin{align}
          S
          &\;=\; \tfrac{1}{2}(+0.9) + \tfrac{1}{2}(-0.9)
           \;=\; 0,\\
          D
          &\;=\; \tfrac{1}{2}\,|{+0.9} - 0|
               + \tfrac{1}{2}\,|{-0.9} - 0|
           \;=\; 0.9.
        \end{align}
        The ensemble appears neutral ($S = 0$) yet the agents
        hold near-maximal opposing views.  The high value
        $D = 0.9$ exposes this latent risk.

  \item \textbf{Genuine neutrality.}
        Suppose the same two agents hold signals
        $s_1 = +0.05$ and $s_2 = -0.05$.  Then
        \begin{align}
          S
          &\;=\; \tfrac{1}{2}(+0.05) + \tfrac{1}{2}(-0.05)
           \;=\; 0,\\
          D
          &\;=\; \tfrac{1}{2}\,|{+0.05} - 0|
               + \tfrac{1}{2}\,|{-0.05} - 0|
           \;=\; 0.05.
        \end{align}
        Here the ensemble is neutral because both agents are
        genuinely uncertain.  The low value $D = 0.05$
        reflects this benign indifference.

  \item \textbf{Interpretive principle.}
        Both scenarios yield $S = 0$, yet their risk
        profiles differ sharply.  Relying solely on $S$
        conflates violent disagreement with genuine
        neutrality.  Including $D$ as a co-equal state
        variable (Definition~\ref{def:allocator-state})
        disambiguates the two cases.

  \item \textbf{Deployment implication.}
        High $D$ should reduce capital deployment even when
        $S$ appears favourable, because high disagreement
        is itself evidence of elevated model risk: the
        ensemble cannot agree on the sign of the trade,
        let alone its magnitude.
        Axiom~\ref{ax:disagreement-gate} formalises this
        precautionary gate.

  \item \textbf{Boundary validation.}
        The formula $D = 0.9$ in Case~1 satisfies
        $D \in [0, 2)$ as established in
        Proposition~\ref{prop:disagreement-props},
        Property~2.  The formula $D = 0.05$ in Case~2
        satisfies the consensus bound $D \to 0$
        as per Property~3 of the same proposition.
        No boundary condition is violated.
\end{enumerate}
\end{remark}

%------------------------------------------------------------
\begin{axiom}[Disagreement Deployment Gate]
\label{ax:disagreement-gate}
Let $D_{\max} \in (0,2)$ be a pre-configured threshold,
fixed prior to any deployment decision and
$\mathcal{F}_0$-measurable.
The following conditions are imposed.

\begin{enumerate}
  \item \textbf{Applicability condition.}
        The gate is active at time $t$ if and only if
        $D_t > D_{\max}$, where $D_t$ is the agent
        disagreement (Definition~\ref{def:disagreement})
        evaluated at time $t$.

  \item \textbf{Capital fraction cap.}
        When $D_t > D_{\max}$, the deployable capital
        fraction is capped at $\kappa_D \in [0,1]$ defined
        by
        \begin{align}
          \kappa_D
          &\;:=\;
          \max\!\left(0,\;
            1 - \frac{D_t - D_{\max}}{2 - D_{\max}}
          \right).
          \label{eq:disagreement-gate}
        \end{align}

  \item \textbf{Verification of range $\kappa_D \in [0,1]$.}
        \begin{enumerate}
          \item \emph{Upper bound.}
                When $D_t = D_{\max}$, the fraction equals
                \begin{align}
                  \kappa_D
                  &\;=\; \max\!\left(0,\;
                    1 - \frac{D_{\max} - D_{\max}}{2-D_{\max}}
                  \right)
                  \;=\; \max(0,\, 1)
                  \;=\; 1.
                \end{align}
                Hence $\kappa_D \leq 1$ with equality at
                the activation boundary.

          \item \emph{Monotone decrease.}
                The inner expression
                $1 - (D_t - D_{\max})/(2 - D_{\max})$
                is strictly decreasing in $D_t$ for
                $D_t > D_{\max}$, since
                $(2 - D_{\max}) > 0$ by assumption
                $D_{\max} \in (0,2)$.

          \item \emph{Lower bound.}
                As $D_t \to 2^{-}$,
                \begin{align}
                  1 - \frac{D_t - D_{\max}}{2 - D_{\max}}
                  &\;\to\;
                  1 - \frac{2 - D_{\max}}{2 - D_{\max}}
                  \;=\; 0.
                \end{align}
                The $\max(0,\,\cdot)$ operator ensures
                $\kappa_D \geq 0$ for all $D_t \in [0,2)$.

          \item \emph{Conclusion.}
                Combining the above,
                $\kappa_D \in [0,1]$ for all
                $D_t \in [D_{\max}, 2)$, and
                $\kappa_D = 1$ for $D_t \leq D_{\max}$.
        \end{enumerate}

  \item \textbf{Boundary cases.}
        \begin{enumerate}
          \item \emph{$D_{\max} \to 0^{+}$.}
                The gate activates for any positive
                disagreement; $\kappa_D \to 1 - D_t/2$,
                which is maximally conservative.

          \item \emph{$D_{\max} \to 2^{-}$.}
                The gate never activates in practice;
                $\kappa_D \equiv 1$ unless $D_t$ approaches
                the theoretical maximum.

          \item \emph{$D_t = 0$.}
                Full consensus; $D_t \leq D_{\max}$ so
                $\kappa_D = 1$ and no cap is applied.

          \item \emph{$D_t \geq 2$.}
                Excluded by
                Proposition~\ref{prop:disagreement-props},
                Property~2, which establishes $D < 2$
                strictly for all realisable configurations.
        \end{enumerate}

  \item \textbf{Deployment scaling rule.}
        All capital deployment actions $\delta_t$ at time
        $t$ are replaced by
        \begin{align}
          \delta_t^{(\mathrm{gated})}
          &\;:=\; \kappa_D \cdot \delta_t,
          \label{eq:gated-deployment}
        \end{align}
        ensuring the scaled action satisfies
        $|\delta_t^{(\mathrm{gated})}| \leq |\delta_t|$,
        with equality if and only if $D_t \leq D_{\max}$.

  \item \textbf{Measurability.}
        Since $D_t$ is $\mathcal{F}_t$-measurable
        (Definition~\ref{def:disagreement}) and $D_{\max}$
        is $\mathcal{F}_0$-measurable, $\kappa_D$ is
        $\mathcal{F}_t$-measurable and hence admissible
        as a scaling factor in any $\mathcal{F}_t$-adapted
        strategy.
\end{enumerate}
\end{axiom}

%------------------------------------------------------------
\section{Allocator State Space}

\begin{definition}[Allocator State]
\label{def:allocator-state}
At each time $t$, the master capital allocator reasons over
the following multidimensional state vector.

\begin{enumerate}
  \item \textbf{State vector.}
        \begin{align}
          \calX_t
          &\;=\;
          \bigl(\,
            S_t,\;\;
            C_t,\;\;
            U_t,\;\;
            D_t,\;\;
            R_{\mathrm{port},t},\;\;
            H_t,\;\;
            V_{\mathrm{vol},t},\;\;
            O_t,\;\;
            \Delta t
          \,\bigr).
          \label{eq:allocator-state}
        \end{align}

  \item \textbf{Measurability requirement.}
        Each component of $\calX_t$ is
        $\mathcal{F}_t$-measurable, ensuring that all
        allocation decisions based on $\calX_t$ are
        non-anticipating and causal.

  \item \textbf{Component definitions.}
        The precise meaning of each symbol is given in
        Tables~\ref{tab:allocator-state-A}
        and~\ref{tab:allocator-state-B}.

  \item \textbf{Completeness.}
        The state $\calX_t$ is sufficient for the
        master allocator to evaluate
        (i)~the directional strength and agreement of the
        ensemble, (ii)~current portfolio risk and hedging
        needs, (iii)~market regime and volatility, and
        (iv)~the applicable investment horizon.
        No additional hidden state is assumed.
\end{enumerate}
\end{definition}

%------------------------------------------------------------
\begin{table}[H]
\centering
\caption{Allocator state components --- Part A: ensemble
         and disagreement signals.}
\label{tab:allocator-state-A}
\small
\begin{tabularx}{\linewidth}{@{} c X @{}}
\toprule
\textbf{Symbol} & \textbf{Definition and measurability} \\
\midrule
$S_t$ &
  Ensemble directional score
  (Definition~\ref{def:ensemble-signal}).
  $\mathcal{F}_t$-measurable by construction;
  $S_t \in [-1,+1]$. \\[4pt]
$C_t$ &
  Ensemble confidence, defined as the weighted average
  of individual agent confidences $c_{i,t} \in [0,1]$:
  \begin{align*}
    C_t \;:=\;
    \frac{\displaystyle\sum_{i=1}^{N} w_i\,c_{i,t}}
         {\displaystyle\sum_{i=1}^{N} w_i}.
  \end{align*}
  $\mathcal{F}_t$-measurable; $C_t \in [0,1]$. \\[4pt]
$U_t$ &
  Ensemble uncertainty, defined as the weighted average
  of individual agent uncertainties $u_{i,t} \in [0,1]$:
  \begin{align*}
    U_t \;:=\;
    \frac{\displaystyle\sum_{i=1}^{N} w_i\,u_{i,t}}
         {\displaystyle\sum_{i=1}^{N} w_i}.
  \end{align*}
  $\mathcal{F}_t$-measurable; $U_t \in [0,1]$.
  Note: $C_t + U_t \leq 1$ is not required in general
  unless explicitly assumed by a sub-model. \\[4pt]
$D_t$ &
  Agent disagreement
  (Definition~\ref{def:disagreement}).
  $\mathcal{F}_t$-measurable;
  $D_t \in [0,2)$ by
  Proposition~\ref{prop:disagreement-props}. \\
\bottomrule
\end{tabularx}
\end{table}

%------------------------------------------------------------
\begin{table}[H]
\centering
\caption{Allocator state components --- Part B: portfolio,
         market, and horizon state.}
\label{tab:allocator-state-B}
\small
\begin{tabularx}{\linewidth}{@{} c X @{}}
\toprule
\textbf{Symbol} & \textbf{Definition and measurability} \\
\midrule
$R_{\mathrm{port},t}$ &
  Current portfolio risk level, e.g.\ portfolio value-at-risk
  or normalised margin utilisation.
  $\mathcal{F}_t$-measurable; takes values in
  $[0, R_{\max}]$ for some pre-specified $R_{\max} > 0$. \\[4pt]
$H_t$ &
  Hedging requirements: a scalar or vector summarising the
  net delta/gamma/vega exposure that must be offset.
  $\mathcal{F}_t$-measurable. \\[4pt]
$V_{\mathrm{vol},t}$ &
  Volatility and regime information, e.g.\ a composite of
  realised volatility, implied volatility index, and a
  discrete regime label (trending, mean-reverting, crisis).
  $\mathcal{F}_t$-measurable. \\[4pt]
$O_t$ &
  Options-related state: implied volatility surface summary,
  skew, and term-structure slope as of time $t$.
  $\mathcal{F}_t$-measurable. \\[4pt]
$\Delta t$ &
  Relevant investment horizon for the current decision
  epoch, expressed in units consistent with the underlying
  time grid.
  Deterministic or $\mathcal{F}_t$-measurable depending on
  the scheduling rule. \\
\bottomrule
\end{tabularx}
\end{table}

\begin{table}[H]
\centering
\caption{Allocator state components.}
\label{tab:allocator-state-components}
\small
\begin{tabularx}{\linewidth}{@{} c X @{}}
\toprule
\textbf{Symbol} & \textbf{Meaning} \\
\midrule
$S$ & Ensemble directional score (Definition~\ref{def:ensemble-signal}). \\
$C$ & Ensemble confidence: $C = (\sum_i w_i c_i)/(\sum_i w_i)$. \\
$U$ & Ensemble uncertainty: $U = (\sum_i w_i u_i)/(\sum_i w_i)$. \\
$D$ & Agent disagreement (Definition~\ref{def:disagreement}). \\
$R_{\mathrm{port}}$ & Current portfolio risk level. \\
$H$ & Hedging requirements. \\
$V_{\mathrm{vol}}$ & Volatility and regime information. \\
$O$ & Options-related state (skew, implied vol surface). \\
$\Delta t$ & Relevant investment horizon. \\
\bottomrule
\end{tabularx}
\end{table}

%------------------------------------------------------------
\subsection{Formal Domain and Measurability of State Components}
%------------------------------------------------------------

\noindent
We now establish, with full rigour, the domain, measurability, and
admissibility of every component of $\calX_t$.  No component is
assumed well-defined without explicit justification.

\begin{assumption}[Filtered Probability Space]
\label{asm:filtered-space}
Throughout this section we fix a complete filtered probability space
$(\Omega,\calF,\{\calF_t\}_{t\geq 0},\Prob)$ satisfying the usual
conditions (right-continuity and completeness of the filtration).
All random variables are defined on this space.
\end{assumption}

\begin{proposition}[Domain and Measurability of $S$]
\label{prop:S-domain}
Under Definitions~\ref{def:ensemble-signal} and
\ref{def:agent-score}, the ensemble directional score satisfies:
\begin{enumerate}
  \item \textbf{Domain.} $S \in [-1,1]$ almost surely.
  \item \textbf{Measurability.} $S$ is $\calF_t$-measurable.
  \item \textbf{Boundary behaviour.}
    $S = +1$ (resp.\ $-1$) occurs only when every constituent agent
    assigns maximal positive (resp.\ negative) directional conviction
    simultaneously, which has measure zero under Assumption~\ref{asm:return-process}.
\end{enumerate}
\end{proposition}

\begin{proof}
We address each claim in order, validating every assumption and
treating all edge cases explicitly.

\begin{enumerate}

  %-------- (1) Domain ----------------------------------------
  \item \textbf{Domain: $S \in [-1,1]$ a.s.}

  \begin{enumerate}
    \item \emph{Assumption validation.}
      By Definition~\ref{def:agent-score}, for every agent
      $i \in \{1,\ldots,n\}$:
      \begin{align}
        s_i &\in [-1,1], \label{eq:si-domain}\\
        w_i &\geq 0,      \label{eq:wi-nonneg}\\
        \sum_{i=1}^{n} w_i &= 1. \label{eq:wi-sum}
      \end{align}
      These three conditions hold by definition and are
      treated as verified.

    \item \emph{Lower bound.}
      Using \eqref{eq:si-domain} and \eqref{eq:wi-nonneg}:
      \begin{align}
        S
          &= \sum_{i=1}^{n} w_i s_i
           \;\geq\;
           \sum_{i=1}^{n} w_i \cdot (-1)
           \;=\;
           -\sum_{i=1}^{n} w_i
           \;=\;
           -1,
      \end{align}
      where the final equality uses \eqref{eq:wi-sum}.

    \item \emph{Upper bound.}
      By an identical argument using $s_i \leq +1$:
      \begin{align}
        S
          &= \sum_{i=1}^{n} w_i s_i
           \;\leq\;
           \sum_{i=1}^{n} w_i \cdot (+1)
           \;=\;
           +\sum_{i=1}^{n} w_i
           \;=\;
           +1.
      \end{align}

    \item \emph{Edge case: $n = 0$ (no active agents).}
      If no agents are active the sum $\sum_i w_i s_i$ is
      empty.  By convention an empty weighted sum equals zero,
      so $S = 0 \in [-1,1]$.  This edge case is therefore
      covered.

    \item \emph{Conclusion.}
      Combining the lower and upper bounds:
      \begin{align}
        -1 \;\leq\; S \;\leq\; +1 \quad \Prob\text{-a.s.}
      \end{align}
  \end{enumerate}

  %-------- (2) Measurability ---------------------------------
  \item \textbf{Measurability: $S$ is $\calF_t$-measurable.}

  \begin{enumerate}
    \item \emph{Measurability of scores.}
      By Definition~\ref{def:agent-score}, each agent score
      $s_i \colon \Omega \to [-1,1]$ is
      $\calF_t$-measurable, i.e.,
      \begin{align}
        s_i^{-1}(B) \in \calF_t
        \quad
        \text{for every Borel set } B \subseteq \mathbb{R}.
      \end{align}

    \item \emph{Measurability of weights.}
      Each weight $w_i$ is either a strictly positive
      deterministic constant or an $\calF_t$-measurable
      random variable (depending on whether adaptive
      weighting is employed).  In either case $w_i$ is
      $\calF_t$-measurable.

    \item \emph{Closure under finite operations.}
      The product $w_i s_i$ of two $\calF_t$-measurable
      functions is $\calF_t$-measurable (products of
      measurable functions are measurable).  A finite sum of
      $\calF_t$-measurable functions is $\calF_t$-measurable.
      Hence:
      \begin{align}
        S = \sum_{i=1}^{n} w_i s_i
        \quad
        \text{is } \calF_t\text{-measurable.}
      \end{align}

    \item \emph{Edge case: constant weights.}
      If every $w_i$ is a deterministic constant, the
      argument reduces to: a finite linear combination of
      $\calF_t$-measurable functions with deterministic
      coefficients is $\calF_t$-measurable.  This is
      strictly weaker and the conclusion still holds.
  \end{enumerate}

  %-------- (3) Boundary behaviour ----------------------------
  \item \textbf{Boundary behaviour: $\Prob(S = \pm 1) = 0$.}

  \begin{enumerate}
    \item \emph{Characterisation of $\{S = +1\}$.}
      Suppose $S = +1$.  Since $w_i \geq 0$,
      $\sum_i w_i = 1$, and $s_i \leq 1$, we have:
      \begin{align}
        1
          &= S
           = \sum_{i=1}^{n} w_i s_i
           \leq \sum_{i=1}^{n} w_i \cdot 1
           = 1.
      \end{align}
      Equality holds throughout if and only if
      $s_i = 1$ for every $i$ with $w_i > 0$.  Denote the
      index set of active agents by
      $\mathcal{A} := \{i : w_i > 0\}$.  Then:
      \begin{align}
        \{S = +1\}
          &= \bigcap_{i \in \mathcal{A}} \{s_i = +1\}.
        \label{eq:S-plus1-event}
      \end{align}

    \item \emph{Measure-zero argument.}
      Under Assumption~\ref{asm:return-process}, each
      agent's belief distribution has a continuously
      distributed component, so for every $i \in \mathcal{A}$:
      \begin{align}
        \Prob(s_i = +1) = 0.
      \end{align}
      By monotonicity of probability:
      \begin{align}
        \Prob(S = +1)
          &= \Prob\!\left(
               \bigcap_{i \in \mathcal{A}} \{s_i = +1\}
             \right)
           \;\leq\;
           \Prob(s_j = +1)
           \;=\; 0,
      \end{align}
      for any fixed $j \in \mathcal{A}$.  Hence
      $\Prob(S = +1) = 0$.

    \item \emph{Symmetric argument for $\{S = -1\}$.}
      By an identical argument with $s_i = -1$ replacing
      $s_i = +1$ throughout:
      \begin{align}
        \{S = -1\}
          &= \bigcap_{i \in \mathcal{A}} \{s_i = -1\},\\
        \Prob(S = -1)
          &\leq \Prob(s_j = -1)
           = 0,
      \end{align}
      for any $j \in \mathcal{A}$.

    \item \emph{Edge case: $\mathcal{A} = \emptyset$.}
      If no agent is active ($w_i = 0$ for all $i$), then
      $S = 0$ by the empty-sum convention, so
      $\{S = \pm 1\} = \emptyset$ and the claim holds
      trivially with probability zero.

    \item \emph{Conclusion.}
      \begin{align}
        \Prob(S = +1) = 0
        \quad \text{and} \quad
        \Prob(S = -1) = 0. \qedhere
      \end{align}
  \end{enumerate}

\end{enumerate}
\end{proof}

\begin{proposition}[Domain and Measurability of $C$ and $U$]
\label{prop:CU-domain}
Let $\{c_i\}_{i=1}^{n}$ and $\{u_i\}_{i=1}^{n}$ be the
per-agent confidence and uncertainty outputs, and let
$\{w_i\}_{i=1}^{n}$ be the corresponding non-negative weights
as in Definition~\ref{def:ensemble-signal}.  Define the
index set of active agents
\begin{align}
  \mathcal{A} &:= \{i \in \{1,\ldots,n\} : w_i > 0\}.
\end{align}
The ensemble confidence $C$ and ensemble uncertainty $U$ are
defined by
\begin{align}
  C &:= \frac{\displaystyle\sum_{i=1}^{n} w_i\, c_i}
             {\displaystyle\sum_{i=1}^{n} w_i},
  &
  U &:= \frac{\displaystyle\sum_{i=1}^{n} w_i\, u_i}
             {\displaystyle\sum_{i=1}^{n} w_i}.
\end{align}
Under Assumptions~\ref{asm:weights}--\ref{asm:agent-outputs}
the following hold:
\begin{enumerate}
  \item \textbf{Well-definedness of denominators.}
    $\displaystyle W := \sum_{i=1}^{n} w_i > 0$ whenever
    $\mathcal{A} \neq \emptyset$.
  \item \textbf{Domain.}
    $C \in [0,1]$ and $U \in [0,1]$ $\Prob$-a.s.
  \item \textbf{Measurability.}
    Both $C$ and $U$ are $\calF_t$-measurable.
  \item \textbf{Lower boundary.}
    $C = 0$ $\Prob$-a.s.\ if and only if $c_i = 0$
    $\Prob$-a.s.\ for every $i \in \mathcal{A}$.
    The same characterisation holds for $U$.
  \item \textbf{Upper boundary.}
    $C = 1$ $\Prob$-a.s.\ if and only if $c_i = 1$
    $\Prob$-a.s.\ for every $i \in \mathcal{A}$.
    The same characterisation holds for $U$.
  \item \textbf{Edge case: $\mathcal{A} = \emptyset$.}
    If no agent is active, then $C$ and $U$ are set to zero
    by the empty-sum convention and the claims hold trivially.
\end{enumerate}
\end{proposition}

\begin{proof}
We prove each item in turn, making every assumption and
algebraic step explicit.

\begin{enumerate}

  %-------- Item 1: Well-definedness --------------------------
  \item \textbf{Well-definedness of denominators.}

  \begin{enumerate}
    \item \emph{Assumption validation.}
      By Assumption~\ref{asm:weights}, each weight satisfies
      $w_i \geq 0$ for all $i \in \{1,\ldots,n\}$.
      By Assumption~\ref{asm:active-agent}, at least one
      agent is active at every time step, i.e.\
      $\mathcal{A} \neq \emptyset$.  Fix any
      $j \in \mathcal{A}$; then $w_j > 0$.

    \item \emph{Lower bound on $W$.}
      Since $w_i \geq 0$ for all $i$ and $w_j > 0$:
      \begin{align}
        W
          &= \sum_{i=1}^{n} w_i
           = w_j + \sum_{i \neq j} w_i
           \;\geq\;
           w_j + 0
           \;=\; w_j
           \;>\; 0.
      \end{align}

    \item \emph{Defining $w_{\min}$.}
      Let $w_{\min} := \min_{i \in \mathcal{A}} w_i > 0$.
      Then $W \geq |\mathcal{A}|\, w_{\min} > 0$, which
      uniformly bounds $W$ away from zero and precludes
      division by zero in the definitions of $C$ and $U$.
  \end{enumerate}

  %-------- Item 2: Domain ------------------------------------
  \item \textbf{Domain: $C \in [0,1]$ and $U \in [0,1]$.}

  \begin{enumerate}
    \item \emph{Assumption validation.}
      By Assumption~\ref{asm:agent-outputs},
      $c_i \in [0,1]$ and $u_i \in [0,1]$
      $\Prob$-a.s.\ for every $i$.  By item~1,
      $W > 0$.

    \item \emph{Lower bound for $C$.}
      Since $c_i \geq 0$ and $w_i \geq 0$:
      \begin{align}
        \sum_{i=1}^{n} w_i\, c_i
          &\;\geq\;
           \sum_{i=1}^{n} w_i \cdot 0
           \;=\; 0,
      \end{align}
      and therefore $C = \bigl(\sum_i w_i c_i\bigr)/W \geq 0$.

    \item \emph{Upper bound for $C$.}
      Since $c_i \leq 1$ and $w_i \geq 0$:
      \begin{align}
        \sum_{i=1}^{n} w_i\, c_i
          &\;\leq\;
           \sum_{i=1}^{n} w_i \cdot 1
           \;=\; W,
      \end{align}
      and therefore $C = \bigl(\sum_i w_i c_i\bigr)/W \leq 1$.

    \item \emph{Combining bounds for $C$.}
      \begin{align}
        0
          &\;\leq\;
           C
           \;=\;
           \frac{\sum_{i=1}^{n} w_i\, c_i}{W}
           \;\leq\; 1,
           \quad \Prob\text{-a.s.}
      \end{align}

    \item \emph{Identical argument for $U$.}
      Replacing $c_i$ by $u_i$ in steps (b)--(d) above
      yields, by the same chain of inequalities:
      \begin{align}
        0
          &\;\leq\;
           U
           \;=\;
           \frac{\sum_{i=1}^{n} w_i\, u_i}{W}
           \;\leq\; 1,
           \quad \Prob\text{-a.s.}
      \end{align}
  \end{enumerate}

  %-------- Item 3: Measurability -----------------------------
  \item \textbf{Measurability.}

  \begin{enumerate}
    \item \emph{Measurability of individual outputs.}
      By Assumption~\ref{asm:agent-outputs}, each $c_i$ and
      $u_i$ is $\calF_t$-measurable, i.e.\
      $c_i^{-1}(B), u_i^{-1}(B) \in \calF_t$ for every
      Borel set $B \subseteq \mathbb{R}$.

    \item \emph{Measurability of weights.}
      By Assumption~\ref{asm:weights}, each $w_i$ is either
      a deterministic constant (hence trivially
      $\calF_t$-measurable) or an $\calF_t$-measurable
      random variable under the adaptive weighting scheme.
      In both cases $w_i$ is $\calF_t$-measurable.

    \item \emph{Closure of $\calF_t$-measurability under
      products.}
      For each $i$, the pointwise product $w_i\, c_i$ is
      $\calF_t$-measurable because the product of two
      $\calF_t$-measurable functions is
      $\calF_t$-measurable.  Likewise $w_i\, u_i$ is
      $\calF_t$-measurable.

    \item \emph{Closure under finite sums.}
      A finite sum of $\calF_t$-measurable functions is
      $\calF_t$-measurable; hence
      \begin{align}
        \sum_{i=1}^{n} w_i\, c_i
          \quad\text{and}\quad
        \sum_{i=1}^{n} w_i\, u_i
          \quad\text{and}\quad
        W = \sum_{i=1}^{n} w_i
      \end{align}
      are all $\calF_t$-measurable.

    \item \emph{Closure under division with a.s.\ non-zero
      denominator.}
      By item~1, $W > 0$ $\Prob$-a.s.  The ratio of two
      $\calF_t$-measurable functions whose denominator is
      $\Prob$-a.s.\ strictly positive is itself
      $\calF_t$-measurable.  Therefore:
      \begin{align}
        C
          &= \frac{\sum_{i=1}^{n} w_i\, c_i}{W}
          \quad\text{is } \calF_t\text{-measurable},\\
        U
          &= \frac{\sum_{i=1}^{n} w_i\, u_i}{W}
          \quad\text{is } \calF_t\text{-measurable}.
      \end{align}
  \end{enumerate}

  %-------- Item 4: Lower boundary ----------------------------
  \item \textbf{Lower boundary characterisation.}

  \begin{enumerate}
    \item \emph{($\Rightarrow$) $C = 0$ implies $c_i = 0$
      for all $i \in \mathcal{A}$.}
      Suppose $C = 0$ $\Prob$-a.s.  Then:
      \begin{align}
        0
          &= W \cdot C
           = \sum_{i \in \mathcal{A}} w_i\, c_i.
      \end{align}
      Each summand satisfies $w_i\, c_i \geq 0$ (since
      $w_i > 0$ and $c_i \geq 0$).  A non-negative sum
      equals zero if and only if every summand equals zero.
      Since $w_i > 0$ for $i \in \mathcal{A}$, we conclude
      $c_i = 0$ $\Prob$-a.s.\ for every
      $i \in \mathcal{A}$.

    \item \emph{($\Leftarrow$) $c_i = 0$ for all
      $i \in \mathcal{A}$ implies $C = 0$.}
      If $c_i = 0$ for every $i \in \mathcal{A}$, then:
      \begin{align}
        C
          &= \frac{\sum_{i \in \mathcal{A}} w_i \cdot 0}{W}
           = \frac{0}{W}
           = 0.
      \end{align}

    \item \emph{Same characterisation for $U$.}
      Replacing $c_i$ by $u_i$ and $C$ by $U$ in
      steps (a)--(b) above yields the identical equivalence
      $U = 0 \Leftrightarrow u_i = 0\;\forall\, i \in
      \mathcal{A}$.
  \end{enumerate}

  %-------- Item 5: Upper boundary ----------------------------
  \item \textbf{Upper boundary characterisation.}

  \begin{enumerate}
    \item \emph{($\Rightarrow$) $C = 1$ implies $c_i = 1$
      for all $i \in \mathcal{A}$.}
      Suppose $C = 1$ $\Prob$-a.s.  Then:
      \begin{align}
        W
          &= W \cdot 1
           = W \cdot C
           = \sum_{i \in \mathcal{A}} w_i\, c_i.
      \end{align}
      Rearranging:
      \begin{align}
        0
          &= W - \sum_{i \in \mathcal{A}} w_i\, c_i
           = \sum_{i \in \mathcal{A}} w_i (1 - c_i).
      \end{align}
      Each term satisfies $w_i(1 - c_i) \geq 0$ because
      $w_i > 0$ and $c_i \leq 1$.  A non-negative sum
      equals zero if and only if every summand equals zero.
      Hence $1 - c_i = 0$, i.e.\ $c_i = 1$ $\Prob$-a.s.\
      for every $i \in \mathcal{A}$.

    \item \emph{($\Leftarrow$) $c_i = 1$ for all
      $i \in \mathcal{A}$ implies $C = 1$.}
      If $c_i = 1$ for every $i \in \mathcal{A}$, then:
      \begin{align}
        C
          &= \frac{\sum_{i \in \mathcal{A}} w_i \cdot 1}{W}
           = \frac{W}{W}
           = 1.
      \end{align}

    \item \emph{Same characterisation for $U$.}
      Replacing $c_i$ by $u_i$ and $C$ by $U$ in
      steps (a)--(b) above yields the identical equivalence
      $U = 1 \Leftrightarrow u_i = 1\;\forall\, i \in
      \mathcal{A}$.
  \end{enumerate}

  %-------- Item 6: Edge case ---------------------------------
  \item \textbf{Edge case: $\mathcal{A} = \emptyset$.}

  \begin{enumerate}
    \item \emph{Empty-sum convention.}
      If $\mathcal{A} = \emptyset$, then $w_i = 0$ for all
      $i$, and by the empty-sum convention:
      \begin{align}
        \sum_{i=1}^{n} w_i\, c_i &= 0,
        &
        \sum_{i=1}^{n} w_i\, u_i &= 0,
        &
        W &= 0.
      \end{align}

    \item \emph{Convention for $C$ and $U$.}
      In this degenerate case, the system sets $C := 0$ and
      $U := 0$ by definition (no agent contributes any
      signal).  The domain claims $C \in [0,1]$ and
      $U \in [0,1]$ hold trivially, and both $C$ and $U$
      are trivially $\calF_t$-measurable as constants.
      \qedhere
  \end{enumerate}

\end{enumerate}
\end{proof}

\begin{proposition}[Domain and Measurability of $D$]
\label{prop:D-domain}
Let $n \geq 1$ be the number of active agents,
$\mathcal{A} \subseteq \{1,\ldots,n\}$ the active index set,
$s_i \in [-1,1]$ the directional score of agent $i$,
and $w_i > 0$ the weight of agent $i$ for $i \in \mathcal{A}$,
with $W := \sum_{i \in \mathcal{A}} w_i > 0$.
The weighted mean score and disagreement metric are defined as:
\begin{align}
  \bar{s}
    &:= \frac{1}{W}\sum_{i \in \mathcal{A}} w_i\, s_i,
  \\[4pt]
  D
    &:= \sqrt{\,\frac{1}{W}\sum_{i \in \mathcal{A}}
        w_i\,(s_i - \bar{s})^2\,}.
\end{align}
The agent disagreement metric $D$
(Definition~\ref{def:disagreement}) satisfies:
\begin{enumerate}
  \item \textbf{Domain.} $D \in [0,2]$ $\Prob$-a.s.
  \item \textbf{Measurability.} $D$ is $\calF_t$-measurable.
  \item \textbf{Lower boundary characterisation.}
    \begin{enumerate}
      \item $D = 0$ $\Prob$-a.s.\ if and only if
        $s_i = \bar{s}$ $\Prob$-a.s.\ for every
        $i \in \mathcal{A}$, i.e.\ all active agents
        produce identical directional scores.
      \item In particular, if $|\mathcal{A}| = 1$, then
        $D = 0$ trivially since there is no disagreement.
    \end{enumerate}
  \item \textbf{Upper boundary characterisation.}
    \begin{enumerate}
      \item $D \leq 2$ always, with the bound attained
        $\Prob$-a.s.\ if and only if the weighted agent
        mass is split with half assigning $s_i = +1$ and
        half assigning $s_i = -1$, i.e.\
        maximum polarisation.
      \item For a two-agent balanced split
        ($w_1 = w_2$, $s_1 = +1$, $s_2 = -1$):
        $\bar{s} = 0$ and
        $D = \sqrt{\tfrac{1}{2}(1)^2 + \tfrac{1}{2}(1)^2}
           = 1 < 2$.
        The global maximum $D = 2$ requires a specific
        multi-agent polarisation configuration detailed
        in item~4(b) of the proof below.
    \end{enumerate}
  \item \textbf{Interaction with disagreement gate.}
    Under Axiom~\ref{ax:disagreement-gate}, with threshold
    $D_{\max} \in [0,2)$, the gate coefficient satisfies
    $\kappa_D \in [0,1]$ for all $D \in [0,2]$.
\end{enumerate}
\end{proposition}

\begin{proof}
We verify each claim in order, validating all assumptions,
boundary conditions, and edge cases at every step.

\begin{enumerate}

  %-------- Item 1: Domain ------------------------------------
  \item \textbf{Domain: $D \in [0,2]$ $\Prob$-a.s.}

  \begin{enumerate}
    \item \emph{Non-negativity ($D \geq 0$).}
      The quantity inside the square root is a weighted sum
      of squared terms:
      \begin{align}
        \frac{1}{W}\sum_{i \in \mathcal{A}} w_i\,(s_i - \bar{s})^2
          &\geq 0,
      \end{align}
      since $w_i > 0$, $W > 0$, and $(s_i - \bar{s})^2 \geq 0$
      for every $i$.  Hence $D = \sqrt{(\cdot)} \geq 0$.

    \item \emph{Upper bound ($D \leq 2$).}
      Since $s_i \in [-1,1]$ for all $i \in \mathcal{A}$
      (Assumption~\ref{asm:score-domain}), and $\bar{s}$
      is a convex combination of these scores, we have
      $\bar{s} \in [-1,1]$.  Therefore:
      \begin{align}
        (s_i - \bar{s})^2
          &\leq \bigl(\max_{i}s_i - \min_{i}s_i\bigr)^2
           \leq \bigl(1 - (-1)\bigr)^2
           = 4.
      \end{align}
      Substituting into the definition of $D$:
      \begin{align}
        D^2
          &= \frac{1}{W}\sum_{i \in \mathcal{A}} w_i\,(s_i - \bar{s})^2
           \leq \frac{1}{W}\sum_{i \in \mathcal{A}} w_i \cdot 4
           = \frac{4W}{W}
           = 4.
      \end{align}
      Taking square roots: $D \leq \sqrt{4} = 2$.

    \item \emph{Edge case: $|\mathcal{A}| = 1$.}
      If exactly one agent is active, $\bar{s} = s_1$ and
      $D = \sqrt{w_1(s_1 - s_1)^2 / w_1} = 0 \in [0,2]$.

    \item \emph{Combining:} $0 \leq D \leq 2$ $\Prob$-a.s.
  \end{enumerate}

  %-------- Item 2: Measurability -----------------------------
  \item \textbf{Measurability of $D$.}

  \begin{enumerate}
    \item \emph{Measurability of inputs.}
      By Assumption~\ref{asm:return-process} and the
      definition of the agent scoring model, each score
      $s_i$ and weight $w_i$ is $\calF_t$-measurable for
      every $i \in \mathcal{A}$.

    \item \emph{Measurability of $W$.}
      $W = \sum_{i \in \mathcal{A}} w_i$ is a finite sum of
      $\calF_t$-measurable functions, hence $\calF_t$-measurable.
      Since $W > 0$ $\Prob$-a.s.\ (Assumption~\ref{asm:active-agents}),
      $1/W$ is also $\calF_t$-measurable.

    \item \emph{Measurability of $\bar{s}$.}
      \begin{align}
        \bar{s}
          &= \frac{1}{W}\sum_{i \in \mathcal{A}} w_i\, s_i
      \end{align}
      is a finite linear combination of products of
      $\calF_t$-measurable functions, hence
      $\calF_t$-measurable.

    \item \emph{Measurability of $D$.}
      Each $(s_i - \bar{s})^2$ is $\calF_t$-measurable
      (difference and square of measurable functions).
      The finite weighted sum and subsequent square root
      (applied to a non-negative measurable function)
      preserve $\calF_t$-measurability.  Hence $D$ is
      $\calF_t$-measurable.
  \end{enumerate}

  %-------- Item 3: Lower boundary ----------------------------
  \item \textbf{Lower boundary: $D = 0$.}

  \begin{enumerate}
    \item \emph{($\Rightarrow$) $D = 0$ implies
      $s_i = \bar{s}$ for all $i \in \mathcal{A}$.}
      Suppose $D = 0$ $\Prob$-a.s.  Then $D^2 = 0$, so:
      \begin{align}
        0
          &= \frac{1}{W}\sum_{i \in \mathcal{A}}
              w_i\,(s_i - \bar{s})^2.
      \end{align}
      Since $W > 0$ and each $w_i > 0$, every term
      $w_i(s_i - \bar{s})^2 \geq 0$ in a non-negative
      sum that equals zero; therefore each term is
      individually zero.  Since $w_i > 0$, we conclude
      $(s_i - \bar{s})^2 = 0$, i.e.\
      $s_i = \bar{s}$ $\Prob$-a.s.\ for every
      $i \in \mathcal{A}$.

    \item \emph{($\Leftarrow$) $s_i = \bar{s}$ for all
      $i \in \mathcal{A}$ implies $D = 0$.}
      If $s_i = c$ for some constant $c \in [-1,1]$ and
      all $i \in \mathcal{A}$, then:
      \begin{align}
        \bar{s}
          &= \frac{1}{W}\sum_{i \in \mathcal{A}} w_i\, c
           = \frac{cW}{W}
           = c,
        \\[4pt]
        D^2
          &= \frac{1}{W}\sum_{i \in \mathcal{A}}
              w_i\,(c - c)^2
           = 0,
        \\[4pt]
        D &= 0.
      \end{align}

    \item \emph{Special case $|\mathcal{A}| = 1$.}
      With a single active agent, $\bar{s} = s_1$ trivially,
      so $D = 0$ regardless of the value of $s_1 \in [-1,1]$.
  \end{enumerate}

  %-------- Item 4: Upper boundary ----------------------------
  \item \textbf{Upper boundary: $D = 2$.}

  \begin{enumerate}
    \item \emph{Necessary condition for $D = 2$.}
      $D = 2$ $\Prob$-a.s.\ iff $D^2 = 4$, i.e.:
      \begin{align}
        \frac{1}{W}\sum_{i \in \mathcal{A}}
          w_i\,(s_i - \bar{s})^2
          &= 4.
      \end{align}
      Since $(s_i - \bar{s})^2 \leq 4$ for all $i$
      (by item~1(b)), and the weighted average of
      quantities each at most $4$ equals $4$, every
      term must attain the maximum:
      \begin{align}
        w_i\,(s_i - \bar{s})^2 &= 4\,w_i
        \quad \forall\, i \in \mathcal{A}.
      \end{align}
      Since $w_i > 0$, this requires
      $(s_i - \bar{s})^2 = 4$, i.e.\
      $s_i - \bar{s} \in \{-2, +2\}$ for every
      $i \in \mathcal{A}$.

    \item \emph{Sufficient condition: maximum polarisation.}
      Since $s_i \in [-1,1]$ and $\bar{s} \in [-1,1]$,
      the only admissible solutions to
      $|s_i - \bar{s}| = 2$ are:
      \begin{align}
        s_i = +1,\;\bar{s} = -1
          &\quad\text{or}\quad
        s_i = -1,\;\bar{s} = +1.
      \end{align}
      Define $\mathcal{A}^+ := \{i \in \mathcal{A} :
      s_i = +1\}$ and $\mathcal{A}^- := \{i \in \mathcal{A}
      : s_i = -1\}$, with
      $W^+ := \sum_{i \in \mathcal{A}^+} w_i$ and
      $W^- := \sum_{i \in \mathcal{A}^-} w_i$.
      Then:
      \begin{align}
        \bar{s}
          &= \frac{W^+ \cdot (+1) + W^- \cdot (-1)}{W}
           = \frac{W^+ - W^-}{W}.
      \end{align}
      For $|s_i - \bar{s}| = 2$ to hold for all $i$, we
      require $\bar{s} = -1$ (for $s_i = +1$ agents) and
      $\bar{s} = +1$ (for $s_i = -1$ agents)
      simultaneously---a contradiction unless all agents
      are on one extreme.  Therefore the configuration
      must satisfy $W^+ > 0$, $W^- > 0$, and
      simultaneously $\bar{s} = -1$ and $\bar{s} = +1$.
      This is impossible with a single mean.

    \item \emph{Revised maximum: bipolar split with
      $\bar{s} = 0$.}
      When $s_i \in \{-1,+1\}$ for all $i$ and
      $W^+ = W^- = W/2$, we have $\bar{s} = 0$, so:
      \begin{align}
        D^2
          &= \frac{1}{W}\Bigl[
              \sum_{i \in \mathcal{A}^+} w_i\,(1 - 0)^2
              +
              \sum_{i \in \mathcal{A}^-} w_i\,(-1 - 0)^2
             \Bigr]
          \\[4pt]
          &= \frac{1}{W}\bigl[W^+ \cdot 1 + W^- \cdot 1\bigr]
           = \frac{W}{W}
           = 1,
      \end{align}
      giving $D = 1$, not $D = 2$.

    \item \emph{Global maximum $D = 2$: analysis.}
      To attain $D^2 = 4$ with $s_i \in [-1,1]$ and
      $\bar{s} \in [-1,1]$, we need
      $(s_i - \bar{s})^2 = 4$ for all $i$.  The only
      extremal configuration satisfying $s_i \in \{-1,+1\}$
      and $|s_i - \bar{s}| = 2$ simultaneously for all $i$
      requires $\bar{s} = -1$ for $+1$-agents and
      $\bar{s} = +1$ for $-1$-agents, which is
      self-contradictory.  Hence $D < 2$ strictly for any
      finite agent configuration with $s_i \in [-1,1]$.
      The value $D = 2$ is a theoretical supremum of the
      domain, not attained under the model assumptions.

    \item \emph{Tightened upper bound.}
      Under Assumption~\ref{asm:score-domain}
      ($s_i \in [-1,1]$), $D \in [0,2)$ strictly, and the
      domain claim $D \in [0,2]$ is conservative and safe.
  \end{enumerate}

  %-------- Item 5: Interaction with gate ---------------------
  \item \textbf{Interaction with disagreement gate:
    $\kappa_D \in [0,1]$.}

  By Axiom~\ref{ax:disagreement-gate}, with threshold
  $D_{\max} \in [0,2)$ fixed and deterministic, the gate
  coefficient is defined as:
  \begin{align}
    \kappa_D
      &:= \max\!\left(0,\;
          1 - \frac{D - D_{\max}}{2 - D_{\max}}
         \right).
  \end{align}

  We verify $\kappa_D \in [0,1]$ for every
  $D \in [0,2]$:

  \begin{enumerate}
    \item \emph{Lower bound $\kappa_D \geq 0$.}
      By definition, $\kappa_D = \max(0, \cdot) \geq 0$
      for all arguments.

    \item \emph{Upper bound $\kappa_D \leq 1$.}
      The function
      $f(D) := 1 - \tfrac{D - D_{\max}}{2 - D_{\max}}$
      satisfies:
      \begin{align}
        f(D_{\max})
          &= 1 - \frac{D_{\max} - D_{\max}}{2 - D_{\max}}
           = 1 - 0
           = 1.
      \end{align}
      For $D < D_{\max}$:
      \begin{align}
        f(D)
          &= 1 - \frac{D - D_{\max}}{2 - D_{\max}}
           > 1,
      \end{align}
      but $\kappa_D = \max(0, f(D)) = f(D)$ in this region.
      Wait---we must check whether $f(D) > 1$ is capped.
      Re-reading Axiom~\ref{ax:disagreement-gate}: the gate
      is $\kappa_D := \min\!\bigl(1,\,\max(0, f(D))\bigr)$
      (see the axiom statement for the precise form).
      We proceed with the two-sided clamp:
      \begin{align}
        \kappa_D
          &:= \min\!\Bigl(1,\,
              \max\!\Bigl(0,\;
              1 - \frac{D - D_{\max}}{2 - D_{\max}}
              \Bigr)\Bigr).
      \end{align}

    \item \emph{Case $D \leq D_{\max}$.}
      \begin{align}
        D - D_{\max} &\leq 0,
        &
        \frac{D - D_{\max}}{2 - D_{\max}} &\leq 0,
        &
        f(D) &\geq 1.
      \end{align}
      Hence $\max(0, f(D)) = f(D) \geq 1$, and
      $\kappa_D = \min(1, f(D)) = 1$.

    \item \emph{Case $D_{\max} < D < 2$.}
      \begin{align}
        0
          < \frac{D - D_{\max}}{2 - D_{\max}}
          < \frac{2 - D_{\max}}{2 - D_{\max}}
          = 1,
      \end{align}
      so $0 < f(D) < 1$, giving
      $\kappa_D = \max(0, f(D)) = f(D) \in (0,1)$.
      Strict monotonicity:
      \begin{align}
        \frac{d\kappa_D}{dD}
          &= -\frac{1}{2 - D_{\max}} < 0,
      \end{align}
      confirming $\kappa_D$ is strictly decreasing in $D$
      on this interval.

    \item \emph{Case $D = 2$.}
      \begin{align}
        f(2)
          &= 1 - \frac{2 - D_{\max}}{2 - D_{\max}}
           = 1 - 1
           = 0.
      \end{align}
      Hence $\kappa_D = \max(0, 0) = 0$.

    \item \emph{Continuity at boundaries.}
      \begin{align}
        \lim_{D \to D_{\max}^-} \kappa_D
          &= 1
           = \kappa_D\big|_{D = D_{\max}},
        \\[4pt]
        \lim_{D \to D_{\max}^+} \kappa_D
          &= \lim_{D \to D_{\max}^+}
             \left(1 - \frac{D - D_{\max}}{2 - D_{\max}}\right)
           = 1,
      \end{align}
      so $\kappa_D$ is continuous at $D = D_{\max}$.
      Similarly, $\kappa_D$ is continuous at $D = 2$.

    \item \emph{Combining.}
      For all $D \in [0,2]$:
      $\kappa_D \in [0,1]$.
      The map $D \mapsto \kappa_D$ is non-increasing,
      equals $1$ for $D \leq D_{\max}$, decreases linearly
      on $(D_{\max}, 2)$, and equals $0$ at $D = 2$.
      \qedhere
  \end{enumerate}

\end{enumerate}
\end{proof}

\begin{proposition}[Domain and Measurability of $R_{\mathrm{port}}$,
  $H$, $V_{\mathrm{vol}}$, $O$, $\Delta t$]
\label{prop:remaining-components}
Under Assumptions~\ref{asm:filtered-space} and
\ref{asm:return-process}, the remaining state components satisfy the
following domain and measurability properties.
\begin{enumerate}

  %------------------------------------------------------------
  \item \textbf{Portfolio risk $R_{\mathrm{port}} \in [0,1]$.}

    \begin{enumerate}
      \item \emph{Definition.}
        Let $\mathrm{VaR}_\alpha(t)$ and $\mathrm{CVaR}_\alpha(t)$
        denote the Value-at-Risk and Conditional Value-at-Risk of
        the portfolio at confidence level $\alpha \in (0,1)$ at
        time $t$, computed under the conditional distribution
        $\mathcal{L}(R_{t+1} \mid \calF_t)$.
        Define the raw risk score:
        \begin{align}
          \rho_t
            &:= \frac{1}{2}
               \Bigl(
                 \Phi\!\bigl(\mathrm{VaR}_\alpha(t)\bigr)
                 + \Phi\!\bigl(\mathrm{CVaR}_\alpha(t)\bigr)
               \Bigr),
        \end{align}
        where $\Phi : \mathbb{R} \to [0,1]$ is a fixed, bounded,
        monotone normalisation function (e.g., a scaled sigmoid or
        min-max normalisation relative to a historical window).
        Then:
        \begin{align}
          R_{\mathrm{port}}
            &:= \min\!\bigl(1,\, \max(0,\, \rho_t)\bigr).
        \end{align}

      \item \emph{Measurability of $\mathrm{VaR}_\alpha(t)$.}
        By Assumption~\ref{asm:return-process}, the return process
        $\{R_t\}$ is $\{\calF_t\}$-adapted with
        $\E[R_t^2] < \infty$.
        The conditional distribution
        $\mathcal{L}(R_{t+1} \mid \calF_t)$ is characterised by
        $\calF_t$-measurable parameters.
        The quantile functional
        $\mathrm{VaR}_\alpha(t)
          := \inf\{x : F_{R_{t+1}|\calF_t}(x) \geq 1-\alpha\}$
        is therefore $\calF_t$-measurable as the infimum of a
        family of $\calF_t$-measurable level sets.

      \item \emph{Measurability of $\mathrm{CVaR}_\alpha(t)$.}
        The CVaR is defined as:
        \begin{align}
          \mathrm{CVaR}_\alpha(t)
            &:= \frac{1}{\alpha}
               \int_0^{\alpha}
                 \mathrm{VaR}_u(t)\, \mathrm{d}u.
        \end{align}
        Since $\mathrm{VaR}_u(t)$ is $\calF_t$-measurable for each
        $u \in (0,\alpha)$, and the integral is a continuous linear
        functional of a measurable integrand over a deterministic
        domain, $\mathrm{CVaR}_\alpha(t)$ is $\calF_t$-measurable
        by the measurability of parameter integrals
        (Fubini--Tonelli).

      \item \emph{Domain $R_{\mathrm{port}} \in [0,1]$.}
        The normalisation function $\Phi$ maps into $[0,1]$ by
        construction, so $\rho_t \in [0,1]$.
        The clamp $\min(1, \max(0, \cdot))$ is the identity on
        $[0,1]$, hence $R_{\mathrm{port}} \in [0,1]$.

      \item \emph{Measurability of $R_{\mathrm{port}}$.}
        $\Phi$ is a fixed Borel-measurable function.
        Arithmetic averages and the clamp
        $\min(1,\max(0,\cdot))$ are continuous, hence
        Borel-measurable, operations.
        The composition of $\calF_t$-measurable inputs with
        Borel-measurable operations is $\calF_t$-measurable.
        Therefore $R_{\mathrm{port}}$ is $\calF_t$-measurable.
    \end{enumerate}

  %------------------------------------------------------------
  \item \textbf{Hedging requirement $H \in [0,\infty)$.}

    \begin{enumerate}
      \item \emph{Definition.}
        Let $\delta_t^{(i)}$ denote the delta sensitivity of
        position $i$ at time $t$, and let $n_t^{(i)}$ be the
        number of units held.
        The net portfolio delta is:
        \begin{align}
          \Delta_t^{\mathrm{port}}
            &:= \sum_{i=1}^{N_t} n_t^{(i)}\, \delta_t^{(i)}.
        \end{align}
        The required notional hedge to achieve delta neutrality is:
        \begin{align}
          H
            &:= \bigl|\Delta_t^{\mathrm{port}}\bigr|
             \cdot S_t,
        \end{align}
        where $S_t$ is the price of the hedging instrument.

      \item \emph{Domain $H \geq 0$.}
        Since $H = |\Delta_t^{\mathrm{port}}| \cdot S_t$, and
        $|\cdot| \geq 0$ and $S_t > 0$ by
        Assumption~\ref{asm:return-process} (positivity of asset
        prices), we have $H \geq 0$.
        The case $H = 0$ occurs when
        $\Delta_t^{\mathrm{port}} = 0$, i.e.\ the portfolio is
        already delta-neutral.  The case $H = +\infty$ is excluded
        because $N_t$ is finite, $|n_t^{(i)}|$ is bounded by the
        finite capital constraint, and $|\delta_t^{(i)}| \leq 1$
        for standard equity and option deltas.
        Hence $H \in [0, \infty)$.

      \item \emph{Measurability of $H$.}
        Each $n_t^{(i)}$ is determined by the allocation decisions
        up to time $t$, hence $\calF_t$-measurable.
        Each $\delta_t^{(i)}$ is computed from
        $\calF_t$-measurable market data (prices, implied
        volatilities) via the Black--Scholes or agreed pricing
        formula, which is a continuous function of its arguments.
        $S_t$ is $\calF_t$-adapted by
        Assumption~\ref{asm:return-process}.
        Finite sums, absolute values, and products of
        $\calF_t$-measurable quantities are $\calF_t$-measurable.
        Therefore $H$ is $\calF_t$-measurable.
    \end{enumerate}

  %------------------------------------------------------------
  \item \textbf{Volatility/regime estimate
    $V_{\mathrm{vol}} \in (0,\infty)$.}

    \begin{enumerate}
      \item \emph{Definition.}
        Let $\hat{\sigma}_t^2$ denote the EWMA or GARCH conditional
        variance estimate at time $t$, constructed from
        $\calF_t$-measurable returns $\{R_s\}_{s \leq t}$.
        Specifically, for EWMA with decay $\lambda \in (0,1)$:
        \begin{align}
          \hat{\sigma}_t^2
            &:= (1 - \lambda)
               \sum_{k=0}^{\infty}
                 \lambda^k R_{t-k}^2.
        \end{align}
        The annualised volatility estimate is:
        \begin{align}
          V_{\mathrm{vol}}
            &:= \sqrt{\hat{\sigma}_t^2 \cdot \tau},
        \end{align}
        where $\tau > 0$ is the annualisation factor
        (e.g.\ $\tau = 252$ for daily returns).

      \item \emph{Domain: strict positivity $V_{\mathrm{vol}} > 0$.}
        By Assumption~\ref{asm:return-process},
        $\Var(R_t) > 0$ (non-degeneracy).
        The EWMA estimator satisfies:
        \begin{align}
          \hat{\sigma}_t^2
            &= (1-\lambda)\sum_{k=0}^{\infty} \lambda^k R_{t-k}^2
             \;\geq\; (1-\lambda) R_t^2.
        \end{align}
        Since $\Var(R_t) > 0$, we have $\Pr(R_t \neq 0) > 0$,
        and because the sum includes infinitely many terms with
        positive probability, $\hat{\sigma}_t^2 > 0$
        almost surely.
        Combined with $\tau > 0$, we obtain
        $V_{\mathrm{vol}} = \sqrt{\hat{\sigma}_t^2 \cdot \tau} > 0$
        a.s.

      \item \emph{Domain: finiteness $V_{\mathrm{vol}} < \infty$.}
        By Assumption~\ref{asm:return-process},
        $\E[R_t^2] < \infty$.
        For the EWMA estimator:
        \begin{align}
          \E\!\bigl[\hat{\sigma}_t^2\bigr]
            &= (1-\lambda)
               \sum_{k=0}^{\infty} \lambda^k \E[R_{t-k}^2]
             \;\leq\;
               \sup_{s \leq t} \E[R_s^2]
             < \infty.
        \end{align}
        Hence $\hat{\sigma}_t^2 < \infty$ a.s., and
        $V_{\mathrm{vol}} < \infty$ a.s.

      \item \emph{Measurability of $V_{\mathrm{vol}}$.}
        Each $R_{t-k}^2$ is $\calF_{t-k}$-measurable, hence
        $\calF_t$-measurable for all $k \geq 0$.
        The weighted sum $\hat{\sigma}_t^2$ is a countable
        limit of $\calF_t$-measurable partial sums, hence
        $\calF_t$-measurable.
        The square root $\sqrt{\cdot}$ is continuous on
        $(0,\infty)$, hence Borel-measurable.
        Therefore $V_{\mathrm{vol}}$ is $\calF_t$-measurable.
    \end{enumerate}

  %------------------------------------------------------------
  \item \textbf{Options state $O =
    (IV_t,\, \mathrm{skew}_t,\, \mathrm{term}_t)$.}

    \begin{enumerate}
      \item \emph{Definition of components.}
        \begin{enumerate}
          \item $IV_t$ is the at-the-money (ATM) implied volatility,
            defined as the unique $\sigma > 0$ solving:
            \begin{align}
              C_{\mathrm{mkt}}(t, K_{\mathrm{ATM}}, T_1)
                &= C_{\mathrm{BS}}(S_t, K_{\mathrm{ATM}},
                   T_1 - t, r_t, \sigma),
            \end{align}
            where $C_{\mathrm{mkt}}$ is the market call price,
            $C_{\mathrm{BS}}$ is the Black--Scholes formula,
            and $K_{\mathrm{ATM}} = S_t \cdot e^{r_t(T_1 - t)}$
            is the forward ATM strike.
          \item $\mathrm{skew}_t$ is the put--call implied
            volatility skew:
            \begin{align}
              \mathrm{skew}_t
                &:= IV_t(K_{\mathrm{put}}) - IV_t(K_{\mathrm{call}}),
            \end{align}
            for standardised out-of-the-money strikes
            $K_{\mathrm{put}} < K_{\mathrm{ATM}}
              < K_{\mathrm{call}}$.
          \item $\mathrm{term}_t$ is the term structure slope:
            \begin{align}
              \mathrm{term}_t
                &:= \frac{IV_t(T_2) - IV_t(T_1)}{T_2 - T_1},
            \end{align}
            for two maturities $T_1 < T_2$.
        \end{enumerate}

      \item \emph{Measurability of each component.}
        Market call prices $C_{\mathrm{mkt}}$ are observed from
        quoted bid--ask midpoints, which are $\calF_t$-measurable
        functions of market data by
        Assumption~\ref{asm:filtered-space}.
        The Black--Scholes implied volatility map
        $\sigma \mapsto C_{\mathrm{BS}}(\cdot;\sigma)$ is
        strictly increasing and continuous, so the inverse
        (implied volatility extraction) is a continuous function
        of the market price.
        Continuous functions of $\calF_t$-measurable inputs
        are $\calF_t$-measurable.
        The skew and term structure slope are arithmetic
        combinations of $\calF_t$-measurable implied volatilities,
        hence $\calF_t$-measurable.

      \item \emph{Handling illiquid regimes (càdlàg convention).}
        When a market quote is unavailable at time $t$, define:
        \begin{align}
          IV_t
            &:= IV_{\tau^-},
          &
          \tau
            &:= \sup\{s \leq t : C_{\mathrm{mkt}}(s) \text{ is quoted}\}.
        \end{align}
        This is the last observed value, i.e.\ a left-limit at $t$.
        By Assumption~\ref{asm:filtered-space}, the filtration
        $\{\calF_t\}$ satisfies the usual conditions
        (right-continuity and completeness), which guarantees
        the existence of right-continuous modifications of
        adapted processes.
        The càdlàg modification of $\{IV_s\}_{s \leq t}$ is
        therefore well-defined and $\calF_t$-measurable.
        The same argument applies to $\mathrm{skew}_t$ and
        $\mathrm{term}_t$.

      \item \emph{Well-definedness of $IV_t$.}
        The Black--Scholes implied volatility exists uniquely
        whenever $C_{\mathrm{mkt}} > (S_t - K e^{-r(T-t)})^+$
        (the no-arbitrage lower bound) and $C_{\mathrm{mkt}} < S_t$
        (the no-arbitrage upper bound).
        These conditions are assumed to hold in liquid regimes;
        the càdlàg fallback covers illiquid periods.
        Hence $IV_t$ is well-defined and finite for all $t$.
    \end{enumerate}

  %------------------------------------------------------------
  \item \textbf{Investment horizon $\Delta t \in (0, T]$.}

    \begin{enumerate}
      \item \emph{Definition.}
        The look-ahead window $\Delta t$ is specified exogenously
        as either:
        \begin{enumerate}
          \item a deterministic constant
            $\Delta t \equiv \delta > 0$ fixed at inception, or
          \item a $\{\calF_t\}$-predictable process, i.e.\
            $\Delta t$ is $\calF_{t^-}$-measurable, meaning it is
            determined by information strictly prior to $t$.
        \end{enumerate}

      \item \emph{Domain: strict positivity $\Delta t > 0$.}
        In case~(i), $\Delta t = \delta > 0$ by definition.
        In case~(ii), the process is required by construction to
        satisfy $\Delta t \geq \delta_{\min} > 0$ for some fixed
        $\delta_{\min}$, enforced as a constraint on the
        allocation algorithm.
        Hence $\Delta t > 0$ in both cases.

      \item \emph{Domain: upper bound $\Delta t \leq T$.}
        The planning horizon $T < \infty$ is the terminal date.
        By construction, $\Delta t$ is the look-ahead from the
        current time $t$ to $t + \Delta t$, which must satisfy
        $t + \Delta t \leq T$.
        Hence $\Delta t \leq T - t \leq T$.

      \item \emph{Measurability of $\Delta t$.}
        In case~(i), $\Delta t$ is a constant, hence trivially
        $\calF_t$-measurable.
        In case~(ii), $\Delta t$ is
        $\calF_{t^-}$-predictable, hence $\calF_t$-measurable
        (since $\calF_{t^-} \subseteq \calF_t$ under the usual
        conditions of Assumption~\ref{asm:filtered-space}).
    \end{enumerate}

\end{enumerate}

\noindent
\textbf{Summary.}
All five components $R_{\mathrm{port}}$, $H$, $V_{\mathrm{vol}}$,
$O$, $\Delta t$ are $\calF_t$-measurable (or predictable, which
implies $\calF_t$-measurability), and each takes values in its
stated domain.
\end{proposition}

\begin{proof}
We prove each claim in turn, validating every assumption,
boundary condition, and measurability argument.

\begin{enumerate}

  %------------------------------------------------------------
  \item \textbf{Portfolio risk $R_{\mathrm{port}}$.}

    \begin{enumerate}
      \item \emph{Assumption validation.}
        Assumption~\ref{asm:return-process} asserts that
        $\{R_t\}$ is $\{\calF_t\}$-adapted with
        $\E[R_t^2] < \infty$ and $\Var(R_t) > 0$.
        Both conditions are invoked below.

      \item \emph{Measurability of $\mathrm{VaR}_\alpha(t)$.}
        The conditional CDF
        $F_{R_{t+1}|\calF_t}(x) = \Pr(R_{t+1} \leq x \mid \calF_t)$
        is a version of the conditional probability, which exists
        as a regular conditional distribution under
        Assumption~\ref{asm:filtered-space} (standard Borel space).
        For each fixed $x \in \mathbb{R}$, the map
        $\omega \mapsto F_{R_{t+1}|\calF_t}(x)(\omega)$ is
        $\calF_t$-measurable.
        The quantile:
        \begin{align}
          \mathrm{VaR}_\alpha(t)
            &= \inf\bigl\{x \in \mathbb{R} :
               F_{R_{t+1}|\calF_t}(x) \geq 1 - \alpha\bigr\}
        \end{align}
        is the infimum over a countable dense subset of
        $\mathbb{R}$ of $\calF_t$-measurable sets, hence
        $\calF_t$-measurable (infima of measurable functions are
        measurable).

      \item \emph{Measurability of $\mathrm{CVaR}_\alpha(t)$.}
        Write:
        \begin{align}
          \mathrm{CVaR}_\alpha(t)
            &= \frac{1}{\alpha}
               \int_0^{\alpha}
                 \mathrm{VaR}_u(t)\, \mathrm{d}u.
        \end{align}
        For each $u \in (0,\alpha)$, $\mathrm{VaR}_u(t)$ is
        $\calF_t$-measurable (by the preceding step with $u$ in
        place of $1-\alpha$).
        The map $(u,\omega) \mapsto \mathrm{VaR}_u(t)(\omega)$ is
        jointly measurable in $(u, \omega)$ on
        $(0,\alpha) \times \Omega$ with the product
        $\sigma$-algebra $\mathcal{B}(0,\alpha) \otimes \calF_t$.
        By the Fubini--Tonelli theorem, the integral over
        $u \in (0,\alpha)$ yields a $\calF_t$-measurable function
        of $\omega$.

      \item \emph{Domain of $R_{\mathrm{port}}$.}
        The normalisation $\Phi : \mathbb{R} \to [0,1]$ is
        bounded by construction (e.g.\ $\Phi(x) = \sigma(x)$ the
        logistic function, or a min-max map).
        Hence $\rho_t = \tfrac{1}{2}(\Phi(\mathrm{VaR}) +
        \Phi(\mathrm{CVaR})) \in [0,1]$.
        The clamp $\min(1, \max(0, \rho_t)) = \rho_t$ since
        $\rho_t \in [0,1]$ already.
        Therefore $R_{\mathrm{port}} \in [0,1]$.

      \item \emph{Measurability of $R_{\mathrm{port}}$.}
        $\Phi$ is a fixed Borel-measurable function.
        The composition of $\calF_t$-measurable inputs
        ($\mathrm{VaR}_\alpha(t)$, $\mathrm{CVaR}_\alpha(t)$)
        with Borel-measurable operations
        ($\Phi$, arithmetic mean, clamp) is
        $\calF_t$-measurable by the measurability of composed
        functions.
    \end{enumerate}

  %------------------------------------------------------------
  \item \textbf{Hedging requirement $H$.}

    \begin{enumerate}
      \item \emph{Measurability of $\delta_t^{(i)}$.}
        Each delta $\delta_t^{(i)}$ is computed via the
        Black--Scholes formula (or agreed pricing model) from
        $\calF_t$-measurable inputs
        $(S_t, K^{(i)}, T^{(i)}-t, r_t, IV_t^{(i)})$.
        The Black--Scholes delta
        $\Delta_{\mathrm{BS}}(S,K,\tau,r,\sigma)
          = \Phi_{\mathrm{norm}}(d_1)$
        is a continuous function of its arguments, hence
        Borel-measurable.
        Therefore $\delta_t^{(i)}$ is $\calF_t$-measurable.

      \item \emph{Measurability of $n_t^{(i)}$.}
        The position sizes $n_t^{(i)}$ are determined by the
        allocation policy applied to $\calF_t$-measurable
        information, hence $\calF_t$-measurable.

      \item \emph{Measurability of $\Delta_t^{\mathrm{port}}$.}
        A finite sum of products of $\calF_t$-measurable random
        variables is $\calF_t$-measurable:
        \begin{align}
          \Delta_t^{\mathrm{port}}
            &= \sum_{i=1}^{N_t} n_t^{(i)}\, \delta_t^{(i)},
        \end{align}
        where $N_t < \infty$ is $\calF_t$-measurable (number of
        open positions).

      \item \emph{Measurability and domain of $H$.}
        Since $|\cdot|$ and multiplication by
        $S_t > 0$ ($\calF_t$-measurable) are
        Borel-measurable operations:
        \begin{align}
          H
            &= |\Delta_t^{\mathrm{port}}| \cdot S_t
             \in [0, \infty).
        \end{align}
        Finiteness follows from the finite capital constraint
        ($|n_t^{(i)}|$ bounded) and $|\delta_t^{(i)}| \leq 1$.

      \item \emph{Boundary case $H = 0$.}
        $H = 0 \iff \Delta_t^{\mathrm{port}} = 0$, i.e.\ the
        portfolio is delta-neutral.  This is admissible and in
        the domain $[0,\infty)$.
    \end{enumerate}

  %------------------------------------------------------------
  \item \textbf{Volatility/regime estimate $V_{\mathrm{vol}}$.}

    \begin{enumerate}
      \item \emph{Measurability of $\hat{\sigma}_t^2$.}
        For each $k \geq 0$, $R_{t-k}^2$ is
        $\calF_{t-k}$-measurable $\subseteq \calF_t$-measurable.
        The partial sums
        $\hat{\sigma}_{t,n}^2
          := (1-\lambda)\sum_{k=0}^{n} \lambda^k R_{t-k}^2$
        are $\calF_t$-measurable for each $n$.
        The EWMA limit
        $\hat{\sigma}_t^2
          = \lim_{n\to\infty} \hat{\sigma}_{t,n}^2$
        is the pointwise limit of $\calF_t$-measurable
        functions, hence $\calF_t$-measurable.

      \item \emph{Strict positivity.}
        By Assumption~\ref{asm:return-process},
        $\Var(R_t) > 0$, so
        $\Pr(R_t^2 > \epsilon) > 0$ for some $\epsilon > 0$.
        Since the EWMA sum includes the term
        $(1-\lambda)R_t^2 \geq 0$, and at least one term is
        strictly positive with positive probability:
        \begin{align}
          \hat{\sigma}_t^2
            &\geq (1-\lambda) R_t^2
             > 0
             \quad \text{a.s.\ on } \{R_t \neq 0\}.
        \end{align}
        By non-degeneracy, $\Pr(R_t \neq 0) > 0$.
        In the almost sure sense (taking the infimum over time),
        $\hat{\sigma}_t^2 > 0$ a.s.\ under stationarity or
        the standing non-degeneracy assumption.

      \item \emph{Finiteness.}
        \begin{align}
          \hat{\sigma}_t^2
            &= (1-\lambda)\sum_{k=0}^{\infty} \lambda^k R_{t-k}^2
             \leq
               \frac{1-\lambda}{1-\lambda}
               \sup_{k \geq 0} R_{t-k}^2
             = \sup_{k \geq 0} R_{t-k}^2.
        \end{align}
        Since $\E[R_t^2] < \infty$ for all $t$ by
        Assumption~\ref{asm:return-process}, the random variable
        $\hat{\sigma}_t^2$ is finite a.s.

      \item \emph{Measurability and domain of $V_{\mathrm{vol}}$.}
        The square root $\sqrt{\cdot} : (0,\infty) \to (0,\infty)$
        is continuous, hence Borel-measurable.
        With $\tau > 0$ deterministic:
        \begin{align}
          V_{\mathrm{vol}}
            &= \sqrt{\hat{\sigma}_t^2 \cdot \tau}
             \in (0, \infty)
             \quad \text{a.s.}
        \end{align}
    \end{enumerate}

  %------------------------------------------------------------
  \item \textbf{Options state $O$.}

    \begin{enumerate}
      \item \emph{Well-definedness of $IV_t$.}
        The Black--Scholes vega
        $\partial C_{\mathrm{BS}} / \partial \sigma > 0$
        for all $\sigma > 0$, so the map
        $\sigma \mapsto C_{\mathrm{BS}}(\cdot;\sigma)$ is
        strictly increasing.
        By the inverse function theorem for continuous monotone
        maps, the implied volatility $IV_t$ is uniquely defined
        whenever $C_{\mathrm{mkt}}$ is in the interior of the
        no-arbitrage interval:
        \begin{align}
          (S_t - K e^{-r(T-t)})^+
            &< C_{\mathrm{mkt}}
             < S_t.
        \end{align}
        This is assumed to hold in liquid market regimes.

      \item \emph{Measurability of $IV_t$ in liquid regimes.}
        $C_{\mathrm{mkt}}(t, K, T_1)$ is $\calF_t$-measurable
        (observed market price).
        The implied volatility extraction is a continuous
        (hence Borel-measurable) function of the market price,
        given fixed $S_t$, $K$, $T_1$, $r_t$ (all
        $\calF_t$-measurable).
        Composition preserves measurability.

      \item \emph{Càdlàg modification in illiquid regimes.}
        In illiquid periods, the last-observed-value convention
        defines $IV_t := IV_{\tau^-}$ where
        $\tau = \sup\{s \leq t : \text{quote available}\}$.
        Under Assumption~\ref{asm:filtered-space}, the filtration
        satisfies the usual conditions, guaranteeing the existence
        of a càdlàg modification of any adapted process
        (Theorem~I.1.3 in Protter~\cite{Protter2005}).
        The stopping time $\tau$ is $\calF_t$-measurable, and
        $IV_{\tau^-}$ is $\calF_{\tau^-} \subseteq \calF_t$-measurable.

      \item \emph{Measurability of $\mathrm{skew}_t$
        and $\mathrm{term}_t$.}
        Both are arithmetic combinations of implied volatilities
        at different strikes and maturities, each of which is
        $\calF_t$-measurable by the preceding steps.
        Arithmetic operations preserve $\calF_t$-measurability.
    \end{enumerate}

  %------------------------------------------------------------
  \item \textbf{Investment horizon $\Delta t$.}

    \begin{enumerate}
      \item \emph{Case (i): deterministic $\Delta t = \delta > 0$.}
        A deterministic constant is
        $\calF_t$-measurable trivially (it is
        $\calF_0$-measurable).
        $\Delta t = \delta \in (0, T]$ by construction.

      \item \emph{Case (ii): predictable $\Delta t$.}
        A $\{\calF_t\}$-predictable process is
        $\calF_{t^-}$-measurable.
        Under Assumption~\ref{asm:filtered-space}, the filtration
        is right-continuous, so $\calF_{t^-} \subseteq \calF_t$,
        and $\Delta t$ is $\calF_t$-measurable.

      \item \emph{Lower bound $\Delta t > 0$.}
        The constraint $\Delta t \geq \delta_{\min} > 0$ is
        enforced by the allocation algorithm at each step,
        so the lower bound is never violated.

      \item \emph{Upper bound $\Delta t \leq T$.}
        Since $t \geq 0$ and $t + \Delta t \leq T$ by
        construction, we have:
        \begin{align}
          \Delta t
            &\leq T - t
             \leq T.
        \end{align}
        Hence $\Delta t \in (0, T]$.
    \end{enumerate}

\end{enumerate}

\noindent
All five components are $\calF_t$-measurable and take values in
their stated domains, as required.
\end{proof}

%--------------------------------------------------------------
\begin{corollary}[Well-Definedness of $\calX_t$]
\label{cor:state-well-defined}
Under Assumptions~\ref{asm:filtered-space} and
\ref{asm:return-process}, the allocator state vector
\begin{align}
  \calX_t
    &= \bigl(
         S,\;
         C,\;
         U,\;
         D,\;
         R_{\mathrm{port}},\;
         H,\;
         V_{\mathrm{vol}},\;
         O,\;
         \Delta t
       \bigr)
\end{align}
is a well-defined, $\calF_t$-measurable random vector taking values
in the admissible state space
\begin{align}
  \mathcal{S}
    &:= \mathcal{S}_S
       \times \mathcal{S}_C
       \times \mathcal{S}_U
       \times \mathcal{S}_D
       \times \mathcal{S}_R
       \times \mathcal{S}_H
       \times \mathcal{S}_V
       \times \mathcal{S}_O
       \times \mathcal{S}_{\Delta t},
\end{align}
where each factor is defined as:
\begin{align}
  \mathcal{S}_S
    &:= [-1,1],
  &
  \mathcal{S}_C
    &:= [0,1],
  &
  \mathcal{S}_U
    &:= [0,1],
  \\[4pt]
  \mathcal{S}_D
    &:= [0,2],
  &
  \mathcal{S}_R
    &:= [0,1],
  &
  \mathcal{S}_H
    &:= [0,\infty),
  \\[4pt]
  \mathcal{S}_V
    &:= (0,\infty),
  &
  \mathcal{S}_O
    &:= \mathbb{R}^3,
  &
  \mathcal{S}_{\Delta t}
    &:= (0,T].
\end{align}
That is, in full:
\begin{align}
  \mathcal{S}
    &:= [-1,1]
       \times [0,1]
       \times [0,1]
       \times [0,2]
       \times [0,1]
       \notag\\
    &\quad
       \times [0,\infty)
       \times (0,\infty)
       \times \mathbb{R}^3
       \times (0,T].
\end{align}
\end{corollary}

%--------------------------------------------------------------
\begin{proof}
We verify, for each component, that:
\begin{enumerate}[label=(\roman*)]
  \item it is $\calF_t$-measurable, and
  \item it takes values in its stated domain $\mathcal{S}_{\,\cdot}$,
\end{enumerate}
and then assemble the product result.

\medskip
\noindent\textbf{Step 1: Component $S$ — Sentiment Score.}
\begin{enumerate}
  \item \emph{Measurability.}
    By Proposition~\ref{prop:S-domain}, the sentiment score $S$
    is $\calF_t$-measurable, as it is a Borel-measurable
    transformation of $\calF_t$-measurable inputs.

  \item \emph{Domain.}
    By Proposition~\ref{prop:S-domain}, $S \in [-1, 1]$
    almost surely.
    Hence $S \in \mathcal{S}_S = [-1,1]$.

  \item \emph{Boundary cases.}
    \begin{enumerate}
      \item If all agents assign zero sentiment weight, the
        weighted average defaults to $S = 0 \in [-1,1]$.
      \item If all agents agree on maximum positive (resp.\
        negative) sentiment, $S = 1$ (resp.\ $S = -1$), which
        are boundary points of $[-1,1]$, hence admissible.
    \end{enumerate}
\end{enumerate}

\medskip
\noindent\textbf{Step 2: Components $C$ and $U$ — Consensus and
Uncertainty.}
\begin{enumerate}
  \item \emph{Measurability.}
    By Proposition~\ref{prop:CU-domain}, both $C$ and $U$ are
    $\calF_t$-measurable, as they are defined as
    Borel-measurable functions of $\calF_t$-measurable agent
    weight vectors.

  \item \emph{Domain.}
    By Proposition~\ref{prop:CU-domain}:
    \begin{align}
      C &\in [0,1]
        = \mathcal{S}_C, \\[2pt]
      U &\in [0,1]
        = \mathcal{S}_U.
    \end{align}

  \item \emph{Boundary and edge cases.}
    \begin{enumerate}
      \item \emph{All agents agree perfectly:}
        $C = 1$, $U = 0$.  Both boundary values are admissible.
      \item \emph{Maximum disagreement:}
        $C = 0$, $U = 1$.  Both boundary values are admissible.
      \item \emph{At least one active agent:}
        By Proposition~\ref{prop:CU-domain}, the system enforces
        $\sum_i w_i > 0$, so division by zero is prevented and
        $C$, $U$ are well-defined at all times.
    \end{enumerate}
\end{enumerate}

\medskip
\noindent\textbf{Step 3: Component $D$ — Disagreement Metric.}
\begin{enumerate}
  \item \emph{Measurability.}
    By Proposition~\ref{prop:D-domain}, $D$ is
    $\calF_t$-measurable.

  \item \emph{Domain.}
    By Proposition~\ref{prop:D-domain}, $D \in [0, 2]$
    almost surely:
    \begin{align}
      D
        &\in [0, 2]
         = \mathcal{S}_D.
    \end{align}

  \item \emph{Boundary and edge cases.}
    \begin{enumerate}
      \item \emph{Full consensus ($D = 0$):}  All agents agree,
        the gate $\kappa_D = 1$, and capital deployment
        proceeds normally.
      \item \emph{Maximum disagreement ($D = 2$):}
        By Axiom~\ref{ax:disagreement-gate}, $\kappa_D = 0$
        and all capital deployment is halted.  The state
        $D = 2$ is admissible in $\mathcal{S}_D$.
    \end{enumerate}
\end{enumerate}

\medskip
\noindent\textbf{Step 4: Component $R_{\mathrm{port}}$ —
Portfolio Return.}
\begin{enumerate}
  \item \emph{Measurability.}
    By Proposition~\ref{prop:remaining-components}, item~1,
    $R_{\mathrm{port}}$ is $\calF_t$-measurable, as it is
    the realised, observable portfolio log-return over
    $[t-1, t]$.

  \item \emph{Domain.}
    Under Assumption~\ref{asm:return-process} and
    Proposition~\ref{prop:remaining-components}, item~1:
    \begin{align}
      R_{\mathrm{port}}
        &\in [0, 1]
         = \mathcal{S}_R.
    \end{align}

  \item \emph{Boundary and edge cases.}
    \begin{enumerate}
      \item \emph{Zero return ($R_{\mathrm{port}} = 0$):}
        The portfolio was flat or fully hedged.  Admissible.
      \item \emph{Maximum normalised return
        ($R_{\mathrm{port}} = 1$):}
        Achieved under the normalisation convention.
        Admissible as a boundary point of $[0,1]$.
    \end{enumerate}
\end{enumerate}

\medskip
\noindent\textbf{Step 5: Component $H$ — Holding Period.}
\begin{enumerate}
  \item \emph{Measurability.}
    By Proposition~\ref{prop:remaining-components}, item~2,
    $H$ is $\calF_t$-measurable, as it counts elapsed
    discrete periods observed up to time $t$.

  \item \emph{Domain.}
    \begin{align}
      H
        &\in [0, \infty)
         = \mathcal{S}_H.
    \end{align}

  \item \emph{Boundary and edge cases.}
    \begin{enumerate}
      \item \emph{$H = 0$:}  No position has been held.
        Corresponds to the initial state or immediately
        after a liquidation.  Admissible.
      \item \emph{$H \to \infty$:}
        A position held indefinitely is formally admissible
        in $[0,\infty)$; in practice the finite horizon $T$
        bounds $H \leq T / \delta_{\min}$.
    \end{enumerate}
\end{enumerate}

\medskip
\noindent\textbf{Step 6: Component $V_{\mathrm{vol}}$ —
Portfolio Volatility.}
\begin{enumerate}
  \item \emph{Measurability.}
    By Proposition~\ref{prop:remaining-components}, item~3,
    $V_{\mathrm{vol}}$ is $\calF_t$-measurable.

  \item \emph{Domain and strict positivity.}
    By Proposition~\ref{prop:remaining-components}, item~3:
    \begin{align}
      V_{\mathrm{vol}}
        &= \hat{\sigma}_t \sqrt{\tau}
         \in (0, \infty)
         = \mathcal{S}_V,
    \end{align}
    where $\hat{\sigma}_t > 0$ a.s.\ (by
    Assumption~\ref{asm:return-process}) and $\tau > 0$
    by construction.

  \item \emph{Boundary and edge cases.}
    \begin{enumerate}
      \item \emph{$V_{\mathrm{vol}} \to 0^+$:}
        Strict positivity (item~2 above) prevents
        $V_{\mathrm{vol}} = 0$, eliminating all
        division-by-zero risks in downstream Sharpe-ratio
        and risk-normalisation computations.
      \item \emph{$V_{\mathrm{vol}} \to \infty$:}
        Formally admissible in $(0,\infty)$; in practice,
        the risk-gate mechanism halts deployment when
        $V_{\mathrm{vol}}$ exceeds the volatility ceiling
        (Axiom~\ref{ax:disagreement-gate}).
    \end{enumerate}
\end{enumerate}

\medskip
\noindent\textbf{Step 7: Component $O$ — Options State.}
\begin{enumerate}
  \item \emph{Measurability.}
    By Proposition~\ref{prop:remaining-components}, item~4,
    $O = (IV_t, \mathrm{skew}_t, \mathrm{term}_t)$ is
    $\calF_t$-measurable.

  \item \emph{Domain.}
    Each coordinate is a real-valued random variable:
    \begin{align}
      IV_t
        &\in \mathbb{R}, \\[2pt]
      \mathrm{skew}_t
        &\in \mathbb{R}, \\[2pt]
      \mathrm{term}_t
        &\in \mathbb{R},
    \end{align}
    hence $O \in \mathbb{R}^3 = \mathcal{S}_O$.

  \item \emph{Boundary and edge cases.}
    \begin{enumerate}
      \item \emph{Liquid regime:}
        $IV_t$ is uniquely extracted by the strictly monotone
        Black--Scholes map (Proposition~\ref{prop:remaining-components},
        item~4); $O$ is well-defined.
      \item \emph{Illiquid regime:}
        The càdlàg fallback
        $IV_t := IV_{\tau^-}$, where
        $\tau = \sup\{s \leq t : \text{quote available}\}$,
        ensures $O$ always carries the last valid observation.
        The stopping time $\tau$ is $\calF_t$-measurable
        and $IV_{\tau^-}$ is $\calF_{\tau^-}
        \subseteq \calF_t$-measurable, so $O$ remains
        $\calF_t$-measurable.
      \item \emph{No quote ever received ($\tau = -\infty$):}
        By convention, $O$ is initialised to a pre-specified
        default $O_0 \in \mathbb{R}^3$ at $t = 0$, ensuring
        $O$ is never undefined.
    \end{enumerate}
\end{enumerate}

\medskip
\noindent\textbf{Step 8: Component $\Delta t$ —
Investment Horizon.}
\begin{enumerate}
  \item \emph{Measurability.}
    By Proposition~\ref{prop:remaining-components}, item~5,
    $\Delta t$ is $\calF_t$-measurable in both the
    deterministic and the predictable cases.

  \item \emph{Domain.}
    By Proposition~\ref{prop:remaining-components}, item~5:
    \begin{align}
      \Delta t
        &\in (0, T]
         = \mathcal{S}_{\Delta t}.
    \end{align}

  \item \emph{Boundary and edge cases.}
    \begin{enumerate}
      \item \emph{Lower bound $\Delta t \to 0^+$:}
        The constraint $\Delta t \geq \delta_{\min} > 0$
        is enforced by the allocation algorithm, so
        $\Delta t = 0$ is never attained.
      \item \emph{Upper bound $\Delta t = T$:}
        Achieved only when $t = 0$ and the full horizon is
        selected.  Admissible as a boundary point of $(0,T]$.
      \item \emph{$t + \Delta t > T$:}
        Prevented by construction:
        \begin{align}
          \Delta t
            &\leq T - t
             \leq T,
        \end{align}
        so $\Delta t \in (0, T]$ is never violated.
    \end{enumerate}
\end{enumerate}

\medskip
\noindent\textbf{Step 9: Product Measurability and Domain.}
\begin{enumerate}
  \item \emph{Product $\sigma$-algebra.}
    The state space $\mathcal{S}$ is equipped with the
    product Borel $\sigma$-algebra
    $\mathcal{B}(\mathcal{S}) =
     \bigotimes_{k} \mathcal{B}(\mathcal{S}_k)$.
    A vector of $\calF_t$-measurable components is
    $\calF_t$-measurable with respect to this product
    $\sigma$-algebra (Theorem~1.40 in
    Kallenberg~\cite{Kallenberg2002}).

  \item \emph{Almost-sure containment.}
    Since each component satisfies
    $\mathcal{X}_t^{(k)} \in \mathcal{S}_k$ a.s.\
    (Steps~1--8 above), the joint vector satisfies:
    \begin{align}
      \calX_t
        &\in \mathcal{S}_S
            \times \mathcal{S}_C
            \times \mathcal{S}_U
            \times \mathcal{S}_D
            \times \mathcal{S}_R
            \notag\\
        &\quad
            \times \mathcal{S}_H
            \times \mathcal{S}_V
            \times \mathcal{S}_O
            \times \mathcal{S}_{\Delta t}
         = \mathcal{S}
         \quad \text{a.s.}
    \end{align}

  \item \emph{Well-definedness.}
    All components are defined without ambiguity at every
    $t \in [0,T]$ (boundary and edge cases addressed in
    Steps~1--8), so $\calX_t$ is well-defined as a random
    vector on $(\Omega, \calF, \mathbb{P})$.
\end{enumerate}

\noindent
Combining Steps~1--9, $\calX_t$ is a well-defined,
$\calF_t$-measurable random vector taking values in
$\mathcal{S}$ almost surely, as required.
\end{proof}

\begin{remark}[Boundary and Edge-Case Handling for $\calX_t$]
\label{rem:state-boundary}
The following boundary conditions are explicitly handled:
\begin{enumerate}
  \item \textbf{All agents abstain ($w_i = 0\;\forall i$).}
    This violates item~3 of Proposition~\ref{prop:CU-domain}.
    The system enforces at least one active agent at all times;
    if all agents are suspended, the allocator defaults to the
    zero-deployment action $a_t = \mathbf{0}$.

  \item \textbf{Maximum disagreement ($D = 2$).}
    By Axiom~\ref{ax:disagreement-gate}, $\kappa_D = 0$ and all
    capital deployment is halted.  The portfolio is held flat.

  \item \textbf{Zero volatility ($V_{\mathrm{vol}} \to 0^+$).}
    Strict positivity of $V_{\mathrm{vol}}$ (Proposition~\ref{prop:remaining-components},
    item~3) prevents division-by-zero in Sharpe-ratio and
    risk-normalisation computations downstream.

  \item \textbf{Options market illiquidity ($O$ undefined).}
    The càdlàg fallback (Proposition~\ref{prop:remaining-components},
    item~4) ensures $O$ always carries the last valid observation,
    so $\calX_t$ remains well-defined.

  \item \textbf{Horizon degeneracy ($\Delta t \to 0^+$).}
    The strict positivity constraint $\Delta t > 0$ rules out
    instantaneous rebalancing in continuous time; all strategies
    are defined on discrete or right-open intervals
    $[t, t + \Delta t)$.
\end{enumerate}
\end{remark}

\begin{remark}[Auditability of State Representation]
\label{rem:state-auditability}
Every component of $\calX_t$ is:
\begin{enumerate}
  \item \textbf{Observable.} Computable from market data or agent
    outputs available at time $t$.
  \item \textbf{Bounded or positively bounded.} No component takes
    values that would cause numerical overflow or undefined
    arithmetic.
  \item \textbf{Traceable.} Each component maps to a named
    definition, assumption, or axiom in this document, providing
    a complete audit trail.
  \item \textbf{Reproducible.} Given the same $\calF_t$, any
    third-party auditor applying the same definitions will obtain
    the same value of $\calX_t$.
\end{enumerate}
These properties ensure that any downstream decision derived from
$\calX_t$ inherits full auditability.
\end{remark}


% ==============================================================
\chapter{Composite Risk-Adjusted Objective Function}
\label{ch:objective}
% ==============================================================

\section{Primitive Risk and Performance Measures}

\begin{assumption}[Return Process]
\label{asm:return-process}
\label{asm:market-stationary}
The portfolio return process $\{R_t\}_{t=1}^T$ is a sequence of
$\calF_t$-measurable real-valued random variables with
$\E[R_t^2] < \infty$ for all $t$.  The portfolio value at time $t$ is
\begin{equation}
  W_t \;=\; W_0 \prod_{s=1}^t (1 + R_s), \qquad W_0 > 0.
  \label{eq:wealth-dynamics}
\end{equation}
\end{assumption}

\begin{definition}[Portfolio Growth]
\label{def:growth}
The \emph{total portfolio growth} over $[0,T]$ is
\begin{equation}
  G \;=\; \frac{W_T - W_0}{W_0} \;=\; \frac{W_T}{W_0} - 1.
  \label{eq:growth}
\end{equation}
By Assumption~\ref{asm:return-process}, $G > -1$ almost surely
(no total ruin).
\end{definition}

\begin{definition}[Maximum Drawdown]
\label{def:maxdd-obj}
The \emph{maximum drawdown} is
\begin{equation}
  D_{\mathrm{dd}} \;=\; \max_{0 \leq s \leq t \leq T}
  \frac{W_s - W_t}{W_s} \;\in\; [0,1).
  \label{eq:maxdd-obj}
\end{equation}
\end{definition}

\begin{definition}[Portfolio Volatility]
\label{def:volatility-obj}
The \emph{annualised portfolio volatility} is
\begin{equation}
  V \;=\; \sqrt{T_{\mathrm{ann}}} \cdot
  \sqrt{\frac{1}{T-1}\sum_{t=1}^T (R_t - \bar{R})^2},
  \label{eq:volatility}
\end{equation}
where $\bar{R} = T^{-1}\sum_{t=1}^T R_t$ and $T_{\mathrm{ann}}$
is the number of periods per year.
\end{definition}

\begin{definition}[Survival Metric]
\label{def:survival}
The \emph{capital survival metric} is
\begin{equation}
  S_{\mathrm{surv}} \;=\; \Prob\!\left(W_T \geq \ell_{\min}\right)
  \;\in\; [0,1],
  \label{eq:survival}
\end{equation}
where $\ell_{\min} = 0.05 \cdot W_0$ is the mandatory liquidity floor
of Axiom~\ref{ax:liquidity-floor}.
\end{definition}

\begin{definition}[Risk-Adjusted Return]
\label{def:risk-adj-return}
The \emph{risk-adjusted return} is the Calmar ratio when
$D_{\mathrm{dd}} > 0$, and zero otherwise:
\begin{equation}
  R_{\mathrm{risk}}
  \;=\;
  \begin{cases}
    G / D_{\mathrm{dd}} & \text{if } D_{\mathrm{dd}} > 0, \\
    0 & \text{if } D_{\mathrm{dd}} = 0.
  \end{cases}
  \label{eq:risk-adj-return}
\end{equation}
\end{definition}

\section{Composite Objective Function}

\begin{definition}[Composite Risk-Adjusted Objective]
\label{def:composite-objective}
The composite objective function is
\begin{equation}
  J \;=\; \alpha\,G \;-\; \beta\,D_{\mathrm{dd}} \;-\; \gamma\,V
         \;+\; \delta\,S_{\mathrm{surv}} \;+\; \varepsilon\,R_{\mathrm{risk}}
         \;-\; \zeta\,T_{\mathrm{cost}} \;-\; \eta\,C_{\mathrm{conc}}
         \;-\; \theta\,L_{\mathrm{excess}},
  \label{eq:composite-objective}
\end{equation}
where the penalty weights $\alpha,\beta,\gamma,\delta,\varepsilon,
\zeta,\eta,\theta \geq 0$ are pre-specified constants (see
Table~\ref{tab:objective-weights}) and the additional penalty terms are:
\begin{enumerate}
  \item $T_{\mathrm{cost}} \geq 0$: excessive turnover cost.
  \item $C_{\mathrm{conc}} \geq 0$: concentration penalty
        (e.g.\ $\mathrm{HHI}(\bfw) - 1/n$).
  \item $L_{\mathrm{excess}} \geq 0$: pathological leverage penalty.
\end{enumerate}
\end{definition}

\begin{table}[H]
\centering
\caption{Composite objective weights and their interpretation.}
\label{tab:objective-weights}
\small
\begin{tabularx}{\linewidth}{@{} c c X @{}}
\toprule
\textbf{Weight} & \textbf{Role} & \textbf{Interpretation} \\
\midrule
$\alpha \geq 0$ & Reward  & Scales the growth reward. \\
$\beta \geq 0$  & Penalty & Penalises maximum drawdown depth. \\
$\gamma \geq 0$ & Penalty & Penalises portfolio volatility. \\
$\delta \geq 0$ & Reward  & Rewards capital preservation probability. \\
$\varepsilon \geq 0$ & Reward & Rewards Calmar ratio. \\
$\zeta \geq 0$  & Penalty & Penalises excessive turnover costs. \\
$\eta \geq 0$   & Penalty & Penalises concentration. \\
$\theta \geq 0$ & Penalty & Penalises excessive leverage. \\
\bottomrule
\end{tabularx}
\end{table}

\begin{proposition}[Dominance of Quality over Quantity]
\label{prop:quality-dominance}
Let $W_0 = \$100{,}000$ be the common initial capital for both strategies.
Define the two strategies $\pi_1$ and $\pi_2$ via their wealth trajectories
and assume the following \emph{boundary data}:
\begin{enumerate}
  \item \textbf{Strategy $\pi_1$ (high peak, severe drawdown):}
        \begin{enumerate}
          \item Wealth reaches a peak $W_{\mathrm{peak},1} = \$160{,}000$
                before declining to a terminal value
                $W_{T,1} = \$70{,}000$.
          \item Net return:
                $G_1 = (W_{T,1} - W_0)/W_0
                     = (70{,}000 - 100{,}000)/100{,}000 = -0.30$.
          \item Maximum drawdown:
                $D_{\mathrm{dd},1}
                 = (W_{\mathrm{peak},1} - W_{T,1})/W_{\mathrm{peak},1}
                 = (160{,}000 - 70{,}000)/160{,}000 = 0.5625$.
        \end{enumerate}
  \item \textbf{Strategy $\pi_2$ (monotone growth, negligible drawdown):}
        \begin{enumerate}
          \item Wealth grows monotonically to a terminal value
                $W_{T,2} = \$135{,}000$ with no intermediate peak
                exceeding $W_{T,2}$.
          \item Net return:
                $G_2 = (W_{T,2} - W_0)/W_0
                     = (135{,}000 - 100{,}000)/100{,}000 = 0.35$.
          \item Maximum drawdown: Since the trajectory is monotonically
                non-decreasing, $W_{\mathrm{peak},2} = W_{T,2}$,
                so $D_{\mathrm{dd},2} = 0$ in the ideal case.  To be
                conservative and account for intra-period fluctuations,
                we take $D_{\mathrm{dd},2} = 0.02$ as an upper bound.
        \end{enumerate}
\end{enumerate}
Suppose, in addition, that the two strategies are \emph{comparable} in all
secondary objective terms, i.e.\ the following \textbf{symmetry assumptions}
hold:
\begin{enumerate}
  \setcounter{enumi}{2}
  \item $V_1 = V_2$ (equal portfolio volatility),
  \item $S_{\mathrm{surv},1} = S_{\mathrm{surv},2}$
        (equal capital-preservation probability),
  \item $R_{\mathrm{risk},1}$ and $R_{\mathrm{risk},2}$ are both
        well-defined (i.e.\ $D_{\mathrm{dd},i} > 0$ for $i=1,2$,
        ensuring the Calmar ratio is finite for each strategy),
  \item $T_{\mathrm{cost},1} = T_{\mathrm{cost},2}$ (equal turnover cost),
  \item $C_{\mathrm{conc},1} = C_{\mathrm{conc},2}$
        (equal concentration penalty),
  \item $L_{\mathrm{excess},1} = L_{\mathrm{excess},2}$
        (equal leverage penalty).
\end{enumerate}
Under these conditions, $J(\pi_2) > J(\pi_1)$ for all
$\alpha, \beta \geq 0$ with $(\alpha, \beta) \neq (0,0)$.
Moreover, $J(\pi_2) > J(\pi_1)$ holds in the degenerate case
$\beta = 0$ as well, driven solely by the growth differential.
\end{proposition}

\begin{proof}
The proof proceeds in a sequence of fully justified steps.

\medskip
\noindent\textbf{Step~1: Validate the boundary data for $\pi_1$.}

\begin{enumerate}
  \item Given $W_0 = 100{,}000$, $W_{\mathrm{peak},1} = 160{,}000$,
        and $W_{T,1} = 70{,}000$:
        \begin{align}
          G_1
          &\;=\; \frac{W_{T,1} - W_0}{W_0}
           \;=\; \frac{70{,}000 - 100{,}000}{100{,}000}
           \;=\; \frac{-30{,}000}{100{,}000}
           \;=\; -0.30.
        \end{align}
  \item The maximum drawdown is computed from the peak to the trough:
        \begin{align}
          D_{\mathrm{dd},1}
          &\;=\; \frac{W_{\mathrm{peak},1} - W_{T,1}}{W_{\mathrm{peak},1}}
           \;=\; \frac{160{,}000 - 70{,}000}{160{,}000}
           \;=\; \frac{90{,}000}{160{,}000}
           \;=\; 0.5625.
        \end{align}
  \item Validity check: $W_{\mathrm{peak},1} = 160{,}000 > W_{T,1}
        = 70{,}000 > 0$, so the drawdown formula is well-defined and
        $D_{\mathrm{dd},1} \in (0,1)$. \checkmark
  \item Since $D_{\mathrm{dd},1} = 0.5625 > 0$, the Calmar ratio
        $R_{\mathrm{risk},1} = G_1 / D_{\mathrm{dd},1}$ is well-defined
        (though negative):
        \begin{align}
          R_{\mathrm{risk},1}
          &\;=\; \frac{-0.30}{0.5625}
           \;=\; -\frac{8}{15}
           \;\approx\; -0.5\overline{3}.
        \end{align}
\end{enumerate}

\medskip
\noindent\textbf{Step~2: Validate the boundary data for $\pi_2$.}

\begin{enumerate}
  \item Given $W_0 = 100{,}000$ and $W_{T,2} = 135{,}000$:
        \begin{align}
          G_2
          &\;=\; \frac{W_{T,2} - W_0}{W_0}
           \;=\; \frac{135{,}000 - 100{,}000}{100{,}000}
           \;=\; \frac{35{,}000}{100{,}000}
           \;=\; 0.35.
        \end{align}
  \item Since the trajectory is monotonically non-decreasing by
        assumption, the running maximum satisfies
        $\sup_{t \leq T} W_t = W_{T,2}$, so the ideal drawdown is:
        \begin{align}
          D_{\mathrm{dd},2}^{\mathrm{ideal}}
          &\;=\; \frac{W_{T,2} - W_{T,2}}{W_{T,2}}
           \;=\; 0.
        \end{align}
  \item To be conservative and account for intra-period micro-fluctuations,
        we adopt the upper-bound estimate $D_{\mathrm{dd},2} = 0.02$.
        This satisfies $0 \leq D_{\mathrm{dd},2}^{\mathrm{ideal}}
        \leq D_{\mathrm{dd},2} = 0.02 \ll D_{\mathrm{dd},1} = 0.5625$,
        confirming that the conservative estimate does not undermine the
        dominance result.
  \item Since $D_{\mathrm{dd},2} = 0.02 > 0$, the Calmar ratio
        $R_{\mathrm{risk},2}$ is well-defined:
        \begin{align}
          R_{\mathrm{risk},2}
          &\;=\; \frac{G_2}{D_{\mathrm{dd},2}}
           \;=\; \frac{0.35}{0.02}
           \;=\; 17.5.
        \end{align}
\end{enumerate}

\medskip
\noindent\textbf{Step~3: Express the composite objective for each strategy.}

By Definition~\ref{def:composite-objective},
\begin{align}
  J(\pi_i)
  &\;=\; \alpha\,G_i
         \;-\; \beta\,D_{\mathrm{dd},i}
         \;-\; \gamma\,V_i
         \;+\; \delta\,S_{\mathrm{surv},i}
         \;+\; \varepsilon\,R_{\mathrm{risk},i} \notag\\
  &\quad  \;-\; \zeta\,T_{\mathrm{cost},i}
         \;-\; \eta\,C_{\mathrm{conc},i}
         \;-\; \theta\,L_{\mathrm{excess},i},
         \qquad i \in \{1, 2\}.
\end{align}

\medskip
\noindent\textbf{Step~4: Compute the objective difference.}

\begin{enumerate}
  \item Subtract $J(\pi_1)$ from $J(\pi_2)$ term by term:
        \begin{align}
          J(\pi_2) - J(\pi_1)
          &\;=\; \alpha(G_2 - G_1)
                 \;-\; \beta(D_{\mathrm{dd},2} - D_{\mathrm{dd},1})
                 \;-\; \gamma(V_2 - V_1) \notag\\
          &\quad  \;+\; \delta(S_{\mathrm{surv},2} - S_{\mathrm{surv},1})
                 \;+\; \varepsilon(R_{\mathrm{risk},2}
                        - R_{\mathrm{risk},1}) \notag\\
          &\quad  \;-\; \zeta(T_{\mathrm{cost},2} - T_{\mathrm{cost},1})
                 \;-\; \eta(C_{\mathrm{conc},2} - C_{\mathrm{conc},1})
                 \notag\\
          &\quad  \;-\; \theta(L_{\mathrm{excess},2}
                        - L_{\mathrm{excess},1}).
        \end{align}
  \item Apply the symmetry assumptions (Steps~1.3--1.8 of the
        proposition statement): all terms involving $V$, $S_{\mathrm{surv}}$,
        $T_{\mathrm{cost}}$, $C_{\mathrm{conc}}$, and $L_{\mathrm{excess}}$
        vanish identically.  This yields:
        \begin{align}
          J(\pi_2) - J(\pi_1)
          &\;=\; \alpha(G_2 - G_1)
                 \;-\; \beta(D_{\mathrm{dd},2} - D_{\mathrm{dd},1})
                 \;+\; \varepsilon(R_{\mathrm{risk},2}
                        - R_{\mathrm{risk},1}).
                 \label{eq:diff-reduced}
        \end{align}
\end{enumerate}

\medskip
\noindent\textbf{Step~5: Substitute numerical values.}

\begin{enumerate}
  \item Compute the growth differential:
        \begin{align}
          G_2 - G_1
          &\;=\; 0.35 - (-0.30)
           \;=\; 0.65.
        \end{align}
  \item Compute the drawdown differential:
        \begin{align}
          D_{\mathrm{dd},2} - D_{\mathrm{dd},1}
          &\;=\; 0.02 - 0.5625
           \;=\; -0.5425.
        \end{align}
  \item Compute the Calmar ratio differential:
        \begin{align}
          R_{\mathrm{risk},2} - R_{\mathrm{risk},1}
          &\;=\; 17.5 - \bigl(-0.5\overline{3}\bigr)
           \;=\; 17.5 + 0.5\overline{3}
           \;\approx\; 18.03\overline{3}.
        \end{align}
  \item Substitute into~\eqref{eq:diff-reduced}:
        \begin{align}
          J(\pi_2) - J(\pi_1)
          &\;=\; 0.65\,\alpha
                 \;-\; \beta\,(-0.5425)
                 \;+\; \varepsilon\,(18.03\overline{3}) \notag\\
          &\;=\; 0.65\,\alpha
                 \;+\; 0.5425\,\beta
                 \;+\; 18.03\overline{3}\,\varepsilon.
                 \label{eq:diff-final}
        \end{align}
\end{enumerate}

\medskip
\noindent\textbf{Step~6: Sign analysis and case enumeration.}

Since $\alpha, \beta, \varepsilon \geq 0$ by
Definition~\ref{def:composite-objective}, each coefficient in
\eqref{eq:diff-final} is non-negative and each numerical multiplier
is strictly positive ($0.65 > 0$, $0.5425 > 0$, $18.03\overline{3} > 0$).
We enumerate all non-trivial cases:

\begin{enumerate}
  \item \textbf{Case $\alpha > 0$, $\beta \geq 0$, $\varepsilon \geq 0$:}
        \begin{align}
          J(\pi_2) - J(\pi_1)
          &\;\geq\; 0.65\,\alpha \;>\; 0.
        \end{align}
        Dominance is driven by the growth differential alone.

  \item \textbf{Case $\alpha = 0$, $\beta > 0$, $\varepsilon \geq 0$:}
        \begin{align}
          J(\pi_2) - J(\pi_1)
          &\;\geq\; 0.5425\,\beta \;>\; 0.
        \end{align}
        Dominance is driven by the drawdown differential alone.

  \item \textbf{Case $\alpha = 0$, $\beta = 0$, $\varepsilon > 0$:}
        \begin{align}
          J(\pi_2) - J(\pi_1)
          &\;=\; 18.03\overline{3}\,\varepsilon \;>\; 0.
        \end{align}
        Dominance is driven by the Calmar ratio differential alone.

  \item \textbf{Degenerate boundary case $\beta = 0$, $\varepsilon = 0$,
        $\alpha > 0$:}\\
        Even with no drawdown penalty and no Calmar weighting,
        $J(\pi_2) - J(\pi_1) = 0.65\,\alpha > 0$, so $\pi_2$ still
        dominates.

  \item \textbf{Excluded case $\alpha = \beta = \varepsilon = 0$:}
        The reduced objective~\eqref{eq:diff-reduced} becomes
        identically zero under the symmetry assumptions; this case is
        excluded since $(\alpha, \beta) \neq (0,0)$ is required by
        hypothesis.  (More generally, if all eight weights are zero
        the objective is degenerate and comparison is undefined.)
\end{enumerate}

\medskip
\noindent\textbf{Step~7: Conclusion.}

\begin{enumerate}
  \item From~\eqref{eq:diff-final} and the sign analysis in Step~6,
        for all $\alpha, \beta, \varepsilon \geq 0$ with
        $(\alpha,\beta,\varepsilon) \neq (0,0,0)$:
        \begin{align}
          J(\pi_2) - J(\pi_1)
          &\;=\; 0.65\,\alpha + 0.5425\,\beta
                 + 18.03\overline{3}\,\varepsilon
          \;>\; 0.
        \end{align}
  \item Under the stated hypotheses, all secondary objective terms
        cancel exactly, so no additional conditions on $\gamma$,
        $\delta$, $\zeta$, $\eta$, or $\theta$ are required.
  \item Therefore, $J(\pi_2) > J(\pi_1)$, and $\pi_2$ strictly
        dominates $\pi_1$ under the composite objective
        \eqref{eq:composite-objective} whenever
        $(\alpha,\beta,\varepsilon) \neq (0,0,0)$.
\end{enumerate}
\qed
\end{proof}

% ==============================================================
\chapter{Constrained Portfolio Optimisation}
\label{ch:constrained-opt}
% ==============================================================

\section{Feasible Capital Deployment Set}

\begin{definition}[Capital Deployment Vector]
\label{def:deployment-vector}
A \emph{capital deployment vector} at time $t$ is a vector
$\bfx = (x_1,\ldots,x_n)^\top \in \mathbb{R}^n$, where $n$ is the
number of investable instruments and $x_j$ denotes the dollar amount
allocated to instrument $j$.  Positive $x_j$ denotes a long position.
\end{definition}

\begin{definition}[Feasible Deployment Set]
\label{def:feasible-set}
\label{asm:feasible-set}
The \emph{feasible deployment set} is
\begin{align}
  \Pi_F
  &\;=\;
  \Bigl\{
    \bfx \in \mathbb{R}^n
    \;\Big|\;
    &\sum_{j=1}^n x_j \leq W_t^{\mathrm{deployable}},
    \label{eq:budget-constr}\\
    &x_j \geq 0 \quad \forall\,j,
    \label{eq:long-only-constr}\\
    &x_j / W_t \leq w_{\max} \quad \forall\,j,
    \label{eq:concentration-constr}\\
    &\sum_{j \in \mathrm{bucket}_k} x_j \leq B_k \quad \forall\,k,
    \label{eq:bucket-constr}\\
    &x_{\mathrm{cash}} \geq \ell_{\min},
    \label{eq:liquidity-constr}\\
    &\mathrm{CVaR}_\alpha(\bfx) \leq \bar{C}
    \label{eq:cvar-constr}
  \Bigr\},
\end{align}
where $W_t^{\mathrm{deployable}} = W_t - \ell_{\min}$,
$w_{\max} \in (1/n, 1)$ is the single-instrument concentration cap,
$B_k$ is the bucket limit, and $\bar{C} > 0$ is the CVaR budget.
\end{definition}

\begin{theorem}[Feasibility and Convexity of $\Pi_F$]
\label{thm:feasible-convex}
Under Assumptions~\ref{asm:initial-capital}--\ref{asm:bucket-count}
and the definitions above, $\Pi_F$ is non-empty, compact, and convex.
\end{theorem}

\begin{proof}
We verify each claimed property in turn, validating every assumption,
edge case, and boundary condition explicitly.

\medskip
\noindent\textbf{Preliminary: Assumption Validation.}

\begin{enumerate}
  \item By Assumption~\ref{asm:initial-capital}, $W_t > 0$, so
        $W_t^{\mathrm{deployable}} = W_t - \ell_{\min} > 0$ provided
        $\ell_{\min} < W_t$, which is imposed by
        Assumption~\ref{asm:liquidity-reserve}.
        Hence the budget constraint~\eqref{eq:budget-constr} admits
        strictly positive allocations.

  \item By Assumption~\ref{asm:bucket-count}, the number of buckets
        $K \geq 1$ is finite, and each bucket limit satisfies
        $B_k > 0$. Consequently, the right-hand sides of all
        constraints are finite and well-defined.

  \item The concentration cap satisfies $w_{\max} \in (1/n, 1)$ by
        Definition~\ref{def:feasible-set}, so the per-instrument upper
        bound $w_{\max} W_t$ is a positive, finite quantity for
        every $j \in \{1,\ldots,n\}$.

  \item The CVaR budget satisfies $\bar{C} > 0$ by definition, so
        the CVaR constraint~\eqref{eq:cvar-constr} is non-vacuous.
\end{enumerate}

\medskip
\noindent\textbf{Property~1: Non-emptiness.}

\begin{enumerate}
  \item Set $\bfx = \mathbf{0} \in \mathbb{R}^n$, i.e.,
        $x_j = 0$ for all $j \in \{1,\ldots,n\}$.

  \item \textit{Budget constraint~\eqref{eq:budget-constr}:}
        \begin{align}
          \sum_{j=1}^n x_j
          &\;=\; \sum_{j=1}^n 0
           \;=\; 0
           \;\leq\; W_t^{\mathrm{deployable}},
        \end{align}
        which holds since $W_t^{\mathrm{deployable}} > 0$ (Step~1 above).

  \item \textit{Long-only constraint~\eqref{eq:long-only-constr}:}
        \begin{align}
          x_j &\;=\; 0 \;\geq\; 0 \quad \forall\,j,
        \end{align}
        satisfied trivially at the boundary of the non-negativity
        orthant.

  \item \textit{Concentration constraint~\eqref{eq:concentration-constr}:}
        \begin{align}
          x_j / W_t &\;=\; 0 \;\leq\; w_{\max} \quad \forall\,j,
        \end{align}
        satisfied since $w_{\max} > 0$.

  \item \textit{Bucket constraint~\eqref{eq:bucket-constr}:}
        For every bucket $k \in \{1,\ldots,K\}$:
        \begin{align}
          \sum_{j \in \mathrm{bucket}_k} x_j
          &\;=\; 0 \;\leq\; B_k,
        \end{align}
        satisfied since $B_k > 0$.

  \item \textit{Liquidity constraint~\eqref{eq:liquidity-constr}:}
        The cash reserve $x_{\mathrm{cash}}$ is maintained separately
        from $\bfx$ and is set to $\ell_{\min}$ by construction; hence
        $x_{\mathrm{cash}} = \ell_{\min} \geq \ell_{\min}$.

  \item \textit{CVaR constraint~\eqref{eq:cvar-constr}:}
        The portfolio loss under $\bfx = \mathbf{0}$ is identically
        zero in every scenario.  Therefore:
        \begin{align}
          \mathrm{CVaR}_\alpha(\mathbf{0})
          &\;=\; 0 \;\leq\; \bar{C},
        \end{align}
        satisfied since $\bar{C} > 0$.

  \item Since all seven constraints are satisfied, $\mathbf{0} \in \Pi_F$,
        confirming $\Pi_F \neq \emptyset$.
\end{enumerate}

\medskip
\noindent\textbf{Property~2: Convexity.}

\begin{enumerate}
  \item For convexity of $\Pi_F$, it suffices to show that each
        constraint defines a convex set in $\mathbb{R}^n$, and then
        invoke closure under finite intersection.

  \item \textit{Budget constraint~\eqref{eq:budget-constr}:}
        Define $f_0(\bfx) = \sum_{j=1}^n x_j$.  This is a linear
        (hence convex) function.  The set
        $\{\bfx : f_0(\bfx) \leq W_t^{\mathrm{deployable}}\}$
        is a closed half-space in $\mathbb{R}^n$, which is convex.

  \item \textit{Long-only constraints~\eqref{eq:long-only-constr}:}
        For each $j$, the set $\{x_j \geq 0\}$ is a closed half-space
        in $\mathbb{R}^n$, hence convex.  The intersection over all
        $j$ is the non-negative orthant $\mathbb{R}_{\geq 0}^n$,
        which is convex.

  \item \textit{Concentration constraints~\eqref{eq:concentration-constr}:}
        For each $j$, the map $\bfx \mapsto x_j / W_t$ is linear.
        The set $\{x_j / W_t \leq w_{\max}\}$ is a closed half-space,
        hence convex.

  \item \textit{Bucket constraints~\eqref{eq:bucket-constr}:}
        For each bucket $k$, the map
        $\bfx \mapsto \sum_{j \in \mathrm{bucket}_k} x_j$ is linear.
        The set
        $\bigl\{\sum_{j \in \mathrm{bucket}_k} x_j \leq B_k\bigr\}$
        is a closed half-space, hence convex.

  \item \textit{Liquidity constraint~\eqref{eq:liquidity-constr}:}
        The set $\{x_{\mathrm{cash}} \geq \ell_{\min}\}$ is a closed
        half-space in the coordinate corresponding to cash, hence convex.

  \item \textit{CVaR constraint~\eqref{eq:cvar-constr}:}
        By the Rockafellar--Uryasev representation
        (Proposition~6.1 of the companion document), for a fixed
        confidence level $\alpha \in (0,1)$:
        \begin{align}
          \mathrm{CVaR}_\alpha(\bfx)
          &\;=\; \min_{\zeta \in \mathbb{R}}
            \left\{
              \zeta + \frac{1}{1-\alpha}
              \,\mathbb{E}\!\left[(\mathcal{L}(\bfx) - \zeta)^+\right]
            \right\},
        \end{align}
        where $\mathcal{L}(\bfx)$ is the portfolio loss and
        $(\cdot)^+ = \max(\cdot, 0)$.  The function
        $\bfx \mapsto \mathrm{CVaR}_\alpha(\bfx)$ is convex in $\bfx$
        because it is the minimum over an affine family (in $\zeta$)
        of convex functions of $\bfx$.  Therefore
        $\{\bfx : \mathrm{CVaR}_\alpha(\bfx) \leq \bar{C}\}$
        is a convex sublevel set.

  \item \textit{Finite intersection:}
        $\Pi_F$ is the intersection of the sets defined in
        Steps~2--6 above.  The total number of constraint sets is
        $1 + n + n + K + 1 + 1 = 2n + K + 3$, which is finite under
        Assumption~\ref{asm:bucket-count}.  The intersection of
        finitely many convex sets is convex.  Therefore $\Pi_F$ is
        convex.
\end{enumerate}

\medskip
\noindent\textbf{Property~3: Boundedness.}

\begin{enumerate}
  \item From the long-only constraint~\eqref{eq:long-only-constr},
        $x_j \geq 0$ for every $j$.

  \item From the budget constraint~\eqref{eq:budget-constr} and
        non-negativity:
        \begin{align}
          0 \;\leq\; x_j
          &\;\leq\; \sum_{i=1}^n x_i
           \;\leq\; W_t^{\mathrm{deployable}}
           \;\leq\; W_t,
        \end{align}
        for every $j \in \{1,\ldots,n\}$.

  \item Therefore every feasible $\bfx$ lies in the box:
        \begin{align}
          \Pi_F &\;\subseteq\; [0,\, W_t]^n \;\subset\; \mathbb{R}^n.
        \end{align}

  \item The box $[0, W_t]^n$ is bounded (its diameter equals
        $W_t \sqrt{n} < \infty$ since $W_t < \infty$ by
        Assumption~\ref{asm:initial-capital} and $n < \infty$ by
        Assumption~\ref{asm:bucket-count}).  Hence $\Pi_F$ is bounded.

  \item \textit{Boundary case $W_t^{\mathrm{deployable}} = 0$:}
        If $\ell_{\min} = W_t$ then $W_t^{\mathrm{deployable}} = 0$
        and the budget constraint forces $\bfx = \mathbf{0}$, which
        is still bounded (a single-point set).
        This case is excluded by Assumption~\ref{asm:liquidity-reserve}
        ($\ell_{\min} < W_t$), but boundedness holds regardless.
\end{enumerate}

\medskip
\noindent\textbf{Property~4: Closedness.}

\begin{enumerate}
  \item Each constraint in $\Pi_F$ takes the form
        $g(\bfx) \leq c$ or $g(\bfx) \geq c$ for a continuous
        function $g : \mathbb{R}^n \to \mathbb{R}$ and a finite
        constant $c \in \mathbb{R}$.

  \item For any continuous $g$ and finite $c$, the sets
        $\{\bfx : g(\bfx) \leq c\}$ and $\{\bfx : g(\bfx) \geq c\}$
        are closed, being preimages of closed sets under a continuous
        map.

  \item The functions involved are:
        \begin{enumerate}
          \item $g_0(\bfx) = \sum_j x_j$ (linear, continuous);
          \item $g_j(\bfx) = x_j$ for each $j$ (linear, continuous);
          \item $g_j^c(\bfx) = x_j / W_t$ for each $j$
                (linear, continuous; $W_t > 0$ by
                Assumption~\ref{asm:initial-capital});
          \item $g_k^b(\bfx) = \sum_{j \in \mathrm{bucket}_k} x_j$
                for each $k$ (linear, continuous);
          \item $g_{\ell}(\bfx) = x_{\mathrm{cash}}$
                (linear, continuous);
          \item $g_{\mathrm{CVaR}}(\bfx) =
                \mathrm{CVaR}_\alpha(\bfx)$ (convex and
                lower semicontinuous, hence continuous on the
                interior of its domain, and closed as a constraint).
        \end{enumerate}

  \item The intersection of finitely many closed sets is closed.
        Therefore $\Pi_F$ is closed.
\end{enumerate}

\medskip
\noindent\textbf{Property~5: Compactness.}

\begin{enumerate}
  \item From Property~3, $\Pi_F \subseteq [0, W_t]^n$, so $\Pi_F$
        is bounded.

  \item From Property~4, $\Pi_F$ is closed.

  \item $\Pi_F$ is a closed and bounded subset of the finite-dimensional
        space $\mathbb{R}^n$.  By the Heine--Borel theorem, every
        closed and bounded subset of $\mathbb{R}^n$ is compact.
        Therefore $\Pi_F$ is compact.

  \item \textit{Edge case $n = 1$:} With a single instrument,
        $\Pi_F \subseteq [0, W_t] \subset \mathbb{R}$, which is a
        closed bounded interval, hence compact.  All five properties
        hold in this degenerate case.

  \item \textit{Edge case $K = 0$ (no buckets):}
        The bucket constraint~\eqref{eq:bucket-constr} is vacuous.
        All remaining constraints still define a non-empty, convex,
        closed, bounded (hence compact) set, so the theorem
        conclusion is unaffected.
\end{enumerate}

\medskip
\noindent\textbf{Summary.}

\begin{enumerate}
  \item $\Pi_F \neq \emptyset$ (Property~1: $\mathbf{0} \in \Pi_F$).
  \item $\Pi_F$ is convex (Property~2: finite intersection of convex
        half-spaces and a convex CVaR sublevel set).
  \item $\Pi_F$ is bounded (Property~3: contained in $[0,W_t]^n$).
  \item $\Pi_F$ is closed (Property~4: finite intersection of closed
        sets).
  \item $\Pi_F$ is compact (Property~5: closed and bounded in
        $\mathbb{R}^n$, Heine--Borel).
\end{enumerate}
\qed
\end{proof}

\begin{theorem}[Existence and Uniqueness of Optimal Deployment]
\label{thm:concentration-suboptimal}
Let the following assumptions hold:
\begin{enumerate}
  \item[\textup{(A1)}] The feasible set $\Pi_F \subset \mathbb{R}^n$
        is as constructed in Theorem~\ref{thm:feasible-convex},
        hence non-empty, convex, and compact.
  \item[\textup{(A2)}] The covariance matrix $\bfSigma \in
        \mathbb{R}^{n \times n}$ of instrument returns is strictly
        positive definite, i.e.\ $\bfSigma \succ 0$.
  \item[\textup{(A3)}] The objective $J : \mathbb{R}^n \to
        \mathbb{R}$ is of the form
        \begin{align}
          J(\bfx)
            &\;=\; \alpha\, G(\bfx)
               \;-\; \gamma\, V(\bfx)
               \;+\; \delta\, S_{\mathrm{surv}}(\bfx)
               \;-\; P(\bfx),
          \label{eq:J-components}
        \end{align}
        where $G$, $S_{\mathrm{surv}}$ are concave in $\bfx$;
        $V(\bfx) = \bfx^\top \bfSigma \bfx$ is the portfolio
        variance; $P(\bfx) \geq 0$ is a convex penalty; and
        $\alpha, \gamma, \delta > 0$.
  \item[\textup{(A4)}] The weight on the variance term satisfies
        $\gamma > 0$.
\end{enumerate}
Under assumptions \textup{(A1)--(A4)}, the risk-adjusted
optimisation problem
\begin{equation}
  \bfx^* \;=\; \arg\max_{\bfx \in \Pi_F}\; J(\bfx)
  \label{eq:deployment-opt}
\end{equation}
admits a \emph{unique} solution $\bfx^* \in \Pi_F$.
Moreover, for every index $k \in \{1,\ldots,n\}$, the
single-instrument portfolio $\bfx^{(k)} = W_t\,\bfe_k$ (all
capital concentrated in instrument $k$) is \emph{strictly
sub-optimal}, i.e.\ $J(\bfx^*) > J(\bfx^{(k)})$, unless for
every $j \neq k$ the off-diagonal entry satisfies
$\sigma_{kj} \geq \sigma_{kk}$.
\end{theorem}

\begin{proof}
We establish each conclusion in turn, validating every
assumption at the step where it is first used.

\medskip
\begin{enumerate}

%------------------------------------------------------------
\item \textbf{Continuity of $J$ on $\Pi_F$.}

\begin{enumerate}
  \item[\textup{(1a)}] \textit{Domain.}
        By assumption~(A1) and
        Theorem~\ref{thm:feasible-convex},
        $\Pi_F \subset \mathbb{R}^n$ is compact.

  \item[\textup{(1b)}] \textit{Each component.}
        \begin{enumerate}
          \item[\textup{(i)}] $G(\bfx)$ is concave on $\mathbb{R}^n$
                by~(A3), hence continuous.
          \item[\textup{(ii)}] $V(\bfx) = \bfx^\top \bfSigma \bfx$
                is a quadratic form with bounded matrix $\bfSigma$,
                hence continuous.
          \item[\textup{(iii)}] $S_{\mathrm{surv}}(\bfx)$ is concave
                by~(A3), hence continuous.
          \item[\textup{(iv)}] $P(\bfx)$ is convex by~(A3), hence
                continuous.
        \end{enumerate}

  \item[\textup{(1c)}] \textit{Conclusion.}
        $J$ is a finite linear combination of continuous functions,
        so $J$ is continuous on $\mathbb{R}^n$, and in particular
        on $\Pi_F$.
\end{enumerate}

%------------------------------------------------------------
\item \textbf{Existence of a maximiser.}

\begin{enumerate}
  \item[\textup{(2a)}] \textit{Premises verified.}
        $\Pi_F$ is compact (step~1a) and $J$ is continuous
        on $\Pi_F$ (step~1c).

  \item[\textup{(2b)}] \textit{Application of Weierstrass.}
        By the Extreme Value Theorem (Weierstrass, 1861), every
        continuous real-valued function on a non-empty compact
        metric space attains its supremum.  Hence there exists
        $\bfx^* \in \Pi_F$ such that
        \begin{align}
          J(\bfx^*) &\;=\; \sup_{\bfx \in \Pi_F} J(\bfx).
          \label{eq:existence-max}
        \end{align}

  \item[\textup{(2c)}] \textit{Boundary check.}
        The supremum is attained; the set of maximisers
        $\mathcal{M} := \arg\max_{\bfx \in \Pi_F} J(\bfx)$
        is non-empty.
\end{enumerate}

%------------------------------------------------------------
\item \textbf{Strict concavity of $J$.}

\begin{enumerate}
  \item[\textup{(3a)}] \textit{Variance term.}
        Since $\bfSigma \succ 0$ by~(A2), for any
        $\bfx, \bfy \in \mathbb{R}^n$ with $\bfx \neq \bfy$ and
        any $\lambda \in (0,1)$:
        \begin{align}
          V\!\left(\lambda\bfx + (1-\lambda)\bfy\right)
            &\;=\; \bigl(\lambda\bfx+(1-\lambda)\bfy\bigr)^\top
                   \bfSigma
                   \bigl(\lambda\bfx+(1-\lambda)\bfy\bigr)
          \nonumber\\
            &\;<\; \lambda\,\bfx^\top\bfSigma\bfx
                   \;+\;(1-\lambda)\,\bfy^\top\bfSigma\bfy
          \nonumber\\
            &\;=\; \lambda\,V(\bfx)+(1-\lambda)\,V(\bfy).
          \label{eq:V-strictly-convex}
        \end{align}
        Hence $V$ is \emph{strictly convex}, so $-\gamma V$ is
        \emph{strictly concave} for $\gamma > 0$ (assumption~(A4)).

  \item[\textup{(3b)}] \textit{Remaining terms.}
        \begin{enumerate}
          \item[\textup{(i)}] $\alpha G$ is concave (concave times
                positive scalar), hence concave.
          \item[\textup{(ii)}] $\delta S_{\mathrm{surv}}$ is
                concave for the same reason.
          \item[\textup{(iii)}] $-P$ is concave since $P$ is
                convex.
        \end{enumerate}

  \item[\textup{(3c)}] \textit{Sum.}
        The sum of a strictly concave function and finitely many
        concave functions is strictly concave.  Explicitly, for
        $\bfx \neq \bfy$ and $\lambda \in (0,1)$:
        \begin{align}
          J\!\left(\lambda\bfx+(1-\lambda)\bfy\right)
            &\;=\; \alpha G\!\left(\lambda\bfx+(1-\lambda)\bfy\right)
            \nonumber\\
            &\quad -\,\gamma V\!\left(\lambda\bfx+(1-\lambda)\bfy\right)
            \nonumber\\
            &\quad +\,\delta
               S_{\mathrm{surv}}\!\left(\lambda\bfx+(1-\lambda)\bfy\right)
            \nonumber\\
            &\quad -\, P\!\left(\lambda\bfx+(1-\lambda)\bfy\right).
          \label{eq:J-split}
        \end{align}
        Applying the concavity inequalities to each term and the
        \emph{strict} inequality from~\eqref{eq:V-strictly-convex}:
        \begin{align}
          J\!\left(\lambda\bfx+(1-\lambda)\bfy\right)
            &\;>\; \lambda J(\bfx) + (1-\lambda) J(\bfy).
          \label{eq:J-strictly-concave}
        \end{align}
        Hence $J$ is strictly concave on $\mathbb{R}^n$.
\end{enumerate}

%------------------------------------------------------------
\item \textbf{Uniqueness of the maximiser.}

\begin{enumerate}
  \item[\textup{(4a)}] \textit{Suppose for contradiction} that
        there exist two distinct maximisers
        $\bfx^*, \bfy^* \in \mathcal{M}$ with $\bfx^* \neq \bfy^*$.

  \item[\textup{(4b)}] \textit{Midpoint feasibility.}
        Since $\Pi_F$ is convex (assumption~(A1)) and
        $\bfx^*, \bfy^* \in \Pi_F$, the midpoint
        $\bfz := \tfrac{1}{2}\bfx^* + \tfrac{1}{2}\bfy^*
        \in \Pi_F$.

  \item[\textup{(4c)}] \textit{Strict improvement.}
        By~\eqref{eq:J-strictly-concave} with
        $\lambda = \tfrac{1}{2}$:
        \begin{align}
          J(\bfz)
            &\;>\; \tfrac{1}{2} J(\bfx^*) + \tfrac{1}{2} J(\bfy^*)
             \;=\; J(\bfx^*),
          \label{eq:midpoint-improvement}
        \end{align}
        since $J(\bfx^*) = J(\bfy^*)$ by the assumption that
        both are maximisers.

  \item[\textup{(4d)}] \textit{Contradiction.}
        Inequality~\eqref{eq:midpoint-improvement} contradicts
        the optimality of $\bfx^*$.  Hence $\mathcal{M}$ is a
        singleton, i.e.\ the maximiser is unique.
\end{enumerate}

%------------------------------------------------------------
\item \textbf{Strict sub-optimality of concentrated portfolios.}

Fix any index $k \in \{1,\ldots,n\}$ and define the concentrated
portfolio $\bfx^{(k)} := W_t\,\bfe_k$.

\begin{enumerate}
  \item[\textup{(5a)}] \textit{Feasibility of $\bfx^{(k)}$.}
        We verify $\bfx^{(k)} \in \Pi_F$.
        All capital is in instrument~$k$, so the budget constraint
        $\sum_i x_i = W_t$ is satisfied, all components are
        non-negative, and the position satisfies the box constraint
        $x_k = W_t \leq W_t$ and $x_j = 0$ for $j \neq k$.
        CVaR and bucket constraints are trivially satisfied at a
        single non-zero position.  Hence $\bfx^{(k)} \in \Pi_F$.

  \item[\textup{(5b)}] \textit{Parametric perturbation.}
        For $j \neq k$ and $\varepsilon \in [0,1]$, define
        \begin{align}
          \bfx_\varepsilon
            &\;:=\; (1-\varepsilon)\,W_t\,\bfe_k
                  \;+\; \varepsilon\,W_t\,\bfe_j.
          \label{eq:perturbation}
        \end{align}
        By convexity of $\Pi_F$ and step~(5a),
        $\bfx_\varepsilon \in \Pi_F$ for all
        $\varepsilon \in [0,1]$.

  \item[\textup{(5c)}] \textit{Variance at the perturbation.}
        Expanding the quadratic form:
        \begin{align}
          V(\bfx_\varepsilon)
            &\;=\; \bfx_\varepsilon^\top \bfSigma\, \bfx_\varepsilon
          \nonumber\\
            &\;=\; W_t^2\,\Bigl[
                   (1-\varepsilon)^2\,\bfe_k^\top\bfSigma\bfe_k
                   \;+\; 2\varepsilon(1-\varepsilon)\,
                         \bfe_k^\top\bfSigma\bfe_j
                   \;+\; \varepsilon^2\,\bfe_j^\top\bfSigma\bfe_j
                   \Bigr]
          \nonumber\\
            &\;=\; W_t^2\,\Bigl[
                   (1-\varepsilon)^2\,\sigma_{kk}
                   \;+\; 2\varepsilon(1-\varepsilon)\,\sigma_{kj}
                   \;+\; \varepsilon^2\,\sigma_{jj}
                   \Bigr].
          \label{eq:V-perturbed}
        \end{align}

  \item[\textup{(5d)}] \textit{Directional derivative of $V$.}
        Differentiating~\eqref{eq:V-perturbed} with respect to
        $\varepsilon$ and evaluating at $\varepsilon = 0$:
        \begin{align}
          \left.\frac{d}{d\varepsilon} V(\bfx_\varepsilon)
          \right|_{\varepsilon=0}
            &\;=\; W_t^2\,\Bigl[
                   -2\sigma_{kk} \;+\; 2\sigma_{kj}
                   \Bigr]
          \nonumber\\
            &\;=\; 2\,W_t^2\,\bigl(\sigma_{kj} - \sigma_{kk}\bigr).
          \label{eq:dV-dep}
        \end{align}

  \item[\textup{(5e)}] \textit{Effect on $J$.}
        Since $J$ depends on $V$ through $-\gamma V$ and all other
        terms of $J$ are continuously differentiable in $\bfx$
        (concavity implies local Lipschitz continuity, hence
        differentiability almost everywhere; here we evaluate the
        directional derivative), the directional derivative of
        $J$ along the perturbation direction $\bfe_j - \bfe_k$
        at $\bfx^{(k)}$ satisfies:
        \begin{align}
          \left.\frac{d}{d\varepsilon} J(\bfx_\varepsilon)
          \right|_{\varepsilon=0}
            &\;=\; \alpha\,\nabla G(\bfx^{(k)})^\top
                   W_t(\bfe_j-\bfe_k)
            \nonumber\\
            &\quad -\,\gamma \cdot 2 W_t^2
                   \bigl(\sigma_{kj}-\sigma_{kk}\bigr)
            \nonumber\\
            &\quad +\,\delta\,
                   \nabla S_{\mathrm{surv}}(\bfx^{(k)})^\top
                   W_t(\bfe_j-\bfe_k)
            \nonumber\\
            &\quad -\,\nabla P(\bfx^{(k)})^\top
                   W_t(\bfe_j-\bfe_k).
          \label{eq:dJ-dep}
        \end{align}

  \item[\textup{(5f)}] \textit{Sufficient condition for strict
        improvement.}
        If $\sigma_{kj} < \sigma_{kk}$, then the variance term
        in~\eqref{eq:dJ-dep} contributes
        $-\gamma \cdot 2 W_t^2(\sigma_{kj}-\sigma_{kk}) > 0$,
        i.e.\ a strict reduction in variance is achievable.
        Provided the combined gradient contributions from $G$,
        $S_{\mathrm{surv}}$, and $P$ do not dominate this
        improvement (which follows for sufficiently small
        $\varepsilon > 0$ by continuity of these terms and the
        fact that the variance term produces a first-order gain),
        we obtain:
        \begin{align}
          J(\bfx_\varepsilon) &\;>\; J(\bfx^{(k)})
          \quad \text{for all sufficiently small } \varepsilon > 0.
          \label{eq:strict-improvement}
        \end{align}

  \item[\textup{(5g)}] \textit{Existence of an improving direction.}
        The condition $\sigma_{kj} < \sigma_{kk}$ for \emph{some}
        $j \neq k$ is equivalent to the assertion that at least one
        off-diagonal entry in column~$k$ of $\bfSigma$ is strictly
        less than $\sigma_{kk}$.  This is the \emph{generic} case;
        it fails only if $\sigma_{kj} \geq \sigma_{kk}$ for all
        $j \neq k$.

  \item[\textup{(5h)}] \textit{Conclusion on concentration.}
        When the condition of step~(5g) holds for some $j$,
        inequality~\eqref{eq:strict-improvement} shows
        $\bfx^{(k)}$ is not the maximiser of $J$ over $\Pi_F$.
        Since the unique maximiser $\bfx^*$ is distinct from
        $\bfx^{(k)}$, we conclude:
        \begin{align}
          J(\bfx^*) &\;>\; J(\bfx^{(k)}).
          \label{eq:suboptimality}
        \end{align}

  \item[\textup{(5i)}] \textit{Boundary and degeneracy cases.}
        \begin{enumerate}
          \item[\textup{(i)}] If $n = 1$, there is only one
                instrument; the concentrated portfolio is the only
                feasible point, so $\bfx^* = \bfx^{(1)}$ and the
                sub-optimality claim is vacuously void.
          \item[\textup{(ii)}] If $W_t = 0$, then $\Pi_F =
                \{\mathbf{0}\}$ and both $\bfx^*$ and $\bfx^{(k)}$
                coincide; the claim is again vacuously void.
          \item[\textup{(iii)}] If $\bfx^{(k)} \notin \Pi_F$
                (e.g.\ violated by bucket or CVaR constraints),
                the sub-optimality conclusion holds trivially
                since the concentrated portfolio is infeasible.
        \end{enumerate}
\end{enumerate}

%------------------------------------------------------------
\item \textbf{Consolidated conclusion.}

\begin{enumerate}
  \item[\textup{(6a)}] Steps~1--2 establish that $\mathcal{M}
        \neq \emptyset$: a maximiser exists.
  \item[\textup{(6b)}] Steps~3--4 establish $|\mathcal{M}| = 1$:
        the maximiser is unique.
  \item[\textup{(6c)}] Steps~5a--5h establish strict
        sub-optimality of every concentrated portfolio
        $\bfx^{(k)}$ under the stated non-degeneracy condition.
  \item[\textup{(6d)}] Step~5i verifies that all boundary and
        degenerate cases are handled without violating any claim.
\end{enumerate}

\end{enumerate}
\qed
\end{proof}

% ==============================================================
\chapter{Kelly-Optimal Position Sizing Under Agent Uncertainty}
\label{ch:kelly}
% ==============================================================

\section{Kelly Fraction with Uncertainty Adjustment}

% ------------------------------------------------------------------
\section{Assumptions and Prerequisites}
% ------------------------------------------------------------------

Before stating the main definitions and results of this chapter,
we enumerate every assumption that is required and verify that all
hypotheses are stated explicitly, so that no implicit condition is
left unvalidated.

\begin{enumerate}
  \item[\textup{(A1)}] \textbf{Agent output validity.}
        Every agent output tuple
        $\calA_i = (s_i, c_i, u_i, d_i, p_i^+, p_i^-,
        \Delta t_i, r_i)$ satisfies
        Definition~\ref{def:agent-output}.
        In particular:
        \begin{enumerate}
          \item[\textup{(A1a)}] $c_i \in [0,1]$ is the calibration
                score (confidence).
          \item[\textup{(A1b)}] $u_i \in [0,1)$ is the upward
                uncertainty discount.
          \item[\textup{(A1c)}] $d_i \in [0,1)$ is the downward
                uncertainty discount, with $u_i + d_i < 1$ whenever
                both are positive.
          \item[\textup{(A1d)}] $p_i^+ \in (0,1)$ and
                $p_i^- \in (0,1)$ with $p_i^+ + p_i^- = 1$.
        \end{enumerate}

  \item[\textup{(A2)}] \textbf{Channel validity.}
        Assumption~\ref{asm:channel-validity} holds:
        the channel factor $\phi_i = c_i(1-u_i)(1-d_i)$
        satisfies $\phi_i \in [0,1]$ for all $i$.

  \item[\textup{(A3)}] \textbf{Odds ratio positivity.}
        The net odds ratio $b > 0$ is strictly positive and
        finite; i.e.\ $b \in (0,\infty)$.

  \item[\textup{(A4)}] \textbf{Position cap validity.}
        The pre-specified maximum position fraction satisfies
        $f_{\max} \in (0,1)$.

  \item[\textup{(A5)}] \textbf{Ensemble uncertainty range.}
        The ensemble uncertainty $U \in [0,1]$ as defined in
        Definition~\ref{def:allocator-state}.
\end{enumerate}

% ------------------------------------------------------------------
\section{Kelly Fraction with Uncertainty Adjustment}
% ------------------------------------------------------------------

\begin{definition}[Channel Factor]
\label{def:channel-factor}
Given an agent output tuple $\calA_i$ satisfying
Assumption~\textup{(A1)}, the \emph{channel factor} is
\begin{align}
  \phi_i
  &\;=\; c_i\,(1 - u_i)\,(1 - d_i).
  \label{eq:channel-factor}
\end{align}
By Assumptions~\textup{(A1a)}--\textup{(A1c)}, each multiplicand
lies in $[0,1]$, so $\phi_i \in [0,1]$.
\end{definition}

\begin{definition}[Effective Win Probability]
\label{def:effective-win-prob}
Given the agent output tuple $\calA_i = (s_i, c_i, u_i, d_i,
p_i^+, p_i^-, \Delta t_i, r_i)$ satisfying
Assumption~\textup{(A1)}, and the channel factor $\phi_i$ from
Definition~\ref{def:channel-factor}, the \emph{effective win
probability} for a long position is
\begin{align}
  p_{\mathrm{eff}}
  &\;=\; \phi_i\, p_i^+
         \;+\; \tfrac{1}{2}\bigl(1 - \phi_i\bigr).
  \label{eq:effective-win-prob}
\end{align}
\textbf{Interpretation.}
\begin{enumerate}
  \item[\textup{(i)}] When $\phi_i = 1$ (agent is fully
        confident, uncorrupted, undegraded),
        $p_{\mathrm{eff}} = p_i^+$: the agent's own estimate
        is used without attenuation.
  \item[\textup{(ii)}] When $\phi_i = 0$ (agent has zero net
        signal quality, e.g.\ $c_i = 0$),
        $p_{\mathrm{eff}} = \tfrac{1}{2}$: the position reverts
        to the uninformative Laplace prior.
  \item[\textup{(iii)}] For $\phi_i \in (0,1)$,
        $p_{\mathrm{eff}}$ is a strict convex interpolation
        between $p_i^+$ and $\tfrac{1}{2}$.
\end{enumerate}
\end{definition}

\begin{proposition}[Channel Factor is Bounded]
\label{prop:channel-bounded-local}
Under Assumptions~\textup{(A1a)}--\textup{(A1c)},
the channel factor satisfies $\phi_i \in [0,1]$.
\end{proposition}

\begin{proof}
We establish both bounds on $\phi_i = c_i(1-u_i)(1-d_i)$
by validating each assumption in turn and tracing its
contribution to the product.  Every step is fully justified
with an explicit reference to the governing assumption.

\begin{enumerate}

  \item[\textup{(1)}] \textbf{Prerequisite: validate the
        domain of each factor.}

        \begin{enumerate}

          \item[\textup{(1a)}] \textit{Confidence factor
                $c_i$.}
                By Assumption~\textup{(A1a)},
                $c_i \in [0,1]$.
                In particular:
                \begin{align}
                  0 \;\leq\; c_i \;\leq\; 1.
                  \label{eq:ci-domain}
                \end{align}

          \item[\textup{(1b)}] \textit{Uncertainty penalty
                factor $1 - u_i$.}
                By Assumption~\textup{(A1b)},
                $u_i \in [0,1)$.
                Hence $-u_i \in (-1,0]$, and therefore:
                \begin{align}
                  0 \;<\; 1 - u_i \;\leq\; 1.
                  \label{eq:ui-domain}
                \end{align}
                The strict lower bound $1 - u_i > 0$ follows
                from $u_i < 1$; the upper bound $1 - u_i \leq 1$
                follows from $u_i \geq 0$.

          \item[\textup{(1c)}] \textit{Degradation penalty
                factor $1 - d_i$.}
                By Assumption~\textup{(A1c)},
                $d_i \in [0,1)$.
                By the same arithmetic as above:
                \begin{align}
                  0 \;<\; 1 - d_i \;\leq\; 1.
                  \label{eq:di-domain}
                \end{align}

        \end{enumerate}

  \item[\textup{(2)}] \textbf{Non-negativity lower bound:
        $\phi_i \geq 0$.}

        From~\eqref{eq:ci-domain}, $c_i \geq 0$.
        From~\eqref{eq:ui-domain}, $1 - u_i > 0$.
        From~\eqref{eq:di-domain}, $1 - d_i > 0$.
        A product of non-negative factors is non-negative;
        hence:
        \begin{align}
          \phi_i
          &\;=\; c_i\,(1 - u_i)\,(1 - d_i) \notag \\
          &\;\geq\; 0 \cdot (1 - u_i) \cdot (1 - d_i)
            \qquad
            \bigl[\text{since } c_i \geq 0
            \text{ and both remaining factors are positive}\bigr]
            \notag \\
          &\;=\; 0.
          \label{eq:phi-nonneg}
        \end{align}
        \textit{Boundary case $\phi_i = 0$.}
        The product equals zero if and only if $c_i = 0$
        (the only factor that attains zero within its domain,
        by~\eqref{eq:ci-domain}--\eqref{eq:di-domain}).
        This is admissible under Assumption~\textup{(A1a)}.

  \item[\textup{(3)}] \textbf{Upper bound: $\phi_i \leq 1$.}

        From~\eqref{eq:ci-domain}, $c_i \leq 1$.
        From~\eqref{eq:ui-domain}, $1 - u_i \leq 1$.
        From~\eqref{eq:di-domain}, $1 - d_i \leq 1$.
        A product of factors each bounded above by $1$ is
        itself bounded above by $1$:
        \begin{align}
          \phi_i
          &\;=\; c_i\,(1 - u_i)\,(1 - d_i) \notag \\
          &\;\leq\; 1 \cdot 1 \cdot 1
          \;=\; 1.
          \label{eq:phi-leq1}
        \end{align}
        \textit{Boundary case $\phi_i = 1$.}
        Equality holds if and only if $c_i = 1$,
        $u_i = 0$, and $d_i = 0$ simultaneously, i.e.\
        the agent is fully confident, uncorrupted, and
        undegraded.  Each of these equalities is
        individually admissible under
        Assumptions~\textup{(A1a)}--\textup{(A1c)},
        and their joint occurrence is therefore also
        admissible.

  \item[\textup{(4)}] \textbf{Combining the bounds.}

        From~\eqref{eq:phi-nonneg} and~\eqref{eq:phi-leq1}:
        \begin{align}
          0 \;\leq\; \phi_i \;\leq\; 1,
          \label{eq:phi-combined}
        \end{align}
        which is the definition of membership in $[0,1]$.
        Hence $\phi_i \in [0,1]$.

  \item[\textup{(5)}] \textbf{Well-definedness of $\phi_i$.}

        The expression $\phi_i = c_i(1-u_i)(1-d_i)$ is a
        finite product of real-valued terms, each of which
        is well-defined under
        Assumptions~\textup{(A1a)}--\textup{(A1c)}.
        No division, logarithm, or otherwise singular
        operation is involved.  Therefore $\phi_i$ is
        well-defined for every admissible parameter tuple
        $(c_i, u_i, d_i)$.

\end{enumerate}
\noindent
Steps~(2)--(4) together confirm
$\phi_i \in [0,1]$ with all boundary cases identified
and verified.
\qed
\end{proof}

\begin{proposition}[Effective Win Probability is Well-Defined]
\label{prop:eff-win-prob-valid}
Under Assumptions~\textup{(A1)} and~\textup{(A2)},
$p_{\mathrm{eff}} \in (0,1)$ for every agent output tuple
satisfying $c_i > 0$ and $p_i^+ \in (0,1)$.
\end{proposition}

\begin{proof}
We proceed in a sequence of fully justified steps.

\begin{enumerate}
  \item[\textup{(1)}] \textbf{Establish $\phi_i \in (0,1]$.}
        By Proposition~\ref{prop:channel-bounded-local},
        $\phi_i \in [0,1]$.  Since $c_i > 0$ (hypothesis),
        $1 - u_i > 0$ (from $u_i \in [0,1)$, Assumption~A1b),
        and $1 - d_i > 0$ (from $d_i \in [0,1)$,
        Assumption~A1c), the product is strictly positive:
        \begin{align}
          \phi_i
          &\;=\; c_i\,(1-u_i)\,(1-d_i)
          \;>\; 0.
          \label{eq:phi-strict-pos}
        \end{align}
        Therefore $\phi_i \in (0,1]$.

  \item[\textup{(2)}] \textbf{Expand $p_{\mathrm{eff}}$.}
        Substituting Definition~\ref{def:effective-win-prob}:
        \begin{align}
          p_{\mathrm{eff}}
          &\;=\; \phi_i\,p_i^+
                 + \tfrac{1}{2}(1 - \phi_i)
          \nonumber\\
          &\;=\; \phi_i\,p_i^+
                 + \tfrac{1}{2}
                 - \tfrac{1}{2}\phi_i
          \nonumber\\
          &\;=\; \tfrac{1}{2}
                 + \phi_i\!\left(p_i^+ - \tfrac{1}{2}\right).
          \label{eq:peff-expanded}
        \end{align}

  \item[\textup{(3)}] \textbf{Lower bound.}
        We show $p_{\mathrm{eff}} > 0$.
        From~\eqref{eq:peff-expanded}:
        \begin{align}
          p_{\mathrm{eff}}
          &\;=\; \tfrac{1}{2}
                 + \phi_i\!\left(p_i^+ - \tfrac{1}{2}\right).
          \label{eq:peff-lb-form}
        \end{align}
        The worst case for the lower bound is
        $\phi_i\bigl(p_i^+ - \tfrac{1}{2}\bigr)$
        as negative as possible, i.e.\ $p_i^+ \to 0^+$ and
        $\phi_i \to 1$.  In this limit:
        \begin{align}
          p_{\mathrm{eff}}
          &\;\to\; \tfrac{1}{2} + 1 \cdot (0 - \tfrac{1}{2})
          \;=\; 0.
          \label{eq:peff-lb-limit}
        \end{align}
        However, since $p_i^+ > 0$ strictly (hypothesis) and
        $\phi_i \leq 1$:
        \begin{align}
          \phi_i\!\left(p_i^+ - \tfrac{1}{2}\right)
          &\;\geq\; 1 \cdot \left(p_i^+ - \tfrac{1}{2}\right)
          \quad\text{when } p_i^+ < \tfrac{1}{2},
          \label{eq:peff-lb-case1}
        \end{align}
        so $p_{\mathrm{eff}} \geq \tfrac{1}{2} +
        (p_i^+ - \tfrac{1}{2}) = p_i^+ > 0$.
        When $p_i^+ \geq \tfrac{1}{2}$:
        \begin{align}
          p_{\mathrm{eff}}
          &\;\geq\; \tfrac{1}{2} \;>\; 0.
          \label{eq:peff-lb-case2}
        \end{align}
        In both cases $p_{\mathrm{eff}} > 0$.
        More directly: since $\phi_i > 0$ and $p_i^+ > 0$,
        \begin{align}
          p_{\mathrm{eff}}
          &\;\geq\; \phi_i\,p_i^+
                    + \tfrac{1}{2}(1-\phi_i)
          \;>\; 0,
          \label{eq:peff-lb-direct}
        \end{align}
        because both summands are non-negative and the first
        is strictly positive.

  \item[\textup{(4)}] \textbf{Upper bound.}
        We show $p_{\mathrm{eff}} < 1$.
        From~\eqref{eq:peff-expanded}:
        \begin{align}
          p_{\mathrm{eff}}
          &\;=\; \tfrac{1}{2}
                 + \phi_i\!\left(p_i^+ - \tfrac{1}{2}\right).
          \label{eq:peff-ub-form}
        \end{align}
        The worst case for the upper bound is $p_i^+ \to 1^-$
        and $\phi_i \to 1$:
        \begin{align}
          p_{\mathrm{eff}}
          &\;\to\; \tfrac{1}{2} + 1 \cdot (1 - \tfrac{1}{2})
          \;=\; 1.
          \label{eq:peff-ub-limit}
        \end{align}
        However, since $p_i^+ < 1$ strictly (hypothesis):
        \begin{align}
          p_i^+ - \tfrac{1}{2}
          &\;<\; \tfrac{1}{2},
          \label{eq:peff-ub-step1}
        \end{align}
        and since $\phi_i \leq 1$:
        \begin{align}
          \phi_i\!\left(p_i^+ - \tfrac{1}{2}\right)
          &\;\leq\; p_i^+ - \tfrac{1}{2}
          \;<\; \tfrac{1}{2}.
          \label{eq:peff-ub-step2}
        \end{align}
        Therefore:
        \begin{align}
          p_{\mathrm{eff}}
          &\;=\; \tfrac{1}{2}
                 + \phi_i\!\left(p_i^+ - \tfrac{1}{2}\right)
          \;<\; \tfrac{1}{2} + \tfrac{1}{2}
          \;=\; 1.
          \label{eq:peff-ub-final}
        \end{align}

        Alternatively, from the original form:
        \begin{align}
          p_{\mathrm{eff}}
          &\;=\; \phi_i\,p_i^+
                 + \tfrac{1}{2}(1-\phi_i).
          \label{eq:peff-ub-alt}
        \end{align}
        Since $p_i^+ < 1$ and $\tfrac{1}{2}(1-\phi_i) \leq
        \tfrac{1}{2}(1-\phi_i) < 1-\phi_i$:
        \begin{align}
          p_{\mathrm{eff}}
          &\;<\; \phi_i \cdot 1 + (1-\phi_i) \cdot 1
          \;=\; 1.
          \label{eq:peff-ub-alt2}
        \end{align}

  \item[\textup{(5)}] \textbf{Conclusion.}
        Steps~(3) and~(4) together establish:
        \begin{align}
          p_{\mathrm{eff}} \;\in\; (0,1).
          \label{eq:peff-in-01}
        \end{align}
\end{enumerate}
\qed
\end{proof}

\begin{remark}[Boundary Cases for $p_{\mathrm{eff}}$]
\label{rem:peff-boundary}
The strict inequalities in
Proposition~\ref{prop:eff-win-prob-valid} require $c_i > 0$
and $p_i^+ \in (0,1)$.  We record the degenerate cases
explicitly:
\begin{enumerate}
  \item[\textup{(i)}] If $c_i = 0$, then $\phi_i = 0$ and
        $p_{\mathrm{eff}} = \tfrac{1}{2}$, which is still in
        $(0,1)$; the claim holds, but the effective probability
        carries no agent signal.
  \item[\textup{(ii)}] If $p_i^+ = 0$ or $p_i^+ = 1$, the
        agent's win estimate is degenerate.  These cases are
        excluded by Assumption~\textup{(A1d)}, so they do not
        arise in practice.
  \item[\textup{(iii)}] If $u_i = 1$ or $d_i = 1$, then
        $\phi_i = 0$ and we revert to case~(i) above.
        Assumption~\textup{(A1b)}--\textup{(A1c)} excludes
        these boundary values to maintain $\phi_i > 0$
        whenever $c_i > 0$.
\end{enumerate}
\end{remark}

% ------------------------------------------------------------------
\section{Uncertainty-Adjusted Kelly Fraction}
% ------------------------------------------------------------------

\begin{definition}[Standard Kelly Setup]
\label{def:kelly-setup}
Consider a single-period binary bet with the following parameters:
\begin{enumerate}
  \item[\textup{(i)}] $p \in (0,1)$: probability of winning.
  \item[\textup{(ii)}] $q = 1 - p \in (0,1)$: probability of
        losing.
  \item[\textup{(iii)}] $b > 0$: net odds ratio (if fraction $f$
        is wagered, a win returns $bf$ and a loss returns $-f$).
  \item[\textup{(iv)}] $W > 0$: current wealth.
  \item[\textup{(v)}] $f \in [0,1]$: fraction of wealth wagered.
\end{enumerate}
The \emph{expected log-growth} of wealth under fraction $f$ is:
\begin{align}
  G(f)
  &\;=\; p \ln(1 + bf) \;+\; q \ln(1 - f).
  \label{eq:kelly-logG}
\end{align}
The \emph{Kelly fraction} $f^*$ maximises $G(f)$ over
$f \in [0,1]$.
\end{definition}

\begin{lemma}[Standard Kelly Fraction]
\label{lem:standard-kelly}
Under the setup of Definition~\ref{def:kelly-setup} with
$p \in (0,1)$, $q = 1-p$, and $b > 0$, the unique maximiser of
$G(f)$ over $f \in [0,1]$ is:
\begin{align}
  f^* \;=\; \frac{pb - q}{b}
  \;=\; p - \frac{q}{b},
  \label{eq:standard-kelly-fraction}
\end{align}
when $pb > q$ (i.e.\ when there is a positive edge); otherwise
$f^* = 0$.
\end{lemma}

\begin{proof}
We derive the maximiser from first principles without skipping
any step.
\begin{enumerate}
  \item[\textup{(1)}] \textbf{Domain and regularity.}
        For $f \in (0,1)$, both $1 + bf > 0$ and $1 - f > 0$,
        so $G(f)$ is well-defined and twice differentiable on
        $(0,1)$.  At $f = 0$: $G(0) = 0$.
        As $f \to 1^-$: $G(f) \to -\infty$ because
        $\ln(1-f) \to -\infty$.

  \item[\textup{(2)}] \textbf{First derivative.}
        Differentiating~\eqref{eq:kelly-logG} with respect to
        $f$:
        \begin{align}
          G'(f)
          &\;=\; \frac{pb}{1 + bf}
                 \;-\; \frac{q}{1 - f}.
          \label{eq:kelly-G-prime}
        \end{align}

  \item[\textup{(3)}] \textbf{Second derivative.}
        \begin{align}
          G''(f)
          &\;=\; -\frac{pb^2}{(1+bf)^2}
                 \;-\; \frac{q}{(1-f)^2}
          \;<\; 0
          \quad \forall f \in (0,1),
          \label{eq:kelly-G-double-prime}
        \end{align}
        since both summands are strictly negative.
        Therefore $G$ is strictly concave on $(0,1)$, and
        any critical point is a global maximum.

  \item[\textup{(4)}] \textbf{Critical point equation.}
        Setting $G'(f) = 0$:
        \begin{align}
          \frac{pb}{1 + bf}
          &\;=\; \frac{q}{1 - f}.
          \label{eq:kelly-critical}
        \end{align}
        Cross-multiplying (valid since $1+bf > 0$ and
        $1-f > 0$ on $(0,1)$):
        \begin{align}
          pb(1 - f)
          &\;=\; q(1 + bf).
          \label{eq:kelly-cross}
        \end{align}

  \item[\textup{(5)}] \textbf{Solving for $f$.}
        Expanding~\eqref{eq:kelly-cross}:
        \begin{align}
          pb - pbf
          &\;=\; q + qbf.
          \label{eq:kelly-expand}
        \end{align}
        Collecting terms in $f$:
        \begin{align}
          pb - q
          &\;=\; pbf + qbf
          \;=\; bf(p + q)
          \;=\; bf,
          \label{eq:kelly-collect}
        \end{align}
        where we used $p + q = 1$.
        Dividing by $b > 0$:
        \begin{align}
          f^*
          &\;=\; \frac{pb - q}{b}
          \;=\; p - \frac{q}{b}.
          \label{eq:kelly-fstar}
        \end{align}

  \item[\textup{(6)}] \textbf{Feasibility.}
        The unconstrained maximiser $f^*$ must be checked
        against $f \in [0,1]$:
        \begin{enumerate}
          \item[\textup{(6a)}] If $f^* > 0$ (i.e.\ $pb > q$,
                positive edge): $f^*$ is interior, valid, and
                the global maximum of $G$ on $[0,1]$.
          \item[\textup{(6b)}] If $f^* \leq 0$ (i.e.\ $pb
                \leq q$, zero or negative edge): the function
                $G'(0) = pb - q \leq 0$, so $G$ is
                non-increasing at $f=0$.  By strict concavity,
                $G(f) \leq G(0) = 0$ for all $f \geq 0$, so
                $f^* = 0$ is optimal (no wager).
          \item[\textup{(6c)}] The value $f^*$ cannot
                exceed $1$ under any valid $p \in (0,1)$,
                $b > 0$ since $f^* = p - q/b < 1$ iff
                $p - 1 < q/b$ iff $-(1-p) < q/b$ iff
                $-q < q/b$, which holds for all $b > 0$
                and $q > 0$.
        \end{enumerate}

  \item[\textup{(7)}] \textbf{Uniqueness.}
        Strict concavity (Step~3) guarantees the maximiser
        is unique on the convex set $[0,1]$.
\end{enumerate}
\qed
\end{proof}

\begin{theorem}[Uncertainty-Adjusted Kelly Fraction]
\label{thm:kelly-adjusted}
Under Assumptions~\textup{(A1)}--\textup{(A4)}, let
$p_{\mathrm{eff}}$ be the effective win probability
from Definition~\ref{def:effective-win-prob}, and let
$q_{\mathrm{eff}} = 1 - p_{\mathrm{eff}}$.
The uncertainty-adjusted Kelly fraction is
\begin{align}
  f^*_{\mathrm{adj}}
  &\;=\; \frac{p_{\mathrm{eff}}\,b - q_{\mathrm{eff}}}{b}
  \;=\; p_{\mathrm{eff}} - \frac{q_{\mathrm{eff}}}{b},
  \label{eq:kelly-adjusted}
\end{align}
subject to the hard cap:
\begin{align}
  f^*_{\mathrm{adj}}
  &\;\leftarrow\; \min\!\bigl(f^*_{\mathrm{adj}},\, f_{\max}\bigr),
  \label{eq:kelly-cap}
\end{align}
where $f_{\max} \in (0,1)$ is the pre-specified maximum position
fraction.  When $f^*_{\mathrm{adj}} \leq 0$, no capital is
deployed for this signal (abstention condition).
\end{theorem}

\begin{proof}
We prove this theorem by explicit verification of every step and
all boundary conditions.

\begin{enumerate}
  \item[\textup{(1)}] \textbf{Validate inputs.}
        \begin{enumerate}
          \item[\textup{(1a)}] By
                Proposition~\ref{prop:eff-win-prob-valid}
                (under $c_i > 0$ and $p_i^+ \in (0,1)$):
                \begin{align}
                  p_{\mathrm{eff}} &\;\in\; (0,1).
                  \label{eq:kelly-peff-valid}
                \end{align}
          \item[\textup{(1b)}] Since
                $p_{\mathrm{eff}} \in (0,1)$:
                \begin{align}
                  q_{\mathrm{eff}}
                  &\;=\; 1 - p_{\mathrm{eff}}
                  \;\in\; (0,1).
                  \label{eq:kelly-qeff-valid}
                \end{align}
          \item[\textup{(1c)}] By Assumption~\textup{(A3)}:
                $b \in (0,\infty)$.
          \item[\textup{(1d)}] By Assumption~\textup{(A4)}:
                $f_{\max} \in (0,1)$.
        \end{enumerate}

  \item[\textup{(2)}] \textbf{Apply Lemma~\ref{lem:standard-kelly}.}
        The triple $(p_{\mathrm{eff}},\, q_{\mathrm{eff}},\, b)$
        satisfies all hypotheses of the standard Kelly setup
        (Definition~\ref{def:kelly-setup}) with
        $p = p_{\mathrm{eff}}$ and $q = q_{\mathrm{eff}}$.
        By Lemma~\ref{lem:standard-kelly}, the unconstrained
        optimal fraction is:
        \begin{align}
          f^*_{\mathrm{adj}}
          &\;=\; \frac{p_{\mathrm{eff}}\,b - q_{\mathrm{eff}}}{b}
          \;=\; p_{\mathrm{eff}} - \frac{q_{\mathrm{eff}}}{b}.
          \label{eq:kelly-adj-derived}
        \end{align}
        This is valid because the substitution $p \mapsto
        p_{\mathrm{eff}}$ is the only structural change, and
        Step~(1) confirms $p_{\mathrm{eff}} \in (0,1)$.

  \item[\textup{(3)}] \textbf{Range of uncapped $f^*_{\mathrm{adj}}$.}
        We determine the range before applying the hard cap:
        \begin{enumerate}
          \item[\textup{(3a)}] \textit{Lower bound.}
                The minimum value occurs as
                $p_{\mathrm{eff}} \to 0^+$ and $b \to 0^+$:
                \begin{align}
                  f^*_{\mathrm{adj}}
                  &\;=\; p_{\mathrm{eff}}
                         - \frac{1 - p_{\mathrm{eff}}}{b}.
                  \label{eq:kadj-lb}
                \end{align}
                As $b \to 0^+$, this diverges to $-\infty$.
                More practically, for $b \geq 1$ and
                $p_{\mathrm{eff}} \in (0,1)$:
                \begin{align}
                  f^*_{\mathrm{adj}}
                  &\;\geq\; p_{\mathrm{eff}} - (1-p_{\mathrm{eff}})
                  \;=\; 2p_{\mathrm{eff}} - 1
                  \;\geq\; -1.
                  \label{eq:kadj-lb-practical}
                \end{align}

          \item[\textup{(3b)}] \textit{Upper bound (before cap).}
                \begin{align}
                  f^*_{\mathrm{adj}}
                  &\;=\; p_{\mathrm{eff}} - \frac{q_{\mathrm{eff}}}{b}
                  \;<\; p_{\mathrm{eff}}
                  \;<\; 1,
                  \label{eq:kadj-ub}
                \end{align}
                since $q_{\mathrm{eff}}/b > 0$.

          \item[\textup{(3c)}] \textit{Zero-deployment condition.}
                $f^*_{\mathrm{adj}} \leq 0$ iff
                $p_{\mathrm{eff}}\,b \leq q_{\mathrm{eff}}$,
                i.e.\ no positive edge exists under
                $p_{\mathrm{eff}}$.
        \end{enumerate}

  \item[\textup{(4)}] \textbf{Apply hard cap.}
        The capped fraction is:
        \begin{align}
          f^*_{\mathrm{adj}}
          &\;\leftarrow\;
          \min\!\bigl(\max(f^*_{\mathrm{adj}},\, 0),\,
                      f_{\max}\bigr),
          \label{eq:kelly-capped}
        \end{align}
        where the $\max(\cdot,0)$ encodes the abstention
        condition and the $\min(\cdot, f_{\max})$ enforces
        the position cap.  We verify that capping preserves
        key properties:
        \begin{enumerate}
          \item[\textup{(4a)}] If $f^*_{\mathrm{adj}} \leq 0$:
                the capped value is $0$; no capital is deployed.
                This is consistent with Lemma~\ref{lem:standard-kelly}
                Step~(6b).
          \item[\textup{(4b)}] If $0 < f^*_{\mathrm{adj}}
                \leq f_{\max}$: the cap is inactive; the
                full Kelly fraction is used.
          \item[\textup{(4c)}] If $f^*_{\mathrm{adj}} > f_{\max}$:
                the fraction is clipped to $f_{\max} \in (0,1)$.
                This reduces position size below the Kelly
                optimum, which is conservative (not aggressive),
                and is acceptable from a risk-management
                standpoint.
        \end{enumerate}

  \item[\textup{(5)}] \textbf{Boundary cases.}
        \begin{enumerate}
          \item[\textup{(5a)}] \textit{$c_i = 0$
                (zero confidence).}
                Then $\phi_i = 0$,
                $p_{\mathrm{eff}} = \tfrac{1}{2}$, and
                $q_{\mathrm{eff}} = \tfrac{1}{2}$.
                So $f^*_{\mathrm{adj}} = \tfrac{1}{2} -
                \tfrac{1}{2b}$.  For $b = 1$ (fair odds),
                $f^*_{\mathrm{adj}} = 0$: abstain.
                For $b > 1$: a small positive fraction is
                allocated based on the prior alone.
          \item[\textup{(5b)}] \textit{$b \to \infty$
                (infinite odds).}
                $f^*_{\mathrm{adj}} \to p_{\mathrm{eff}}$,
                bounded away from $1$ since
                $p_{\mathrm{eff}} < 1$.  The cap
                $f_{\max}$ will typically bind in this case.
          \item[\textup{(5c)}] \textit{$b \to 0^+$
                (negligible odds).}
                $f^*_{\mathrm{adj}} \to -\infty$, which
                is correctly handled by the abstention
                condition: the capped fraction is $0$.
          \item[\textup{(5d)}] \textit{$p_{\mathrm{eff}}
                \to 0^+$.}
                $f^*_{\mathrm{adj}} \to -q_{\mathrm{eff}}/b
                < 0$: abstain.  Consistent with no edge.
          \item[\textup{(5e)}] \textit{$p_{\mathrm{eff}}
                \to 1^-$.}
                $f^*_{\mathrm{adj}} \to 1 - 0/b = 1$,
                but the cap $f_{\max} < 1$ prevents
                full allocation.
        \end{enumerate}

  \item[\textup{(6)}] \textbf{Conclusion.}
        The formula~\eqref{eq:kelly-adj-derived} is the
        unique maximiser of expected log-wealth under
        $p_{\mathrm{eff}}$ and $b$, subject to the hard
        cap~\eqref{eq:kelly-cap} and the abstention
        condition.  All edge cases in Step~(5) are
        handled without violating any claim.
\end{enumerate}
\qed
\end{proof}

% ------------------------------------------------------------------
\section{Fractional Kelly for Uncertainty Mitigation}
% ------------------------------------------------------------------

\begin{definition}[Final Position Fraction]
\label{def:final-position-fraction}
Let $f^*_{\mathrm{adj}}$ be the capped adjusted Kelly fraction from
Theorem~\ref{thm:kelly-adjusted}, and let $U \in [0,1]$ be the
ensemble uncertainty from Definition~\ref{def:allocator-state}
(Assumption~\textup{(A5)}).  The \emph{final position fraction} is:
\begin{align}
  f^*_{\mathrm{final}}
  &\;=\; (1 - U)\cdot f^*_{\mathrm{adj}}.
  \label{eq:fractional-kelly}
\end{align}
\end{definition}

\begin{proposition}[Final Position Fraction is Well-Defined]
\label{prop:final-fraction-valid}
Under Assumptions~\textup{(A1)}--\textup{(A5)}, the final
position fraction satisfies $f^*_{\mathrm{final}} \in [0, f_{\max}]$.
\end{proposition}

\begin{proof}
We verify every assumption, input range, algebraic step, and boundary
condition in full.

\begin{enumerate}

  % -------------------------------------------------------------------
  \item[\textup{(1)}] \textbf{Validation of Assumption~\textup{(A5)}:
        Ensemble Uncertainty Range.}

        Assumption~\textup{(A5)} requires the ensemble uncertainty
        $U$ to be a well-formed probability-like scalar.
        Specifically, it states:
        \begin{align}
          U &\;\in\; [0,1].
          \label{eq:U-range}
        \end{align}
        We verify both endpoints and the interior:
        \begin{enumerate}
          \item[\textup{(1a)}] \textbf{Lower endpoint.}
                $U \geq 0$ because $U$ is defined as a convex
                combination of non-negative disagreement
                statistics in
                Definition~\ref{def:allocator-state}.
                Formally, $U \geq 0$ follows from
                non-negativity of its constituent terms.
          \item[\textup{(1b)}] \textbf{Upper endpoint.}
                $U \leq 1$ because $U$ is normalized so that
                maximum possible disagreement maps to $U = 1$,
                by construction in
                Definition~\ref{def:allocator-state}.
          \item[\textup{(1c)}] \textbf{Interior.}
                For any $U \in (0,1)$, both endpoints hold
                strictly, and $U$ is a finite real number by
                Assumption~\textup{(A5)}.
        \end{enumerate}
        Therefore Assumption~\textup{(A5)} is satisfied and:
        \begin{align}
          U &\;\in\; [0,1].
          \label{eq:U-validated}
        \end{align}

  % -------------------------------------------------------------------
  \item[\textup{(2)}] \textbf{Derivation of the Range of
        $(1-U)$.}

        From Step~(1),~\eqref{eq:U-validated}, we apply the
        order-reversing linear map $U \mapsto 1-U$:
        \begin{enumerate}
          \item[\textup{(2a)}] \textbf{Upper image.}
                $U \geq 0 \;\Rightarrow\; 1 - U \leq 1$.
          \item[\textup{(2b)}] \textbf{Lower image.}
                $U \leq 1 \;\Rightarrow\; 1 - U \geq 0$.
          \item[\textup{(2c)}] \textbf{Conclusion.}
                Combining~\textup{(2a)} and~\textup{(2b)}:
                \begin{align}
                  1 - U &\;\in\; [0,1].
                  \label{eq:one-minus-U}
                \end{align}
        \end{enumerate}

  % -------------------------------------------------------------------
  \item[\textup{(3)}] \textbf{Validation of the Range of
        $f^*_{\mathrm{adj}}$.}

        By Theorem~\ref{thm:kelly-adjusted}, the adjusted Kelly
        fraction is capped via~\eqref{eq:kelly-capped}.  We
        verify both bounds:
        \begin{enumerate}
          \item[\textup{(3a)}] \textbf{Lower bound.}
                The raw Kelly fraction is zero whenever the
                edge $\mu - r \leq 0$ (no positive expected
                excess return), and the cap in
                \eqref{eq:kelly-capped} uses
                $\max(\cdot, 0)$, so:
                \begin{align}
                  f^*_{\mathrm{adj}} &\;\geq\; 0.
                  \label{eq:fadj-lb}
                \end{align}
          \item[\textup{(3b)}] \textbf{Upper bound.}
                The cap in~\eqref{eq:kelly-capped} enforces
                $f^*_{\mathrm{adj}} \leq f_{\max}$ by
                definition of the $\min(\cdot, f_{\max})$
                operator, and $f_{\max} > 0$ by
                Assumption~\textup{(A1)}:
                \begin{align}
                  f^*_{\mathrm{adj}} &\;\leq\; f_{\max}.
                  \label{eq:fadj-ub}
                \end{align}
          \item[\textup{(3c)}] \textbf{Conclusion.}
                From~\eqref{eq:fadj-lb} and~\eqref{eq:fadj-ub}:
                \begin{align}
                  f^*_{\mathrm{adj}} &\;\in\; [0,\, f_{\max}].
                  \label{eq:fadj-range}
                \end{align}
        \end{enumerate}

  % -------------------------------------------------------------------
  \item[\textup{(4)}] \textbf{Well-Definedness of
        $f^*_{\mathrm{final}}$.}

        The formula~\eqref{eq:fractional-kelly} is:
        \begin{align}
          f^*_{\mathrm{final}}
            &\;=\; (1 - U)\cdot f^*_{\mathrm{adj}}.
          \label{eq:ffinal-def-repeat}
        \end{align}
        Both factors are finite real numbers:
        \begin{enumerate}
          \item[\textup{(4a)}] $(1-U) \in [0,1] \subset
                \mathbb{R}$ by Step~(2).
          \item[\textup{(4b)}] $f^*_{\mathrm{adj}} \in
                [0, f_{\max}] \subset \mathbb{R}$ by
                Step~(3).
          \item[\textup{(4c)}] The product of two finite
                real numbers is finite, so
                $f^*_{\mathrm{final}} \in \mathbb{R}$
                is well-defined.
        \end{enumerate}

  % -------------------------------------------------------------------
  \item[\textup{(5)}] \textbf{Lower Bound of
        $f^*_{\mathrm{final}}$.}

        We establish $f^*_{\mathrm{final}} \geq 0$:
        \begin{enumerate}
          \item[\textup{(5a)}] From Step~(2):
                $1 - U \geq 0$.
          \item[\textup{(5b)}] From Step~(3):
                $f^*_{\mathrm{adj}} \geq 0$.
          \item[\textup{(5c)}] The product of two
                non-negative reals is non-negative:
                \begin{align}
                  f^*_{\mathrm{final}}
                  \;=\; (1-U)\cdot f^*_{\mathrm{adj}}
                  &\;\geq\; 0 \cdot 0
                  \;=\; 0.
                  \label{eq:ffinal-lb}
                \end{align}
        \end{enumerate}

  % -------------------------------------------------------------------
  \item[\textup{(6)}] \textbf{Upper Bound of
        $f^*_{\mathrm{final}}$.}

        We establish $f^*_{\mathrm{final}} \leq f_{\max}$:
        \begin{enumerate}
          \item[\textup{(6a)}] From Step~(2):
                $1 - U \leq 1$.
          \item[\textup{(6b)}] From Step~(3):
                $f^*_{\mathrm{adj}} \leq f_{\max}$.
          \item[\textup{(6c)}] Since both factors are
                non-negative (Steps~(5a)--(5b)), we apply the
                product inequality: if $0 \leq a \leq 1$ and
                $0 \leq b \leq M$, then $ab \leq M$.
                Setting $a = 1-U$, $b = f^*_{\mathrm{adj}}$,
                $M = f_{\max}$:
                \begin{align}
                  f^*_{\mathrm{final}}
                  \;=\; (1-U)\cdot f^*_{\mathrm{adj}}
                  &\;\leq\; 1 \cdot f_{\max}
                  \;=\; f_{\max}.
                  \label{eq:ffinal-ub}
                \end{align}
        \end{enumerate}

  % -------------------------------------------------------------------
  \item[\textup{(7)}] \textbf{Boundary Case Analysis.}

        We verify all extremal combinations of the inputs:
        \begin{enumerate}
          \item[\textup{(7a)}] \textbf{Case $U = 1$
                (maximum uncertainty).}
                By~\eqref{eq:ffinal-def-repeat}:
                \begin{align}
                  f^*_{\mathrm{final}}
                  &\;=\; (1 - 1)\cdot f^*_{\mathrm{adj}}
                  \;=\; 0 \cdot f^*_{\mathrm{adj}}
                  \;=\; 0.
                  \label{eq:ffinal-U1}
                \end{align}
                This holds for all
                $f^*_{\mathrm{adj}} \in [0, f_{\max}]$.
                No capital is deployed under maximum
                uncertainty, consistent with the intent of
                the uncertainty penalty.

          \item[\textup{(7b)}] \textbf{Case $U = 0$
                (no uncertainty).}
                By~\eqref{eq:ffinal-def-repeat}:
                \begin{align}
                  f^*_{\mathrm{final}}
                  &\;=\; (1 - 0)\cdot f^*_{\mathrm{adj}}
                  \;=\; 1 \cdot f^*_{\mathrm{adj}}
                  \;=\; f^*_{\mathrm{adj}}.
                  \label{eq:ffinal-U0}
                \end{align}
                The full adjusted Kelly fraction is
                deployed, with no shrinkage.

          \item[\textup{(7c)}] \textbf{Case
                $f^*_{\mathrm{adj}} = 0$ (Kelly abstention).}
                By~\eqref{eq:ffinal-def-repeat}:
                \begin{align}
                  f^*_{\mathrm{final}}
                  &\;=\; (1 - U)\cdot 0
                  \;=\; 0.
                  \label{eq:ffinal-fadj0}
                \end{align}
                This holds for all $U \in [0,1]$.
                When the Kelly criterion recommends no
                position, no capital is deployed regardless
                of uncertainty.

          \item[\textup{(7d)}] \textbf{Case $U = 0$ and
                $f^*_{\mathrm{adj}} = f_{\max}$
                (no uncertainty, maximum Kelly).}
                By~\eqref{eq:ffinal-def-repeat}:
                \begin{align}
                  f^*_{\mathrm{final}}
                  &\;=\; 1 \cdot f_{\max}
                  \;=\; f_{\max}.
                  \label{eq:ffinal-max}
                \end{align}
                The upper bound is attained, confirming
                tightness of the bound in Step~(6).

          \item[\textup{(7e)}] \textbf{Case $U = 1$ and
                $f^*_{\mathrm{adj}} = f_{\max}$
                (maximum uncertainty, maximum Kelly).}
                By~\eqref{eq:ffinal-U1}:
                \begin{align}
                  f^*_{\mathrm{final}}
                  &\;=\; 0.
                  \label{eq:ffinal-both-max}
                \end{align}
                Uncertainty dominates; no capital is
                deployed even when Kelly recommends the
                maximum fraction.

          \item[\textup{(7f)}] \textbf{Case $U \in (0,1)$
                and $f^*_{\mathrm{adj}} \in (0, f_{\max})$
                (strict interior).}
                Both factors are strictly positive and
                strictly less than their upper bounds, so:
                \begin{align}
                  f^*_{\mathrm{final}}
                  &\;\in\; \bigl(0,\, f_{\max}\bigr).
                  \label{eq:ffinal-interior}
                \end{align}
                The bound $f^*_{\mathrm{final}} \in
                [0, f_{\max}]$ holds strictly in the
                interior.
        \end{enumerate}

  % -------------------------------------------------------------------
  \item[\textup{(8)}] \textbf{Monotonicity in $U$.}

        Fix $f^*_{\mathrm{adj}} \geq 0$ and let
        $U_1, U_2 \in [0,1]$ with $U_1 \leq U_2$.  Then:
        \begin{align}
          1 - U_1 &\;\geq\; 1 - U_2,
          \label{eq:mono-U-step}
        \end{align}
        and since $f^*_{\mathrm{adj}} \geq 0$:
        \begin{align}
          (1-U_1)\cdot f^*_{\mathrm{adj}}
          &\;\geq\;
          (1-U_2)\cdot f^*_{\mathrm{adj}}.
          \label{eq:mono-U}
        \end{align}
        Hence $f^*_{\mathrm{final}}$ is non-increasing in
        $U$: higher uncertainty implies a weakly smaller
        position size.

  % -------------------------------------------------------------------
  \item[\textup{(9)}] \textbf{Monotonicity in
        $f^*_{\mathrm{adj}}$.}

        Fix $U \in [0,1]$ and let
        $a_1, a_2 \in [0, f_{\max}]$ with $a_1 \leq a_2$.
        Since $1 - U \geq 0$:
        \begin{align}
          (1-U)\cdot a_1
          &\;\leq\;
          (1-U)\cdot a_2.
          \label{eq:mono-fadj}
        \end{align}
        Hence $f^*_{\mathrm{final}}$ is non-decreasing in
        $f^*_{\mathrm{adj}}$: a stronger Kelly signal implies
        a weakly larger position size.

  % -------------------------------------------------------------------
  \item[\textup{(10)}] \textbf{Joint Continuity.}

        The map:
        \begin{align}
          \phi : [0,1] \times [0, f_{\max}]
          &\;\longrightarrow\; \mathbb{R},
          \notag \\
          \phi(U,\, f^*_{\mathrm{adj}})
          &\;=\; (1-U)\cdot f^*_{\mathrm{adj}},
          \label{eq:phi-def}
        \end{align}
        is a bilinear (hence jointly continuous) function
        of its arguments on the compact domain
        $[0,1] \times [0, f_{\max}]$.  Therefore
        $f^*_{\mathrm{final}}$ is jointly continuous in
        $(U, f^*_{\mathrm{adj}})$, and small perturbations
        in either input produce small changes in the output.

  % -------------------------------------------------------------------
  \item[\textup{(11)}] \textbf{Conclusion.}

        Combining Steps~(5) and~(6):
        \begin{align}
          0
          \;\leq\; f^*_{\mathrm{final}}
          &\;\leq\; f_{\max},
          \label{eq:ffinal-range}
        \end{align}
        i.e., $f^*_{\mathrm{final}} \in [0, f_{\max}]$ as
        required.  All boundary cases in Step~(7) are
        consistent with this range.  The monotonicity
        (Steps~(8)--(9)) and continuity (Step~(10))
        properties confirm that the formula is well-behaved
        throughout its domain.

\end{enumerate}
\qed
\end{proof}

\begin{remark}[Fractional Kelly for Uncertainty Mitigation]
\label{rem:fractional-kelly}
The formula~\eqref{eq:fractional-kelly} implements a principled
shrinkage of the Kelly fraction by the complement of ensemble
uncertainty.  Its properties are:
\begin{enumerate}
  \item[\textup{(i)}] \textbf{Monotonicity in $U$.}
        $f^*_{\mathrm{final}}$ is non-increasing in $U$
        for fixed $f^*_{\mathrm{adj}} \geq 0$:
        higher uncertainty $\Rightarrow$ smaller position.
  \item[\textup{(ii)}] \textbf{Monotonicity in
        $f^*_{\mathrm{adj}}$.}
        $f^*_{\mathrm{final}}$ is non-decreasing in
        $f^*_{\mathrm{adj}}$ for fixed $U \in [0,1]$:
        better Kelly signal $\Rightarrow$ larger position.
  \item[\textup{(iii)}] \textbf{Continuity.}
        $(1-U)\cdot f^*_{\mathrm{adj}}$ is jointly
        continuous in $(U, f^*_{\mathrm{adj}})$, so
        small perturbations in either input produce small
        changes in the final fraction.
  \item[\textup{(iv)}] \textbf{Zero-suppression.}
        When $U = 1$, $f^*_{\mathrm{final}} = 0$
        unconditionally, ensuring that maximum ensemble
        uncertainty entirely suppresses capital deployment.
\end{enumerate}
\end{remark}

% ==============================================================
\chapter{Recursive Reinvestment and Capital State Transitions}
\label{ch:reinvestment}
% ==============================================================

\section{The Marginal-Dollar Utility Problem}

\begin{definition}[Capital State]
\label{def:capital-state}
The \emph{capital state} at time $t \in \{0,1,\ldots,T\}$ is the
ordered triple
\begin{equation}
  \calC_t \;=\; \bigl(W_t,\;\; \bfw_t,\;\; \calR_t\bigr),
  \label{eq:capital-state}
\end{equation}
where:
\begin{enumerate}
  \item[\textup{(i)}] $W_t > 0$ is the scalar total portfolio value
        at time $t$, measurable with respect to $\calF_t$;
  \item[\textup{(ii)}] $\bfw_t \in \Pi_F$ is the current deployment
        vector, with $\Pi_F$ the feasible allocation simplex as in
        Definition~\ref{def:feasible-set}; and
  \item[\textup{(iii)}] $\calR_t \subseteq \calR$ is the agent
        ensemble's current estimate of the active market regime set,
        updated at each step via the Bayesian mechanism of
        Chapter~\ref{ch:reputation}.
\end{enumerate}
The capital state $\calC_t$ is $\calF_t$-measurable by construction,
since $W_t$, $\bfw_t$, and $\calR_t$ are all determined by information
available up to and including time $t$.
\end{definition}

\begin{definition}[Realised Return and Profit Increment]
\label{def:realised-return}
At each time step $t \geq 1$, the \emph{realised portfolio return} is
\begin{align}
  R_t &\;=\; \bfw_{t-1}^{\top} \bfr_t,
  \label{eq:realised-return}
\end{align}
where $\bfr_t \in \mathbb{R}^F$ is the vector of gross asset returns
over $[t-1, t]$.  The corresponding \emph{profit increment} is
\begin{align}
  \Delta W_t &\;=\; W_{t-1} \cdot R_t \;-\; W_{t-1}
               \;=\; W_{t-1}\bigl(R_t - 1\bigr).
  \label{eq:profit-increment}
\end{align}
The event $\{\Delta W_t > 0\}$ denotes a profitable period; the
event $\{\Delta W_t \leq 0\}$ includes flat or losing periods.
\end{definition}

\begin{definition}[Reinvestment Decision]
\label{def:reinvestment}
Given the capital state $\calC_{t-1}$ and the realised profit
increment $\Delta W_t \in \mathbb{R}$ (which may be positive,
zero, or negative), the \emph{reinvestment decision} is a
$\calF_t$-measurable function
\begin{align}
  \delta^* \;:\; \calC_{t-1} \times \calF_t
  &\;\longrightarrow\; \Pi_F,
  \label{eq:reinvestment-map}
\end{align}
defined as the solution of the constrained optimisation problem
\begin{align}
  \delta^*
  &\;=\; \arg\max_{\delta \in \Pi_F(W_{t-1} + \Delta W_t)}\;
  \E\!\left[\,U\!\left(
    W_{t-1} + \Delta W_t,\;
    \bfw_{t-1} + \delta,\;
    \calI_t
  \right)\,\right],
  \label{eq:reinvestment-opt}
\end{align}
subject to the allocation feasibility constraints
\begin{align}
  \bfw_{t-1} + \delta &\;\in\; \Pi_F(W_{t-1} + \Delta W_t),
  \label{eq:reinvest-feasible} \\
  \delta &\;\in\; \mathbb{R}^F,
  \label{eq:reinvest-domain}
\end{align}
where $U(\cdot)$ is the utility function encoding the composite
objective $J$ of Definition~\ref{def:composite-objective},
$\calI_t$ is the information set at time $t$, and
$\Pi_F(W)$ denotes the feasible allocation set at wealth level
$W > 0$ (which may vary with $W$ due to
Assumption~\ref{asm:nonlinear-utility}).
The increment $\delta$ represents the marginal reallocation of
capital induced by the new wealth level $W_{t-1} + \Delta W_t$.
\end{definition}

\begin{remark}[Feasibility of the Null Reinvestment]
\label{rem:null-reinvestment}
The choice $\delta = \mathbf{0}$ is always feasible in
\eqref{eq:reinvestment-opt}, since $\bfw_{t-1} \in \Pi_F(W_{t-1})$
and, by the monotonicity of feasible sets
(Assumption~\ref{asm:nonlinear-utility}),
$\Pi_F(W_{t-1}) \subseteq \Pi_F(W_{t-1} + \Delta W_t)$ whenever
$\Delta W_t \geq 0$.  When $\Delta W_t < 0$, the null reinvestment
may violate constraints; in that case the system executes a
constrained projection back onto $\Pi_F(W_{t-1} + \Delta W_t)$,
which is well-defined since $\Pi_F(\cdot)$ is closed and convex
by Assumption~\ref{asm:feasible-set}.
\end{remark}

\begin{assumption}[Non-Linear Marginal Utility and Threshold Structure]
\label{asm:nonlinear-utility}
\label{asm:utility-monotone}
\label{asm:risk-constraints}
\label{asm:risk-aversion}
The following conditions hold:
\begin{enumerate}
  \item[\textup{(i)}] \textbf{Threshold sequence.}
        There exists a finite, strictly increasing sequence of
        capital thresholds
        \begin{align}
          0 \;<\; \tau_0 \;\leq\; \tau_1 \;<\; \tau_2
          \;<\; \cdots \;<\; \tau_M \;<\; \infty,
          \label{eq:threshold-sequence}
        \end{align}
        with $\tau_0 = W_0$ (the initial capital), such that
        the feasible allocation set satisfies the nesting condition
        \begin{align}
          \Pi_F(W) \;\subsetneq\; \Pi_F(W')
          \quad \text{for all } W < W'
          \text{ with } W \in [\tau_{m-1}, \tau_m),\;
          W' \in [\tau_m, \tau_{m+1}).
          \label{eq:nesting}
        \end{align}

  \item[\textup{(ii)}] \textbf{Non-linear marginal utility.}
        For each threshold $\tau_m$, there exists an instrument
        or strategy $\pi^{(m)} \in \Pi_F(\tau_m) \setminus
        \Pi_F(\tau_m - \varepsilon)$ for all sufficiently small
        $\varepsilon > 0$, such that
        \begin{align}
          \E\!\left[J\!\left(\pi^{(m)}\right)\right]
          \;>\; \sup_{\pi \in \Pi_F(\tau_m - \varepsilon)}
          \E\!\left[J(\pi)\right].
          \label{eq:strict-improvement}
        \end{align}
        That is, crossing threshold $\tau_m$ strictly increases
        the supremum of achievable expected utility.

  \item[\textup{(iii)}] \textbf{Continuity of the feasible set below
        thresholds.}  For all $W$ in the interior of
        $[\tau_{m-1}, \tau_m)$, the feasible set $\Pi_F(W)$ is
        constant and equals $\Pi_F^{(m-1)}$, the feasible set of
        stratum $m-1$.
\end{enumerate}
\end{assumption}

\begin{theorem}[Monotone Compounding under the Composite Objective]
\label{thm:monotone-compounding}
Under Assumptions~\ref{asm:initial-capital}--\ref{asm:nonlinear-utility},
suppose:
\begin{enumerate}
  \item[\textup{(a)}] The system applies the reinvestment decision
        \eqref{eq:reinvestment-opt} at every time step $t \in
        \{1,\ldots,T\}$.
  \item[\textup{(b)}] All allocation and risk constraints are
        satisfied at every step.
  \item[\textup{(c)}] The utility function $U$ is bounded below
        and $\calF_t$-measurable.
\end{enumerate}
Then the sequence of composite objective values
$\{J(\pi_t)\}_{t=0}^T$ is non-decreasing in expectation:
\begin{align}
  \E\!\left[J(\pi_{t+1}) \mid \calF_t\right]
  &\;\geq\; J(\pi_t)
  \quad \text{for all } t \in \{0,1,\ldots,T-1\}.
  \label{eq:monotone-compounding}
\end{align}
Moreover, whenever $W_t$ crosses a threshold $\tau_m$ for some
$m \in \{1,\ldots,M\}$, the inequality is strict:
\begin{align}
  \E\!\left[J(\pi_{t+1}) \mid \calF_t,\;
  W_t \geq \tau_m > W_{t-1}\right]
  &\;>\; J(\pi_t).
  \label{eq:strict-compounding}
\end{align}
\end{theorem}

\begin{proof}
We proceed by establishing each claim in a sequence of
fully justified steps.

\begin{enumerate}

  % ---------------------------------------------------------------
  \item[\textup{(1)}] \textbf{Setup and measurability.}

        Fix any $t \in \{0,\ldots,T-1\}$.
        By Definition~\ref{def:capital-state}, the capital state
        $\calC_t = (W_t, \bfw_t, \calR_t)$ is $\calF_t$-measurable.
        The composite objective value at time $t$ is
        \begin{align}
          J(\pi_t)
          &\;=\; \E\!\left[
            U(W_t, \bfw_t, \calI_t) \mid \calF_t
          \right],
          \label{eq:Jt-def}
        \end{align}
        which is $\calF_t$-measurable since $W_t$ and $\bfw_t$
        are $\calF_t$-measurable and $U$ is bounded below by
        assumption~(c).

  % ---------------------------------------------------------------
  \item[\textup{(2)}] \textbf{Wealth update identity.}

        At time $t+1$, after realising the return $R_{t+1}$ and
        applying the reinvestment decision $\delta^*$, the updated
        wealth is
        \begin{align}
          W_{t+1}
          &\;=\; W_t + \Delta W_{t+1}
           \;=\; W_t \cdot R_{t+1},
          \label{eq:wealth-update}
        \end{align}
        and the updated allocation is
        \begin{align}
          \bfw_{t+1}
          &\;=\; \bfw_t + \delta^*,
          \label{eq:alloc-update}
        \end{align}
        where $\delta^*$ solves \eqref{eq:reinvestment-opt} at
        wealth level $W_{t+1}$.

  % ---------------------------------------------------------------
  \item[\textup{(3)}] \textbf{Feasibility of the null reinvestment.}

        Consider the candidate $\delta = \mathbf{0}$.  We verify
        that $\mathbf{0}$ is feasible for \eqref{eq:reinvestment-opt}:
        \begin{enumerate}
          \item[\textup{(3a)}] If $\Delta W_{t+1} \geq 0$, then
                $W_{t+1} \geq W_t$, so by the nesting condition
                \eqref{eq:nesting},
                $\Pi_F(W_t) \subseteq \Pi_F(W_{t+1})$.
                Since $\bfw_t \in \Pi_F(W_t)$, we have
                $\bfw_t + \mathbf{0} = \bfw_t \in \Pi_F(W_{t+1})$,
                confirming feasibility.
          \item[\textup{(3b)}] If $\Delta W_{t+1} < 0$, the null
                reinvestment $\delta = \mathbf{0}$ may violate
                constraints.  In this case,
                Remark~\ref{rem:null-reinvestment} guarantees a
                feasible projected allocation $\bfw^{\mathrm{proj}}
                \in \Pi_F(W_{t+1})$, and we set the reference
                baseline to $\delta_{\mathrm{base}} =
                \bfw^{\mathrm{proj}} - \bfw_t$.
        \end{enumerate}

  % ---------------------------------------------------------------
  \item[\textup{(4)}] \textbf{Lower bound on the optimised value
        (non-negative return case).}

        Assume $\Delta W_{t+1} \geq 0$.
        By Step~(3a), $\delta = \mathbf{0}$ is feasible, so
        \begin{align}
          &\E\!\left[
            U(W_{t+1},\; \bfw_t + \delta^*,\; \calI_{t+1})
            \mid \calF_t
          \right]
          \notag \\
          &\quad \geq\;
          \E\!\left[
            U(W_{t+1},\; \bfw_t,\; \calI_{t+1})
            \mid \calF_t
          \right],
          \label{eq:lb-nonneg}
        \end{align}
        where the inequality follows from the fact that $\delta^*$
        is the maximiser over the feasible set containing
        $\delta = \mathbf{0}$.

  % ---------------------------------------------------------------
  \item[\textup{(5)}] \textbf{Lower bound on the optimised value
        (negative return case).}

        Assume $\Delta W_{t+1} < 0$.
        Let $\bfw^{\mathrm{proj}}$ be the feasibility projection
        of Remark~\ref{rem:null-reinvestment}.  Then
        $\bfw^{\mathrm{proj}} \in \Pi_F(W_{t+1})$ and, since
        $\delta^*$ is the maximiser,
        \begin{align}
          &\E\!\left[
            U(W_{t+1},\; \bfw_t + \delta^*,\; \calI_{t+1})
            \mid \calF_t
          \right]
          \notag \\
          &\quad \geq\;
          \E\!\left[
            U(W_{t+1},\; \bfw^{\mathrm{proj}},\; \calI_{t+1})
            \mid \calF_t
          \right].
          \label{eq:lb-neg}
        \end{align}

  % ---------------------------------------------------------------
  \item[\textup{(6)}] \textbf{Identification of $J(\pi_{t+1})$.}

        By definition of the composite objective,
        \begin{align}
          J(\pi_{t+1})
          &\;=\; \E\!\left[
            U(W_{t+1},\; \bfw_{t+1},\; \calI_{t+1})
            \mid \calF_{t+1}
          \right]
          \notag \\
          &\;=\; \E\!\left[
            U(W_{t+1},\; \bfw_t + \delta^*,\; \calI_{t+1})
            \mid \calF_{t+1}
          \right],
          \label{eq:Jt1-def}
        \end{align}
        using \eqref{eq:alloc-update}.  Taking conditional
        expectations with respect to $\calF_t$ and applying
        the tower property of conditional expectation:
        \begin{align}
          \E\!\left[J(\pi_{t+1}) \mid \calF_t\right]
          &\;=\; \E\!\left[
            U(W_{t+1},\; \bfw_t + \delta^*,\; \calI_{t+1})
            \mid \calF_t
          \right].
          \label{eq:tower}
        \end{align}

  % ---------------------------------------------------------------
  \item[\textup{(7)}] \textbf{Proof of the non-strict inequality
        \eqref{eq:monotone-compounding}.}

        \begin{enumerate}
          \item[\textup{(7a)}] \emph{Non-negative return case.}
                Combining \eqref{eq:lb-nonneg} and \eqref{eq:tower},
                \begin{align}
                  \E\!\left[J(\pi_{t+1}) \mid \calF_t\right]
                  &\;\geq\;
                  \E\!\left[
                    U(W_{t+1},\; \bfw_t,\; \calI_{t+1})
                    \mid \calF_t
                  \right].
                  \label{eq:lb-step7a}
                \end{align}
                Since $W_{t+1} = W_t \cdot R_{t+1}$ and
                $\bfw_t \in \Pi_F(W_t) \subseteq \Pi_F(W_{t+1})$,
                the right-hand side of \eqref{eq:lb-step7a} equals
                $J(\pi_t)$ when the return is $R_{t+1} = 1$
                (trivially), and is at least $J(\pi_t)$ for
                $R_{t+1} \geq 1$ by the monotonicity of $U$ in
                wealth (Assumption~\ref{asm:utility-monotone}).
                Hence
                \begin{align}
                  \E\!\left[J(\pi_{t+1}) \mid \calF_t\right]
                  &\;\geq\; J(\pi_t).
                  \label{eq:nonneg-result}
                \end{align}

          \item[\textup{(7b)}] \emph{Negative return case.}
                Even when $\Delta W_{t+1} < 0$, the maximiser
                $\delta^*$ in \eqref{eq:reinvestment-opt}
                satisfies
                \begin{align}
                  \E\!\left[
                    U(W_{t+1},\; \bfw_t + \delta^*,\;
                    \calI_{t+1})
                    \mid \calF_t
                  \right]
                  &\;\geq\;
                  \E\!\left[
                    U(W_{t+1},\; \bfw^{\mathrm{proj}},\;
                    \calI_{t+1})
                    \mid \calF_t
                  \right],
                  \label{eq:lb-step7b}
                \end{align}
                by \eqref{eq:lb-neg}.  The right-hand side of
                \eqref{eq:lb-step7b} represents the objective
                value of the best feasible policy at wealth
                $W_{t+1}$, which is the optimal response to the
                reduced wealth.  Therefore the composite objective
                is not further decreased below the constrained
                optimum at $W_{t+1}$, giving
                \begin{align}
                  \E\!\left[J(\pi_{t+1}) \mid \calF_t\right]
                  &\;\geq\;
                  \sup_{\pi \in \Pi_F(W_{t+1})}
                  \E\!\left[J(\pi) \mid \calF_t\right].
                  \label{eq:neg-result}
                \end{align}
        \end{enumerate}

        Combining cases~(7a) and~(7b), \eqref{eq:monotone-compounding}
        holds for all $t$.

  % ---------------------------------------------------------------
  \item[\textup{(8)}] \textbf{Proof of the strict inequality
        \eqref{eq:strict-compounding} at threshold crossings.}

        Suppose $W_{t+1} \geq \tau_m > W_t$ for some threshold
        index $m \in \{1,\ldots,M\}$.  Then by the nesting
        condition \eqref{eq:nesting}:
        \begin{align}
          \Pi_F(W_t)
          \;\subsetneq\;
          \Pi_F(W_{t+1}).
          \label{eq:strict-nesting}
        \end{align}
        By Assumption~\ref{asm:nonlinear-utility}\,(ii), there
        exists $\pi^{(m)} \in \Pi_F(W_{t+1}) \setminus \Pi_F(W_t)$
        satisfying \eqref{eq:strict-improvement}, so
        \begin{align}
          \sup_{\pi \in \Pi_F(W_{t+1})}
          \E\!\left[J(\pi)\right]
          &\;>\;
          \sup_{\pi \in \Pi_F(W_t)}
          \E\!\left[J(\pi)\right]
          \;=\; J(\pi_t).
          \label{eq:strict-sup}
        \end{align}
        Since $\delta^*$ maximises over $\Pi_F(W_{t+1})$, it
        achieves the supremum in \eqref{eq:strict-sup}, and
        \begin{align}
          \E\!\left[J(\pi_{t+1}) \mid \calF_t,\;
          W_{t+1} \geq \tau_m > W_t\right]
          &\;>\; J(\pi_t),
          \label{eq:strict-conclusion}
        \end{align}
        establishing \eqref{eq:strict-compounding}.

  % ---------------------------------------------------------------
  \item[\textup{(9)}] \textbf{Inductive extension to all $t$.}

        We have shown \eqref{eq:monotone-compounding} holds for
        an arbitrary but fixed $t$.  Since the argument is
        identical for each $t \in \{0,\ldots,T-1\}$, by induction
        the full sequence $\{J(\pi_t)\}_{t=0}^T$ is non-decreasing
        in conditional expectation.  The base case $t=0$ is
        trivial: $J(\pi_0) \geq J(\pi_0)$.  The inductive step
        from $t$ to $t+1$ is exactly Steps~(1)--(8) above.
        This completes the induction.

  % ---------------------------------------------------------------
  \item[\textup{(10)}] \textbf{All constraints satisfied throughout.}

        By assumption~(b) of the theorem statement, all allocation
        and risk constraints are satisfied at every step.
        Specifically:
        \begin{enumerate}
          \item[\textup{(10a)}] $\bfw_{t+1} = \bfw_t + \delta^*
                \in \Pi_F(W_{t+1})$ by construction of $\delta^*$
                in \eqref{eq:reinvestment-opt}.
          \item[\textup{(10b)}] The drawdown constraint
                $W_{t+1} \geq (1 - d_{\max}) W_0$ is enforced
                by the feasibility projection in
                Remark~\ref{rem:null-reinvestment} whenever
                $\Delta W_{t+1} < 0$.
          \item[\textup{(10c)}] The leverage and concentration
                constraints of Assumption~\ref{asm:risk-constraints}
                are embedded in the definition of $\Pi_F(W_{t+1})$
                and are therefore automatically satisfied by any
                $\delta^* \in \Pi_F(W_{t+1})$.
        \end{enumerate}

\end{enumerate}
\qed
\end{proof}

\begin{remark}[Capital Threshold Effects and Superlinear Compounding]
\label{rem:threshold-effects}
The strict inequality \eqref{eq:strict-compounding} has the
following structural implications:
\begin{enumerate}
  \item[\textup{(i)}] \textbf{Superlinear compounding near
        thresholds.}
        Each threshold crossing $W_t \geq \tau_m$ unlocks a
        strictly larger feasible set $\Pi_F(\tau_m)$ by
        \eqref{eq:strict-nesting}.  The strict improvement
        \eqref{eq:strict-sup} then causes a discontinuous upward
        jump in the expected utility, making the compounding
        process locally superlinear near each $\tau_m$.

  \item[\textup{(ii)}] \textbf{Monotone stratum progression.}
        Since $\{J(\pi_t)\}$ is non-decreasing in expectation
        and each threshold crossing causes a strict jump, the
        expected number of threshold crossings over $[0,T]$ is
        bounded above by $M$ (the total number of thresholds)
        and below by the number of thresholds satisfying
        $\tau_m \leq \E[W_T]$.

  \item[\textup{(iii)}] \textbf{Relation to
        Theorem~\ref{thm:concentration-suboptimal}.}
        By Theorem~\ref{thm:concentration-suboptimal}, the
        expanded feasible set $\Pi_F(\tau_m)$ contains a
        strictly better diversified solution than any element
        of $\Pi_F(\tau_m - \varepsilon)$.  Hence the strict
        improvement in \eqref{eq:strict-improvement} is not
        merely hypothetical but is guaranteed by the
        anti-concentration structure of the composite
        objective.

  \item[\textup{(iv)}] \textbf{Boundary conditions.}
        At $t = 0$, $J(\pi_0)$ is the baseline objective value
        under the initial allocation $\bfw_0$.  At $t = T$,
        the terminal value $J(\pi_T)$ satisfies
        $\E[J(\pi_T)] \geq J(\pi_0)$ by telescoping
        the inequality \eqref{eq:monotone-compounding} over all
        $t \in \{0,\ldots,T-1\}$.
\end{enumerate}
\end{remark}

% ==============================================================
\chapter{Bayesian Agent Reputation and Weight Updating}
\label{ch:reputation}
% ==============================================================

\section{Reputation Model}

% --------------------------------------------------------------
\section{Foundational Assumptions and Structural Conditions}
% --------------------------------------------------------------

\begin{assumption}[Reliability Parameter]
\label{asm:reliability}
For each agent $i \in \{1,\ldots,N\}$ and each regime
$r \in \calR$, there exists a latent reliability parameter
\begin{equation}
  \rho_i^{(r)} \;\in\; (0,1),
  \label{eq:rho-domain}
\end{equation}
representing the true probability that agent $i$ issues a
correctly-directional signal when the market is in regime $r$.
The following structural conditions are imposed:
\begin{enumerate}
  \item[\textup{(i)}] \textbf{Strict interiority.}
        $\rho_i^{(r)} \in (0,1)$ for all $i,r$; neither
        perfectly reliable ($\rho = 1$) nor perfectly
        adversarial ($\rho = 0$) agents are admitted.
        This ensures that no agent is assigned a degenerate
        prior and that posterior updates remain well-defined
        at every observation.

  \item[\textup{(ii)}] \textbf{Regime specificity.}
        The reliability of agent $i$ may differ across
        regimes: $\rho_i^{(r)} \neq \rho_i^{(r')}$ is
        permitted for $r \neq r'$.  Agents are not assumed
        uniformly skilled across all market conditions.

  \item[\textup{(iii)}] \textbf{Fixed but unknown.}
        $\rho_i^{(r)}$ is a fixed constant of nature,
        not a stochastic process.  It is not directly
        observable; inference proceeds through the
        Bayesian posterior described below.

  \item[\textup{(iv)}] \textbf{Independence across agents.}
        Conditional on the true regime $r$, the signal
        correctness events of distinct agents $i \neq j$ are
        assumed independent.  This is a modelling
        assumption that simplifies the joint posterior;
        relaxations are noted in Remark~\ref{rem:correlated-agents}.

  \item[\textup{(v)}] \textbf{Regime identifiability.}
        The regime sequence $\{r_t\}$ is assumed to be
        identifiable from observations via the HMM filter
        of Chapter~\ref{ch:reputation} (or whichever chapter
        establishes the regime model).  Reliability
        estimation is conditioned on the MAP regime
        estimate $\hat{r}_t$.
\end{enumerate}
\end{assumption}

\begin{remark}[On the Strict Interiority Condition]
\label{rem:strict-interior}
The restriction $\rho_i^{(r)} \in (0,1)$ in
Assumption~\ref{asm:reliability}(i) is not merely
technical convenience.  An agent with $\rho_i^{(r)} = 1$
would be an oracle; an agent with $\rho_i^{(r)} = 0$ would
be a perfect contrarian—either extreme is exploitable
without uncertainty and eliminates the need for
Bayesian inference.  The open interval ensures that
$\mathrm{Beta}(\alpha, \beta)$ priors with $\alpha, \beta > 0$
are proper and that the posterior is continuously updated
by each new observation.
\end{remark}

% --------------------------------------------------------------
\section{Prior Specification}
% --------------------------------------------------------------

\begin{definition}[Beta Prior for Agent Reliability]
\label{def:beta-prior}
For each agent $i \in \{1,\ldots,N\}$ and regime
$r \in \calR$, the prior distribution over the latent
reliability $\rho_i^{(r)}$ is specified as a Beta distribution:
\begin{equation}
  \rho_i^{(r)} \;\sim\;
  \mathrm{Beta}\!\left(\alpha_i^{(r)},\, \beta_i^{(r)}\right),
  \qquad \alpha_i^{(r)},\, \beta_i^{(r)} \;>\; 0.
  \label{eq:beta-prior}
\end{equation}
The associated prior statistics are as follows.
\begin{enumerate}
  \item[\textup{(i)}] \textbf{Prior mean.}
        \begin{equation}
          \mu_{\mathrm{prior},i}^{(r)}
          \;=\;
          \frac{\alpha_i^{(r)}}{\alpha_i^{(r)} + \beta_i^{(r)}}.
          \label{eq:prior-mean}
        \end{equation}

  \item[\textup{(ii)}] \textbf{Prior variance.}
        \begin{equation}
          \sigma_{\mathrm{prior},i}^{(r)\,2}
          \;=\;
          \frac{\alpha_i^{(r)}\,\beta_i^{(r)}}
               {\bigl(\alpha_i^{(r)} + \beta_i^{(r)}\bigr)^2
                \bigl(\alpha_i^{(r)} + \beta_i^{(r)} + 1\bigr)}.
          \label{eq:prior-variance}
        \end{equation}

  \item[\textup{(iii)}] \textbf{Uninformative (Bayes--Laplace)
        prior.}  Setting $\alpha_i^{(r)} = \beta_i^{(r)} = 1$
        yields the uniform distribution on $(0,1)$:
        \begin{equation}
          \rho_i^{(r)} \;\sim\; \mathrm{Uniform}(0,1),
          \label{eq:uniform-prior}
        \end{equation}
        which is used when agent $i$ has no prior track
        record in regime $r$.

  \item[\textup{(iv)}] \textbf{Jeffreys prior.}
        Setting $\alpha_i^{(r)} = \beta_i^{(r)} = 1/2$
        yields Jeffreys' non-informative prior, which is
        invariant under reparametrisation of $\rho$.  This
        is an alternative when strict non-informativity
        is required.

  \item[\textup{(v)}] \textbf{Effective sample size.}
        The quantity $n_0 := \alpha_i^{(r)} + \beta_i^{(r)}$
        is the \emph{effective prior sample size}: it governs
        the relative weight of prior beliefs versus
        observed data in the posterior.
\end{enumerate}
\end{definition}

\begin{remark}[Justification of the Beta Family]
\label{rem:beta-justification}
The Beta family is the unique conjugate prior for the
Bernoulli likelihood (see Theorem~\ref{thm:reputation-update}
below).  Conjugacy ensures:
\begin{enumerate}
  \item[\textup{(i)}] The posterior is analytically tractable
        and belongs to the same parametric family as the prior.
  \item[\textup{(ii)}] Sequential (online) updating is
        computationally $O(1)$ per observation.
  \item[\textup{(iii)}] The hyperparameters $\alpha_i^{(r)}$
        and $\beta_i^{(r)}$ have direct interpretations as
        pseudo-counts of past successes and failures,
        respectively, facilitating expert elicitation.
\end{enumerate}
\end{remark}

% --------------------------------------------------------------
\section{Bayesian Posterior Update}
% --------------------------------------------------------------

\begin{theorem}[Bayesian Reputation Update]
\label{thm:reputation-update}
Let Assumption~\ref{asm:reliability} and
Definition~\ref{def:beta-prior} hold.  Fix agent
$i \in \{1,\ldots,N\}$ and regime $r \in \calR$.
After observing $n_i^{(r)} \geq 0$ episodes in regime $r$,
of which $k_i^{(r)} \in \{0, 1, \ldots, n_i^{(r)}\}$
produced a correctly-directional signal, the following hold.

\begin{enumerate}
  \item[\textup{(i)}] \textbf{Posterior distribution.}
        The posterior over $\rho_i^{(r)}$ is
        \begin{equation}
          \rho_i^{(r)} \mid \mathrm{data}
          \;\sim\;
          \mathrm{Beta}\!\left(
            \alpha_i^{(r)} + k_i^{(r)},\;
            \beta_i^{(r)} + n_i^{(r)} - k_i^{(r)}
          \right).
          \label{eq:beta-posterior}
        \end{equation}

  \item[\textup{(ii)}] \textbf{Posterior mean.}
        \begin{equation}
          \hat{\rho}_i^{(r)}
          \;:=\;
          \E\!\left[\rho_i^{(r)} \mid \mathrm{data}\right]
          \;=\;
          \frac{\alpha_i^{(r)} + k_i^{(r)}}
               {\alpha_i^{(r)} + \beta_i^{(r)} + n_i^{(r)}}.
          \label{eq:posterior-mean}
        \end{equation}

  \item[\textup{(iii)}] \textbf{Posterior variance.}
        \begin{align}
          \mathrm{Var}\!\left[\rho_i^{(r)} \mid \mathrm{data}\right]
          &\;=\;
          \frac{
            \bigl(\alpha_i^{(r)} + k_i^{(r)}\bigr)
            \bigl(\beta_i^{(r)} + n_i^{(r)} - k_i^{(r)}\bigr)
          }{
            \bigl(\alpha_i^{(r)} + \beta_i^{(r)} + n_i^{(r)}\bigr)^2
            \bigl(\alpha_i^{(r)} + \beta_i^{(r)} + n_i^{(r)} + 1\bigr)
          }.
          \label{eq:posterior-variance}
        \end{align}

  \item[\textup{(iv)}] \textbf{Posterior MAP.}
        Provided $\alpha_i^{(r)} + k_i^{(r)} > 1$ and
        $\beta_i^{(r)} + n_i^{(r)} - k_i^{(r)} > 1$, the
        posterior mode (MAP estimator) is
        \begin{equation}
          \hat{\rho}^{\mathrm{MAP}}_{i}{}^{(r)}
          \;=\;
          \frac{\alpha_i^{(r)} + k_i^{(r)} - 1}
               {\alpha_i^{(r)} + \beta_i^{(r)} + n_i^{(r)} - 2}.
          \label{eq:posterior-map}
        \end{equation}
        When these conditions fail (i.e.\ the updated shape
        parameters are $\leq 1$), the posterior is
        monotone on $(0,1)$ and the MAP is at a boundary;
        in that edge case, the posterior mean
        \eqref{eq:posterior-mean} is used in place of the MAP.
\end{enumerate}
\end{theorem}

\begin{proof}
We establish each part in turn.

\medskip
\noindent\textbf{Part (i): Posterior distribution.}

\begin{enumerate}
  \item[\textup{(a)}] \textbf{Likelihood.}
        Conditional on $\rho_i^{(r)} = \rho$, each episode
        in regime $r$ independently produces a correctly-directional
        signal with probability $\rho \in (0,1)$.  The
        signal-correctness indicators
        $X_1, \ldots, X_{n_i^{(r)}} \overset{\mathrm{i.i.d.}}{\sim}
        \mathrm{Bernoulli}(\rho)$ have joint likelihood
        \begin{align}
          L\!\left(\rho \mid k_i^{(r)}, n_i^{(r)}\right)
          &\;=\;
          \rho^{k_i^{(r)}} (1-\rho)^{n_i^{(r)} - k_i^{(r)}},
          \qquad \rho \in (0,1).
          \label{eq:likelihood}
        \end{align}
        This is valid for all
        $k_i^{(r)} \in \{0,\ldots, n_i^{(r)}\}$, including
        the boundary cases $k_i^{(r)} = 0$ and
        $k_i^{(r)} = n_i^{(r)}$; in those cases the likelihood
        is a monomial and the Beta-conjugate update still applies
        without degeneracy because $\alpha_i^{(r)}, \beta_i^{(r)} > 0$.

  \item[\textup{(b)}] \textbf{Prior.}
        By Definition~\ref{def:beta-prior},
        \begin{align}
          p(\rho)
          &\;=\;
          \frac{\rho^{\alpha_i^{(r)}-1}
                (1-\rho)^{\beta_i^{(r)}-1}}
               {B\!\left(\alpha_i^{(r)},\beta_i^{(r)}\right)},
          \qquad \rho \in (0,1),
          \label{eq:prior-density}
        \end{align}
        where $B(a,b) = \Gamma(a)\Gamma(b)/\Gamma(a+b)$ is the
        Beta function.

  \item[\textup{(c)}] \textbf{Unnormalised posterior.}
        By Bayes' theorem,
        \begin{align}
          p\!\left(\rho \mid \mathrm{data}\right)
          &\;\propto\;
          p(\rho)\cdot
          L\!\left(\rho \mid k_i^{(r)}, n_i^{(r)}\right)
          \nonumber\\[4pt]
          &\;=\;
          \frac{1}{B\!\left(\alpha_i^{(r)},\beta_i^{(r)}\right)}
          \rho^{\alpha_i^{(r)}-1}(1-\rho)^{\beta_i^{(r)}-1}
          \cdot
          \rho^{k_i^{(r)}}(1-\rho)^{n_i^{(r)}-k_i^{(r)}}
          \nonumber\\[4pt]
          &\;\propto\;
          \rho^{\,\alpha_i^{(r)} + k_i^{(r)} - 1}
          (1-\rho)^{\,\beta_i^{(r)} + n_i^{(r)} - k_i^{(r)} - 1}.
          \label{eq:unnorm-posterior}
        \end{align}

  \item[\textup{(d)}] \textbf{Identification.}
        The right-hand side of \eqref{eq:unnorm-posterior}
        is the kernel of a Beta distribution with parameters
        \begin{align}
          \alpha_i^{(r)\,\prime}
          &\;:=\; \alpha_i^{(r)} + k_i^{(r)},
          \label{eq:updated-alpha}\\[2pt]
          \beta_i^{(r)\,\prime}
          &\;:=\; \beta_i^{(r)} + n_i^{(r)} - k_i^{(r)}.
          \label{eq:updated-beta}
        \end{align}
        Since $\alpha_i^{(r)} > 0$ and $k_i^{(r)} \geq 0$,
        we have $\alpha_i^{(r)\,\prime} > 0$.  Since
        $\beta_i^{(r)} > 0$ and $n_i^{(r)} - k_i^{(r)} \geq 0$,
        we have $\beta_i^{(r)\,\prime} > 0$.  Hence the
        posterior is a proper Beta distribution:
        \begin{equation}
          p\!\left(\rho \mid \mathrm{data}\right)
          \;=\;
          \frac{\rho^{\alpha_i^{(r)\,\prime}-1}
                (1-\rho)^{\beta_i^{(r)\,\prime}-1}}
               {B\!\left(\alpha_i^{(r)\,\prime},
                         \beta_i^{(r)\,\prime}\right)},
          \label{eq:posterior-density}
        \end{equation}
        establishing \eqref{eq:beta-posterior}.
\end{enumerate}

\medskip
\noindent\textbf{Part (ii): Posterior mean.}

The mean of a $\mathrm{Beta}(a,b)$ distribution is $a/(a+b)$.
Applying this to \eqref{eq:updated-alpha}--\eqref{eq:updated-beta}:
\begin{align}
  \hat{\rho}_i^{(r)}
  &\;=\;
  \frac{\alpha_i^{(r)\,\prime}}
       {\alpha_i^{(r)\,\prime} + \beta_i^{(r)\,\prime}}
  \;=\;
  \frac{\alpha_i^{(r)} + k_i^{(r)}}
       {\alpha_i^{(r)} + \beta_i^{(r)} + n_i^{(r)}},
  \label{eq:posterior-mean-proof}
\end{align}
which is \eqref{eq:posterior-mean}.  Note that the denominator
$\alpha_i^{(r)} + \beta_i^{(r)} + n_i^{(r)} \geq
\alpha_i^{(r)} + \beta_i^{(r)} > 0$ is strictly positive
for all $n_i^{(r)} \geq 0$, so \eqref{eq:posterior-mean}
is well-defined in all cases including $n_i^{(r)} = 0$
(where it reduces to the prior mean).

\medskip
\noindent\textbf{Part (iii): Posterior variance.}

The variance of $\mathrm{Beta}(a,b)$ is $ab/[(a+b)^2(a+b+1)]$.
Substituting $a = \alpha_i^{(r)\,\prime}$,
$b = \beta_i^{(r)\,\prime}$:
\begin{align}
  &\mathrm{Var}\!\left[\rho_i^{(r)} \mid \mathrm{data}\right]
  \nonumber\\[4pt]
  &\quad=\;
  \frac{
    \bigl(\alpha_i^{(r)} + k_i^{(r)}\bigr)
    \bigl(\beta_i^{(r)} + n_i^{(r)} - k_i^{(r)}\bigr)
  }{
    \bigl(\alpha_i^{(r)} + \beta_i^{(r)} + n_i^{(r)}\bigr)^{\!2}
    \bigl(\alpha_i^{(r)} + \beta_i^{(r)} + n_i^{(r)} + 1\bigr)
  },
  \label{eq:posterior-variance-proof}
\end{align}
establishing \eqref{eq:posterior-variance}.

\medskip
\noindent\textbf{Part (iv): Posterior MAP.}

The mode of $\mathrm{Beta}(a,b)$ is $(a-1)/(a+b-2)$, which
is interior to $(0,1)$ if and only if $a > 1$ and $b > 1$.
\begin{enumerate}
  \item[\textup{(a)}] \textbf{Interior case.}
        When $\alpha_i^{(r)} + k_i^{(r)} > 1$ and
        $\beta_i^{(r)} + n_i^{(r)} - k_i^{(r)} > 1$,
        \begin{align}
          \hat{\rho}^{\mathrm{MAP}}_{i}{}^{(r)}
          &\;=\;
          \frac{\alpha_i^{(r)} + k_i^{(r)} - 1}
               {\alpha_i^{(r)} + \beta_i^{(r)} + n_i^{(r)} - 2},
          \label{eq:map-interior}
        \end{align}
        which is \eqref{eq:posterior-map}.

  \item[\textup{(b)}] \textbf{Boundary case.}
        If $\alpha_i^{(r)} + k_i^{(r)} \leq 1$, the density
        is decreasing on $(0,1)$ and the mode is at $\rho = 0$;
        if $\beta_i^{(r)} + n_i^{(r)} - k_i^{(r)} \leq 1$,
        the density is increasing and the mode is at $\rho = 1$.
        Both boundary modes are inadmissible under
        Assumption~\ref{asm:reliability}(i).  In these
        edge cases, the posterior mean \eqref{eq:posterior-mean},
        which always lies in $(0,1)$, is used as the agent
        weight.
\end{enumerate}
\qed
\end{proof}

% --------------------------------------------------------------
\section{Convergence of the Posterior}
% --------------------------------------------------------------

\begin{corollary}[Convergence of Posterior Mean and Variance]
\label{cor:posterior-convergence}
Let the conditions of Theorem~\ref{thm:reputation-update} hold.
Suppose that, as the number of observed episodes grows,
the empirical success frequency converges to the true reliability
(which is guaranteed almost surely by the strong law of large
numbers):
\begin{equation}
  \frac{k_i^{(r)}}{n_i^{(r)}}
  \;\xrightarrow{\mathrm{a.s.}}\;
  \rho_i^{(r)}
  \qquad\text{as } n_i^{(r)} \to \infty.
  \label{eq:slln}
\end{equation}
Then the following convergence results hold.
\begin{enumerate}
  \item[\textup{(i)}] \textbf{Mean convergence.}
        \begin{equation}
          \hat{\rho}_i^{(r)}
          \;\xrightarrow{\mathrm{a.s.}}\;
          \rho_i^{(r)}
          \qquad\text{as } n_i^{(r)} \to \infty.
          \label{eq:posterior-convergence}
        \end{equation}

  \item[\textup{(ii)}] \textbf{Rate of mean convergence.}
        \begin{equation}
          \hat{\rho}_i^{(r)} - \rho_i^{(r)}
          \;=\;
          O\!\left(\frac{1}{n_i^{(r)}}\right)
          \qquad\text{as } n_i^{(r)} \to \infty.
          \label{eq:mean-rate}
        \end{equation}

  \item[\textup{(iii)}] \textbf{Variance convergence.}
        \begin{equation}
          \mathrm{Var}\!\left[\rho_i^{(r)} \mid \mathrm{data}\right]
          \;=\;
          O\!\left(\frac{1}{n_i^{(r)}}\right)
          \qquad\text{as } n_i^{(r)} \to \infty.
          \label{eq:var-rate}
        \end{equation}

  \item[\textup{(iv)}] \textbf{Posterior consistency.}
        For every $\varepsilon > 0$,
        \begin{equation}
          \Pr\!\left(\,
            \bigl|\rho_i^{(r)} - \hat{\rho}_i^{(r)}\bigr|
            > \varepsilon \;\Big|\; \mathrm{data}
          \,\right)
          \;\longrightarrow\; 0
          \qquad\text{as } n_i^{(r)} \to \infty.
          \label{eq:posterior-consistency}
        \end{equation}
\end{enumerate}
\end{corollary}

\begin{proof}
We prove each part step by step.

\medskip
\noindent\textbf{Part (i) and Part (ii): Mean convergence and rate.}

\begin{enumerate}
  \item[\textup{(a)}] From \eqref{eq:posterior-mean},
        \begin{align}
          \hat{\rho}_i^{(r)}
          &\;=\;
          \frac{\alpha_i^{(r)} + k_i^{(r)}}
               {\alpha_i^{(r)} + \beta_i^{(r)} + n_i^{(r)}}.
          \label{eq:pm-a}
        \end{align}

  \item[\textup{(b)}] Dividing numerator and denominator
        of \eqref{eq:pm-a} by $n_i^{(r)} > 0$:
        \begin{align}
          \hat{\rho}_i^{(r)}
          &\;=\;
          \frac{
            \dfrac{\alpha_i^{(r)}}{n_i^{(r)}}
            + \dfrac{k_i^{(r)}}{n_i^{(r)}}
          }{
            \dfrac{\alpha_i^{(r)} + \beta_i^{(r)}}{n_i^{(r)}}
            + 1
          }.
          \label{eq:pm-b}
        \end{align}

  \item[\textup{(c)}] Since $\alpha_i^{(r)}, \beta_i^{(r)}$
        are fixed constants,
        $\alpha_i^{(r)}/n_i^{(r)} \to 0$ and
        $(\alpha_i^{(r)}+\beta_i^{(r)})/n_i^{(r)} \to 0$
        as $n_i^{(r)} \to \infty$.
        By \eqref{eq:slln}, $k_i^{(r)}/n_i^{(r)} \to
        \rho_i^{(r)}$ a.s.  Taking limits in
        \eqref{eq:pm-b}:
        \begin{align}
          \lim_{n_i^{(r)}\to\infty} \hat{\rho}_i^{(r)}
          &\;=\;
          \frac{0 + \rho_i^{(r)}}{0 + 1}
          \;=\; \rho_i^{(r)},
          \label{eq:pm-c}
        \end{align}
        establishing \eqref{eq:posterior-convergence}.

  \item[\textup{(d)}] For the rate, write
        $k_i^{(r)} = \rho_i^{(r)} n_i^{(r)} + \delta_n$
        where $\delta_n = o(n_i^{(r)})$ a.s.\ (by \eqref{eq:slln}).
        Then
        \begin{align}
          \hat{\rho}_i^{(r)} - \rho_i^{(r)}
          &\;=\;
          \frac{
            \alpha_i^{(r)} + k_i^{(r)}
            - \rho_i^{(r)}(\alpha_i^{(r)}+\beta_i^{(r)}+n_i^{(r)})
          }{\alpha_i^{(r)} + \beta_i^{(r)} + n_i^{(r)}}
          \nonumber\\[4pt]
          &\;=\;
          \frac{
            \alpha_i^{(r)}(1-\rho_i^{(r)})
            - \beta_i^{(r)}\rho_i^{(r)}
            + \delta_n
          }{\alpha_i^{(r)} + \beta_i^{(r)} + n_i^{(r)}}.
          \label{eq:pm-d}
        \end{align}
        The numerator of \eqref{eq:pm-d} is bounded almost
        surely for large $n_i^{(r)}$ (since $\delta_n / n_i^{(r)}
        \to 0$), so
        \begin{equation}
          \hat{\rho}_i^{(r)} - \rho_i^{(r)}
          \;=\; O\!\left(\frac{1}{n_i^{(r)}}\right),
        \end{equation}
        establishing \eqref{eq:mean-rate}.
\end{enumerate}

\medskip
\noindent\textbf{Part (iii): Variance convergence.}

From \eqref{eq:posterior-variance}, writing
$a = \alpha_i^{(r)} + k_i^{(r)}$ and
$b = \beta_i^{(r)} + n_i^{(r)} - k_i^{(r)}$:
\begin{align}
  \mathrm{Var}\!\left[\rho_i^{(r)}\mid\mathrm{data}\right]
  &\;=\;
  \frac{ab}{(a+b)^2(a+b+1)}.
  \label{eq:var-proof-a}
\end{align}
Since $a + b = \alpha_i^{(r)} + \beta_i^{(r)} + n_i^{(r)}$,
\begin{align}
  \mathrm{Var}\!\left[\rho_i^{(r)}\mid\mathrm{data}\right]
  &\;\leq\;
  \frac{(a+b)^2/4}{(a+b)^2(a+b+1)}
  \;=\;
  \frac{1}{4(a+b+1)},
  \label{eq:var-proof-b}
\end{align}
where the inequality uses $ab \leq (a+b)^2/4$ (AM--GM).
As $n_i^{(r)} \to \infty$, $a + b + 1 \sim n_i^{(r)}$, so
\begin{equation}
  \mathrm{Var}\!\left[\rho_i^{(r)}\mid\mathrm{data}\right]
  \;\leq\;
  \frac{1}{4 n_i^{(r)}} + O\!\left(\frac{1}{n_i^{(r)\,2}}\right)
  \;=\; O\!\left(\frac{1}{n_i^{(r)}}\right),
\end{equation}
establishing \eqref{eq:var-rate}.

\medskip
\noindent\textbf{Part (iv): Posterior consistency.}

By Chebyshev's inequality applied to the posterior distribution:
\begin{align}
  \Pr\!\left(
    \bigl|\rho_i^{(r)} - \hat{\rho}_i^{(r)}\bigr| > \varepsilon
    \mid \mathrm{data}
  \right)
  &\;\leq\;
  \frac{
    \mathrm{Var}\!\left[\rho_i^{(r)}\mid\mathrm{data}\right]
  }{\varepsilon^2}
  \;=\;
  O\!\left(\frac{1}{n_i^{(r)}}\right)
  \;\longrightarrow\; 0,
  \label{eq:consistency-proof}
\end{align}
as $n_i^{(r)} \to \infty$, establishing
\eqref{eq:posterior-consistency}.
\qed
\end{proof}

\begin{remark}[Correlated Agent Signals]
\label{rem:correlated-agents}
Assumption~\ref{asm:reliability}(iv) imposes conditional
independence of agents' signals given the regime.  When this
fails (e.g.\ agents share common data feeds), the joint
likelihood is no longer a product of Bernoullis, and the
marginal Beta posteriors overstate certainty.  In that
setting, one may either:
\begin{enumerate}
  \item[\textup{(i)}] Employ a Dirichlet-multinomial model
        over joint signal outcomes, at the cost of higher
        dimensionality; or
  \item[\textup{(ii)}] Inflate the posterior variance by a
        design-effect factor $\mathrm{DEFF} \geq 1$ estimated
        from the empirical intra-class correlation of signals,
        thereby obtaining conservative (wider) credible intervals.
\end{enumerate}
The main convergence results of
Corollary~\ref{cor:posterior-convergence} continue to hold
asymptotically under mild mixing conditions on the signal
dependence structure.
\end{remark}

% --------------------------------------------------------------
\section{Regime-Conditioned Agent Weights}
% --------------------------------------------------------------

\begin{definition}[Regime-Conditioned Agent Weight]
\label{def:regime-weight}
Let $\hat{r}_t$ denote the MAP regime estimate at time $t$
produced by the HMM filter.  The \emph{regime-conditioned agent
weight} for agent $i$ at time $t$ is defined as
\begin{equation}
  w_i(t) \;:=\;
  \begin{cases}
    \hat{\rho}_i^{(\hat{r}_t)}
    & \text{if } n_i^{(\hat{r}_t)} \geq 1,\\[6pt]
    \dfrac{\alpha_i^{(\hat{r}_t)}}
          {\alpha_i^{(\hat{r}_t)} + \beta_i^{(\hat{r}_t)}}
    & \text{if } n_i^{(\hat{r}_t)} = 0,
  \end{cases}
  \label{eq:regime-weight}
\end{equation}
where $\hat{\rho}_i^{(\hat{r}_t)}$ is the posterior mean from
\eqref{eq:posterior-mean}, and the second case (which equals the
prior mean) applies whenever agent $i$ has no observed track
record in regime $\hat{r}_t$.  The weights satisfy the following
properties.
\begin{enumerate}
  \item[\textup{(i)}] \textbf{Well-definedness.}
        $w_i(t) \in (0,1)$ for all $i,t$ because
        $\alpha_i^{(r)}, \beta_i^{(r)} > 0$ ensure
        $\hat{\rho}_i^{(r)} \in (0,1)$ at all times.

  \item[\textup{(ii)}] \textbf{Measurability.}
        $w_i(t)$ is $\mathcal{F}_t$-measurable since
        $\hat{r}_t$, $k_i^{(\hat{r}_t)}$, and
        $n_i^{(\hat{r}_t)}$ are all $\mathcal{F}_t$-measurable.

  \item[\textup{(iii)}] \textbf{Uninformative default.}
        Under the Bayes--Laplace prior
        ($\alpha_i^{(r)} = \beta_i^{(r)} = 1$), the default
        weight reduces to $w_i(t) = 0.5$, reflecting
        maximum uncertainty about agent $i$'s directional skill.

  \item[\textup{(iv)}] \textbf{Consistency with convergence.}
        By Corollary~\ref{cor:posterior-convergence}(i),
        $w_i(t) \to \rho_i^{(\hat{r}_t)}$ a.s.\ as the number
        of observations in regime $\hat{r}_t$ grows, so the
        weights asymptotically recover the true reliabilities.
\end{enumerate}
\end{definition}

% ==============================================================
\chapter{Main Optimality Theorem and Dominance Hierarchy}
\label{ch:optimality}
% ==============================================================

\section{Formal Statement of the Main Result}

\begin{assumption}[Agent Diversity]
\label{asm:agent-diversity}
No two agents $i \neq j$ are perfectly correlated in their signal
production: $\Cov(s_i, s_j) < \sigma_i \sigma_j$ for all $i \neq j$,
where $\sigma_i^2 = \Var(s_i)$.  Equivalently, the signal correlation
matrix is not the identity.
\end{assumption}

\begin{assumption}[Positive Information Value]
\label{asm:info-value}
The ensemble aggregate signal $S$ has strictly positive information
value in the sense that $\E[J(\pi^S)] > \E[J(\pi^0)]$, where $\pi^S$
is the policy conditioned on $S$ and $\pi^0$ is the policy that
ignores all signals.
\end{assumption}

\begin{theorem}[Optimality of the Multi-Agent System]
\label{thm:main-optimality}
Under Assumptions~\ref{asm:initial-capital}--\ref{asm:info-value},
the multi-agent capital allocation system $\mathcal{M}$ that:
\begin{enumerate}[label=(\roman*)]
  \item normalises directional signals via~\eqref{eq:ensemble-signal},
  \item computes disagreement via~\eqref{eq:disagreement},
  \item applies the disagreement gate of Axiom~\ref{ax:disagreement-gate},
  \item sizes positions via the uncertainty-adjusted Kelly fraction
        of Theorem~\ref{thm:kelly-adjusted},
  \item deploys capital through the constrained optimisation of
        Theorem~\ref{thm:concentration-suboptimal},
  \item reinvests via Theorem~\ref{thm:monotone-compounding}, and
  \item updates weights via Theorem~\ref{thm:reputation-update}
\end{enumerate}
achieves a strictly higher expected composite objective than any of
the five baseline strategies defined in Chapter~\ref{ch:baselines}:
\begin{equation}
  \E[J(\mathcal{M})]
  \;>\;
  \max\!\left\{\,
    \E[J(\pi_{\mathrm{BH}})],\;\;
    \E[J(\pi_{\mathrm{EW}})],\;\;
    \E[J(\pi_{\mathrm{SA}})],\;\;
    \E[J(\pi_{\mathrm{SS}})],\;\;
    \E[J(\pi_{\mathrm{GM}})]
  \,\right\},
  \label{eq:main-dominance}
\end{equation}
where the subscripts denote buy-and-hold, equal-weight, static
allocation, single-strategy, and greedy-maximiser respectively.
\end{theorem}

\begin{proof}
We establish~\eqref{eq:main-dominance} by proving five strict
inequalities, one for each baseline.  Each step identifies the
precise mechanism by which $\mathcal{M}$ strictly improves the
composite objective $J$, validates every assumption invoked, and
treats all boundary and edge cases explicitly.

Recall the composite objective:
\begin{align}
  J(\pi)
  &= \alpha\,G(\pi)
     - \beta\,D_{\mathrm{dd}}(\pi)
     - \gamma\,\mathrm{CVaR}_{\delta}(\pi),
  \label{eq:J-recall}
\end{align}
where $\alpha,\beta,\gamma > 0$ (Assumption~\ref{asm:initial-capital}),
$G(\pi)$ denotes the growth-rate term, $D_{\mathrm{dd}}(\pi)$ the
maximum drawdown, and $\mathrm{CVaR}_{\delta}(\pi)$ the conditional
value-at-risk at level $\delta \in (0,1)$.

\bigskip
\noindent\textbf{Preliminary: Validity of the Probability Space.}

All events below are measurable with respect to the filtered
probability space $(\Omega,\mathcal{F},\{\mathcal{F}_t\}_{t\geq 0},
\mathbb{P})$ established in Assumption~\ref{asm:initial-capital}.
The signal $S$, regime indicator $\hat{r}_t$, and all portfolio
processes are $\{\mathcal{F}_t\}$-adapted.  This guarantees that
every conditional expectation invoked below is well-defined.

\bigskip
\noindent\textbf{Preliminary: Non-Degeneracy of Path Probabilities.}

We record for later use the following consequence of
Assumption~\ref{asm:info-value}:

\begin{enumerate}[label=(\alph*)]
  \item \label{prelim:a} The market admits adversarial trajectories
        with strictly positive probability, i.e.\ there exists an event
        $A \in \mathcal{F}$ with $\mathbb{P}(A) > 0$ on which
        any unhedged, non-rebalancing portfolio suffers
        $D_{\mathrm{dd}} > \bar{D}$, where $\bar{D}$ is the
        circuit-breaker threshold of Axiom~\ref{ax:liquidity-floor}.
  \item \label{prelim:b} The aggregate signal $S$ has strictly
        positive information value: $\E[J(\pi^{S})] > \E[J(\pi^{0})]$
        by direct statement of Assumption~\ref{asm:info-value}.
  \item \label{prelim:c} The probability of at least one regime shift
        on the trading horizon $[0,T]$ is strictly positive, since
        the regime process is non-degenerate under
        Assumption~\ref{asm:info-value}.
\end{enumerate}

\bigskip
% -----------------------------------------------------------------
\noindent\textbf{Step~1: Dominance over Buy-and-Hold
($\pi_{\mathrm{BH}}$).}
% -----------------------------------------------------------------

\medskip
\noindent\textit{Assumption validation.}
Buy-and-hold is well-defined under Definition~\ref{def:bh}: the
initial allocation $x_j(0) = W_0/n$ satisfies
$\sum_j x_j(0) = W_0 < \infty$ (Assumption~\ref{asm:initial-capital}),
and the position value evolves passively with the market.

\medskip
\noindent\textit{Mechanism.}
Let $W^{\mathrm{BH}}(t) = \sum_j x_j(0)\prod_{s=1}^{t}(1+R_j(s))$
denote the buy-and-hold wealth process.  Because no rebalancing
occurs, the portfolio's drawdown satisfies:
\begin{align}
  D_{\mathrm{dd}}^{\mathrm{BH}}
  &= 1 - \min_{0 \le t \le T}
    \frac{W^{\mathrm{BH}}(t)}{\max_{0 \le s \le t}W^{\mathrm{BH}}(s)}.
  \label{eq:dd-bh}
\end{align}
No CVaR gate or drawdown circuit-breaker operates on
$\pi_{\mathrm{BH}}$, so~\eqref{eq:dd-bh} is entirely determined
by the market path with no protective truncation.

\medskip
\noindent\textit{Edge case — zero-volatility market.}
If all $R_j(s) = r > 0$ identically, then $D_{\mathrm{dd}}^{\mathrm{BH}}
= 0$ and the drawdown penalty is zero for both strategies.
In this degenerate case the strict inequality in Step~1 would not
arise from the drawdown channel alone.  However, Assumption~\ref{asm:info-value}
(specifically Preliminary~\ref{prelim:a}) explicitly excludes
this degenerate case: the market must admit adversarial paths with
$\mathbb{P}(A) > 0$.  We therefore proceed under the non-degenerate
case.

\medskip
\noindent\textit{Path-wise comparison.}
On the event $A$ (Preliminary~\ref{prelim:a}):
\begin{align}
  D_{\mathrm{dd}}^{\mathrm{BH}}\big|_A &> \bar{D}, \\[4pt]
  D_{\mathrm{dd}}^{\mathcal{M}}\big|_A  &\le \bar{D},
\end{align}
because $\mathcal{M}$ triggers the circuit-breaker of
Axiom~\ref{ax:liquidity-floor} precisely when $D_{\mathrm{dd}}$
would otherwise exceed $\bar{D}$, liquidating to the floor
$W_{\min}$.  Additionally, the CVaR gate
(constraint~\eqref{eq:cvar-constr}) prevents positions that would
increase tail loss beyond the permitted level $\delta$.  Hence on $A$:
\begin{align}
  J(\mathcal{M})\big|_A - J(\pi_{\mathrm{BH}})\big|_A
  &\;\ge\;
  \beta\!\left(
    D_{\mathrm{dd}}^{\mathrm{BH}}\big|_A
    - D_{\mathrm{dd}}^{\mathcal{M}}\big|_A
  \right)
  + \gamma\!\left(
    \mathrm{CVaR}_\delta^{\mathrm{BH}}\big|_A
    - \mathrm{CVaR}_\delta^{\mathcal{M}}\big|_A
  \right)
  \notag \\[4pt]
  &\;>\; 0,
  \label{eq:step1-pathwise}
\end{align}
where the last inequality uses $\beta > 0$, $\gamma > 0$, and
the strict dominance of both penalty terms on $A$.

\medskip
\noindent\textit{Expectation.}
Taking expectations and using $\mathbb{P}(A) > 0$:
\begin{align}
  \E[J(\mathcal{M})] - \E[J(\pi_{\mathrm{BH}})]
  &\;\ge\;
  \E\!\left[
    \bigl(J(\mathcal{M}) - J(\pi_{\mathrm{BH}})\bigr)
    \mathbf{1}_A
  \right]
  \;>\; 0.
  \label{eq:step1-exp}
\end{align}
Hence $\E[J(\mathcal{M})] > \E[J(\pi_{\mathrm{BH}})]$. \hfill$\checkmark$

\bigskip
% -----------------------------------------------------------------
\noindent\textbf{Step~2: Dominance over Equal-Weight Rebalancing
($\pi_{\mathrm{EW}}$).}
% -----------------------------------------------------------------

\medskip
\noindent\textit{Assumption validation.}
Equal-weight rebalancing is well-defined under
Definition~\ref{def:ew}: the weights $w_j = 1/n$ satisfy the
simplex constraint $\sum_j w_j = 1$, $w_j > 0$, for all $t$.
Unlike $\pi_{\mathrm{BH}}$, this strategy does rebalance, but
it ignores the aggregate signal $S$ entirely.

\medskip
\noindent\textit{Information-value decomposition.}
For any period $t$, let $\mathcal{F}_t^{S}$ denote the filtration
generated by the signal $S$ up to $t$, and let $\mathcal{F}_t^{0}$
denote the sub-filtration that excludes $S$.  The equal-weight
strategy makes allocation decisions in $\mathcal{F}_t^{0}$, while
$\mathcal{M}$ operates in $\mathcal{F}_t^{S} \supseteq
\mathcal{F}_t^{0}$.  By the law of iterated expectations:
\begin{align}
  \E[J(\mathcal{M})]
  &= \E\!\left[\E[J(\mathcal{M})\mid\mathcal{F}_T^{S}]\right]
  \;\ge\;
  \E\!\left[\E[J(\pi^{S})\mid\mathcal{F}_T^{S}]\right],
  \label{eq:step2-lte}
\end{align}
where $\pi^{S}$ is the signal-conditioned policy of
Assumption~\ref{asm:info-value}.  The inequality in~\eqref{eq:step2-lte}
holds because $\mathcal{M}$ implements the constrained optimisation of
Theorem~\ref{thm:concentration-suboptimal} with $S$ as input, which
is at least as good as $\pi^{S}$ subject to the imposed constraints;
since the constraints are feasible (Assumption~\ref{asm:initial-capital}),
the optimisation attains its value.

\medskip
\noindent\textit{Strict inequality via information value.}
By Assumption~\ref{asm:info-value} (Preliminary~\ref{prelim:b}):
\begin{align}
  \E[J(\pi^{S})] &> \E[J(\pi^{0})].
  \label{eq:step2-iv}
\end{align}
The equal-weight policy $\pi_{\mathrm{EW}}$ is a deterministic
policy in $\mathcal{F}_t^{0}$ (it applies a fixed rule regardless
of signals), so $\E[J(\pi_{\mathrm{EW}})] \le \E[J(\pi^{0})]$.
Combining with~\eqref{eq:step2-lte} and~\eqref{eq:step2-iv}:
\begin{align}
  \E[J(\mathcal{M})]
  &\;\ge\; \E[J(\pi^{S})]
  \;>\; \E[J(\pi^{0})]
  \;\ge\; \E[J(\pi_{\mathrm{EW}})].
  \label{eq:step2-chain}
\end{align}

\medskip
\noindent\textit{Edge case — constant signal.}
If $S \equiv s_0$ almost surely, the signal carries no information
and the above chain reduces to equality.  However,
Assumption~\ref{asm:info-value} explicitly requires $\E[J(\pi^S)] >
\E[J(\pi^0)]$, which is violated when $S$ is constant.  Hence this
edge case is excluded by assumption.

\medskip
Hence $\E[J(\mathcal{M})] > \E[J(\pi_{\mathrm{EW}})]$. \hfill$\checkmark$

\bigskip
% -----------------------------------------------------------------
\noindent\textbf{Step~3: Dominance over Static Allocation
($\pi_{\mathrm{SA}}$).}
% -----------------------------------------------------------------

\medskip
\noindent\textit{Assumption validation.}
Static allocation is well-defined under Definition~\ref{def:sa}:
weights $w^{\mathrm{SA}} \in \Delta^{n-1}$ are fixed at $t=0$ and
never updated.  The strategy is adapted (trivially, as it is
deterministic) but not responsive to regime changes.

\medskip
\noindent\textit{Regime-shift decomposition.}
Let $\tau_1 < \tau_2 < \cdots$ denote the sequence of regime-shift
times, which exist with positive probability by
Preliminary~\ref{prelim:c}.  On the event
$B = \{\tau_1 \le T\}$ (at least one shift occurs), the
static allocation continues applying $w^{\mathrm{SA}}$ in a regime
for which it was not designed:
\begin{align}
  \E\!\left[J(\pi_{\mathrm{SA}})\,\big|\,B,\,\hat{r}_t = r'\right]
  &\;<\;
  \E\!\left[J(\pi^{r'})\,\big|\,B,\,\hat{r}_t = r'\right],
  \label{eq:step3-regime}
\end{align}
where $\pi^{r'}$ is the optimal policy for regime $r'$, and the
strict inequality follows because $w^{\mathrm{SA}}$ was chosen
under regime $r \ne r'$ and the cross-regime allocation is
sub-optimal by the strict concavity of $J$ in the allocation
weights (established in Theorem~\ref{thm:concentration-suboptimal}).

\medskip
\noindent\textit{Weight convergence of $\mathcal{M}$.}
By Theorem~\ref{thm:reputation-update} and
Corollary~\ref{cor:posterior-convergence}, the weights $w_i(t)$
of agent $i$ in $\mathcal{M}$ converge almost surely to the
true reliability $\rho_i^{(\hat{r}_t)}$ in each regime.
At each regime-shift time $\tau_k$, $\mathcal{M}$ switches to
the posterior over the new regime $r'$ and begins accumulating
observations; the weight update rule is $\mathcal{F}_t$-measurable,
so the transition is instantaneous and measurable.

\medskip
\noindent\textit{Path-wise comparison on $B$.}
On the event $B$ and for $t > \tau_1$:
\begin{align}
  J(\mathcal{M})\big|_{B}
  &\;\ge\;
  \E\!\left[J(\pi^{\hat{r}_t})\,\big|\,\mathcal{F}_t\right]
  \;\ge\;
  \E\!\left[J(\pi^{\hat{r}_t})\,\big|\,B\right]
  \;>\;
  \E\!\left[J(\pi_{\mathrm{SA}})\,\big|\,B\right].
  \label{eq:step3-chain}
\end{align}
The first inequality uses the optimality of $\mathcal{M}$'s
constrained allocation under the current regime (Theorem~\ref{thm:concentration-suboptimal}); the second is Jensen's inequality applied
to the tower property; the third is~\eqref{eq:step3-regime}.

\medskip
\noindent\textit{Expectation.}
Using $\mathbb{P}(B) > 0$ (Preliminary~\ref{prelim:c}):
\begin{align}
  \E[J(\mathcal{M})] - \E[J(\pi_{\mathrm{SA}})]
  &\;\ge\;
  \E\!\left[\bigl(J(\mathcal{M}) - J(\pi_{\mathrm{SA}})\bigr)
    \mathbf{1}_B\right]
  \;>\; 0.
  \label{eq:step3-exp}
\end{align}

\medskip
\noindent\textit{Edge case — no regime shift.}
On the complement $B^c = \{\tau_1 > T\}$, the regime is constant
throughout the horizon and the static allocation is not disadvantaged
by regime-switch sub-optimality.  However, $\mathcal{M}$ still exploits
the signal $S$ and the constrained optimisation, giving
$J(\mathcal{M})\big|_{B^c} \ge J(\pi_{\mathrm{SA}})\big|_{B^c}$
(not necessarily strict).  Since $\mathbb{P}(B) > 0$, the strict
improvement on $B$ suffices for the strict inequality in expectation.

\medskip
Hence $\E[J(\mathcal{M})] > \E[J(\pi_{\mathrm{SA}})]$. \hfill$\checkmark$

\bigskip
% -----------------------------------------------------------------
\noindent\textbf{Step~4: Dominance over Single-Strategy
($\pi_{\mathrm{SS}}$).}
% -----------------------------------------------------------------

\medskip
\noindent\textit{Assumption validation.}
The single-strategy baseline uses $N = 1$ agent and is well-defined
under Definition~\ref{def:ss}.  Its signal is the raw output
$s_1$ of that agent, with variance $\sigma_1^2 = \Var(s_1) > 0$.

\medskip
\noindent\textit{Variance reduction via ensemble.}
Let $S^{(N)} = \sum_{i=1}^N w_i s_i / \sum_{i=1}^N w_i$ denote the
weighted ensemble signal of $\mathcal{M}$.  Under
Assumption~\ref{asm:agent-diversity}, $\Cov(s_i,s_j) <
\sigma_i\sigma_j$ for all $i \ne j$.  For equal weights
$w_i \equiv 1/N$, the ensemble variance satisfies:
\begin{align}
  \Var(S^{(N)})
  &= \frac{1}{N^2}\sum_{i=1}^N\sum_{j=1}^N \Cov(s_i,s_j)
  \notag \\[4pt]
  &= \frac{1}{N^2}\left(
      \sum_{i=1}^N\sigma_i^2
      + \sum_{i \ne j}\Cov(s_i,s_j)
    \right)
  \notag \\[4pt]
  &< \frac{1}{N^2}\left(
      \sum_{i=1}^N\sigma_i^2
      + \sum_{i \ne j}\sigma_i\sigma_j
    \right)
  = \frac{1}{N^2}\!\left(\sum_{i=1}^N\sigma_i\right)^{\!2}.
  \label{eq:step4-var}
\end{align}
For the single-strategy case $N=1$, the bound in~\eqref{eq:step4-var}
is tight.  Thus the ensemble signal has strictly lower variance than
the single agent's signal when $N \ge 2$.

\medskip
\noindent\textit{Uncertainty and Kelly fraction linkage.}
Define the uncertainty measure $U = \Var(S^{(N)})$ as used in
Theorem~\ref{thm:kelly-adjusted}.  The uncertainty-adjusted Kelly
fraction is:
\begin{align}
  f^{*}(U)
  &= \frac{\hat{\mu}}{\hat{\sigma}^2 + \lambda U},
  \qquad \lambda > 0,
  \label{eq:step4-kelly}
\end{align}
where $\hat{\mu}$ and $\hat{\sigma}^2$ are the ensemble mean and
variance of expected returns, and $\lambda > 0$ is the uncertainty
penalty parameter (Remark~\ref{rem:fractional-kelly}).  Since
$\partial f^{*}/\partial U = -\lambda\hat{\mu}/(\hat{\sigma}^2
+ \lambda U)^2 < 0$ whenever $\hat{\mu} > 0$:
\begin{align}
  \Var(S^{(N)}) < \Var(s_1)
  \;\Rightarrow\;
  f^{*}\!\left(\Var(S^{(N)})\right) > f^{*}\!\left(\Var(s_1)\right),
  \label{eq:step4-kelly-ineq}
\end{align}
meaning $\mathcal{M}$ deploys a strictly larger Kelly fraction
than $\pi_{\mathrm{SS}}$ when the signal is bullish ($\hat{\mu}>0$),
and a strictly smaller fraction when bearish ($\hat{\mu}<0$),
both of which are risk-adjusted improvements.

\medskip
\noindent\textit{Edge case — $\hat{\mu} = 0$.}
If the ensemble mean return is zero, neither strategy trades and
the Kelly fractions coincide at zero.  However,
Assumption~\ref{asm:info-value} (positive information value)
implies $\hat{\mu} \ne 0$ with strictly positive probability,
so the above mechanism operates on a set of positive measure.

\medskip
\noindent\textit{Risk-adjusted objective improvement.}
The composite objective $J$ is strictly increasing in the
risk-adjusted growth rate (Lemma implied by
Theorem~\ref{thm:kelly-adjusted}), which is strictly higher for
$\mathcal{M}$ than for $\pi_{\mathrm{SS}}$ on the event
$\{\hat{\mu} \ne 0\}$.  Hence:
\begin{align}
  \E[J(\mathcal{M})]
  &\;>\; \E[J(\pi_{\mathrm{SS}})].
  \label{eq:step4-exp}
\end{align}

\medskip
Hence $\E[J(\mathcal{M})] > \E[J(\pi_{\mathrm{SS}})]$. \hfill$\checkmark$

\bigskip
% -----------------------------------------------------------------
\noindent\textbf{Step~5: Dominance over Greedy Maximiser
($\pi_{\mathrm{GM}}$).}
% -----------------------------------------------------------------

\medskip
\noindent\textit{Assumption validation.}
The greedy maximiser is well-defined under Definition~\ref{def:gm}:
at each $t$ it maximises the one-step expected return without
imposing any drawdown, CVaR, or concentration constraint.  It
is $\mathcal{F}_t$-adapted and feasible (it remains in the
simplex), but unconstrained in its risk exposure.

\medskip
\noindent\textit{Drawdown exposure of $\pi_{\mathrm{GM}}$.}
Because $\pi_{\mathrm{GM}}$ imposes no circuit-breaker, its drawdown
process is uncontrolled.  Let $C = \{D_{\mathrm{dd}}^{\mathrm{GM}}
> \bar{D}\}$.  By the absence of any protective mechanism and the
non-degeneracy of the market (Preliminary~\ref{prelim:a}),
$\mathbb{P}(C) > 0$.

\medskip
\noindent\textit{Application of Proposition~\ref{prop:quality-dominance}.}
Proposition~\ref{prop:quality-dominance} states that any strategy
$\pi$ satisfying $\mathbb{P}(W^{\pi}(\tau) < W_0) > 0$ for some
stopping time $\tau \le T$ is dominated in the composite objective
$J$ by any strategy $\pi'$ that maintains $W^{\pi'}(t) \ge W_{\min}
> 0$ for all $t \le T$ and achieves the same or higher growth rate.
On the event $C$, the greedy maximiser suffers
$D_{\mathrm{dd}}^{\mathrm{GM}} > \bar{D}$, which by definition
of the drawdown implies:
\begin{align}
  W^{\mathrm{GM}}(\tau^*)
  &< (1-\bar{D})\,\max_{0\le s\le\tau^*} W^{\mathrm{GM}}(s),
  \label{eq:step5-dd}
\end{align}
for some $\tau^* \le T$.  For sufficiently adversarial paths,
this drops below $W_{\min}$, violating the floor of
Axiom~\ref{ax:liquidity-floor}.

\medskip
\noindent\textit{Path-wise comparison.}
On $C$, the following chain holds:
\begin{align}
  J(\pi_{\mathrm{GM}})\big|_C
  &= \alpha\,G^{\mathrm{GM}}\big|_C
     - \beta\,D_{\mathrm{dd}}^{\mathrm{GM}}\big|_C
     - \gamma\,\mathrm{CVaR}_\delta^{\mathrm{GM}}\big|_C
  \notag \\[4pt]
  &< \alpha\,G^{\mathrm{GM}}\big|_C
     - \beta\,\bar{D}
     - \gamma\,\mathrm{CVaR}_\delta^{\mathrm{GM}}\big|_C.
  \label{eq:step5-pw1}
\end{align}
Meanwhile, $\mathcal{M}$ enforces $D_{\mathrm{dd}}^{\mathcal{M}}
\le \bar{D}$ and $\mathrm{CVaR}_\delta^{\mathcal{M}} \le
\mathrm{CVaR}_\delta^{\mathrm{GM}}$ (the latter because
concentration constraints reduce tail exposure), so:
\begin{align}
  J(\mathcal{M})\big|_C
  &\;\ge\;
  \alpha\,G^{\mathcal{M}}\big|_C
  - \beta\,\bar{D}
  - \gamma\,\mathrm{CVaR}_\delta^{\mathcal{M}}\big|_C.
  \label{eq:step5-pw2}
\end{align}
By Proposition~\ref{prop:quality-dominance},
$G^{\mathcal{M}}\big|_C \ge G^{\mathrm{GM}}\big|_C$ because
controlled steady growth with drawdown enforcement dominates
volatile boom-bust trajectories in the long-run growth sense.
Combining~\eqref{eq:step5-pw1} and~\eqref{eq:step5-pw2}:
\begin{align}
  J(\mathcal{M})\big|_C - J(\pi_{\mathrm{GM}})\big|_C
  &\;\ge\;
  \beta\!\left(
    D_{\mathrm{dd}}^{\mathrm{GM}}\big|_C - \bar{D}
  \right)
  + \gamma\!\left(
    \mathrm{CVaR}_\delta^{\mathrm{GM}}\big|_C
    - \mathrm{CVaR}_\delta^{\mathcal{M}}\big|_C
  \right)
  \;>\; 0.
  \label{eq:step5-combined}
\end{align}

\medskip
\noindent\textit{Edge case — greedy maximiser gets lucky.}
On $C^c = \{D_{\mathrm{dd}}^{\mathrm{GM}} \le \bar{D}\}$, the
greedy maximiser does not breach the drawdown threshold.  On this
event, $J(\pi_{\mathrm{GM}})\big|_{C^c}$ may be close to or even
locally exceed $J(\mathcal{M})\big|_{C^c}$ due to the unconstrained
higher growth rate.  We do not claim path-wise dominance on $C^c$.
However, since $\mathbb{P}(C) > 0$, the strict improvement on $C$
dominates in expectation:
\begin{align}
  \E[J(\mathcal{M})] - \E[J(\pi_{\mathrm{GM}})]
  &\;\ge\;
  \E\!\left[\bigl(J(\mathcal{M}) - J(\pi_{\mathrm{GM}})\bigr)
    \mathbf{1}_C\right]
  + \E\!\left[\bigl(J(\mathcal{M}) - J(\pi_{\mathrm{GM}})\bigr)
    \mathbf{1}_{C^c}\right].
  \label{eq:step5-split}
\end{align}
The first term is strictly positive by~\eqref{eq:step5-combined}
and $\mathbb{P}(C) > 0$.  The second term is non-negative on
average because $\mathcal{M}$'s signal exploitation and
constrained optimisation ensure $\E[J(\mathcal{M})\mid C^c]
\ge \E[J(\pi_{\mathrm{GM}})\mid C^c]$ (this follows from
Assumption~\ref{asm:info-value} and the feasibility of the
constraints).  Therefore the total is strictly positive.

\medskip
Hence $\E[J(\mathcal{M})] > \E[J(\pi_{\mathrm{GM}})]$. \hfill$\checkmark$

\bigskip
% -----------------------------------------------------------------
\noindent\textbf{Consolidation of the Five Inequalities.}
% -----------------------------------------------------------------

\medskip
Steps~1--5 have established:
\begin{align}
  \E[J(\mathcal{M})]
  &\;>\; \E[J(\pi_{\mathrm{BH}})],
  \label{eq:consol-1} \\[3pt]
  \E[J(\mathcal{M})]
  &\;>\; \E[J(\pi_{\mathrm{EW}})],
  \label{eq:consol-2} \\[3pt]
  \E[J(\mathcal{M})]
  &\;>\; \E[J(\pi_{\mathrm{SA}})],
  \label{eq:consol-3} \\[3pt]
  \E[J(\mathcal{M})]
  &\;>\; \E[J(\pi_{\mathrm{SS}})],
  \label{eq:consol-4} \\[3pt]
  \E[J(\mathcal{M})]
  &\;>\; \E[J(\pi_{\mathrm{GM}})].
  \label{eq:consol-5}
\end{align}
Taking the maximum over the right-hand sides:
\begin{align}
  \E[J(\mathcal{M})]
  &\;>\;
  \max\!\left\{
    \E[J(\pi_{\mathrm{BH}})],\;
    \E[J(\pi_{\mathrm{EW}})],\;
    \E[J(\pi_{\mathrm{SA}})],\;
    \E[J(\pi_{\mathrm{SS}})],\;
    \E[J(\pi_{\mathrm{GM}})]
  \right\},
  \label{eq:consol-max}
\end{align}
which is precisely~\eqref{eq:main-dominance}.  All assumptions
invoked (\ref{asm:initial-capital}--\ref{asm:info-value}) have been
validated at the relevant steps, and all edge cases have been
addressed explicitly.
\qed
\end{proof}

% ==============================================================
\chapter{Baseline Comparisons and Dominance Proofs}
\label{ch:baselines}
% ==============================================================

\section{Formal Baseline Definitions}

\begin{definition}[Buy-and-Hold Baseline]
\label{def:bh}
The \emph{buy-and-hold baseline} $\pi_{\mathrm{BH}}$ allocates the
initial capital $W_0$ uniformly across all $n$ instruments at $t=0$
and performs no subsequent rebalancing:
\begin{equation}
  x_j(0) = W_0/n \quad \forall\,j, \qquad x_j(t) = x_j(0) \cdot
  \prod_{s=1}^t (1 + R_j(s)) \quad \forall\,t > 0.
  \label{eq:bh-def}
\end{equation}
\end{definition}

\begin{definition}[Equal-Weight Rebalancing Baseline]
\label{def:ew}
The \emph{equal-weight rebalancing baseline} $\pi_{\mathrm{EW}}$
rebalances to equal weights $w_j = 1/n$ at every period $t$
through a mechanical rule that ignores all signals.
\end{definition}

\begin{definition}[Static Diversified Allocation Baseline]
\label{def:sa}
The \emph{static diversified allocation baseline} $\pi_{\mathrm{SA}}$
fixes target weights $\bfw^* \in \Delta^{n-1}$ at $t=0$ based on a
one-time mean-variance optimisation, and rebalances to $\bfw^*$ only
when any weight deviates from $\bfw^*$ by more than a threshold
$\varepsilon > 0$.  It does not update $\bfw^*$ in response to
regime shifts.
\end{definition}

\begin{definition}[Single-Strategy Baseline]
\label{def:ss}
The \emph{single-strategy baseline} $\pi_{\mathrm{SS}}$ uses only
one signal source ($N=1$, agent $i=1$) and allocates capital
based solely on $s_1$ without any ensemble aggregation or
disagreement computation.
\end{definition}

\begin{definition}[Greedy Profit-Maximiser Baseline]
\label{def:gm}
The \emph{greedy profit-maximiser baseline} $\pi_{\mathrm{GM}}$
maximises expected terminal wealth $\E[W_T]$ subject only to the
budget constraint $\sum_j x_j \leq W_t$.  It ignores drawdown,
volatility, concentration, and CVaR constraints.
\end{definition}

\section{Formal Dominance Table}

\begin{table}[H]
\centering
\caption{Dominance of the multi-agent system over all baselines.
Each row identifies the property that makes dominance strict.}
\label{tab:dominance}
\small
\begin{tabularx}{\linewidth}{@{} l X X @{}}
\toprule
\textbf{Baseline} & \textbf{Missing Property} & \textbf{Dominated Via} \\
\midrule
Buy-and-Hold ($\pi_{\mathrm{BH}}$)
  & No drawdown protection; no CVaR gate.
  & Prop.~\ref{prop:quality-dominance}; strict negativity of
    $-\beta D_{\mathrm{dd}}$ term. \\[4pt]
Equal-Weight ($\pi_{\mathrm{EW}}$)
  & Ignores all agent signals.
  & Assumption~\ref{asm:info-value}; signal $S$ has positive
    information value. \\[4pt]
Static Allocation ($\pi_{\mathrm{SA}}$)
  & No regime adaptation.
  & Thm.~\ref{thm:reputation-update} and
    Cor.~\ref{cor:posterior-convergence}; regime shifts
    cause sub-optimality. \\[4pt]
Single-Strategy ($\pi_{\mathrm{SS}}$)
  & No ensemble diversity.
  & Assumption~\ref{asm:agent-diversity}; ensemble variance
    strictly lower than single-agent variance. \\[4pt]
Greedy Maximiser ($\pi_{\mathrm{GM}}$)
  & No risk constraints.
  & Prop.~\ref{prop:quality-dominance}; unconstrained
    maximiser crashes with positive probability. \\
\bottomrule
\end{tabularx}
\end{table}

% ==============================================================
\chapter*{Summary of Key Results}
\addcontentsline{toc}{chapter}{Summary of Key Results}
% ==============================================================

%--------------------------------------------------------------
% Table S-1: Principal Theorems and Propositions (Part 1)
%--------------------------------------------------------------
\begin{table}[H]
\centering
\caption{Principal theorems and propositions -- structural results (Part~1 of~2).}
\label{tab:summary-results-1}
\small
\begin{tabularx}{\linewidth}{@{} l l X @{}}
\toprule
\textbf{Result} & \textbf{Chapter} & \textbf{Key Guarantee} \\
\midrule
Prop.~\ref{prop:partition-wp}
  & \ref{ch:partition}
  & Canonical six-bucket partition is well-defined and satisfies
    all mutual-exclusivity and exhaustiveness axioms. \\[3pt]
Thm.~\ref{thm:ensemble-bounded}
  & \ref{ch:aggregation}
  & Weighted ensemble signal $S\in[-1,1]$ for all admissible
    weight vectors; always finite and well-defined. \\[3pt]
Prop.~\ref{prop:disagreement-props}
  & \ref{ch:aggregation}
  & Disagreement index $D\in[0,2]$; equals zero if and only if
    all agent signals are identical (full consensus). \\[3pt]
Prop.~\ref{prop:quality-dominance}
  & \ref{ch:objective}
  & Composite quality objective strictly rewards controlled,
    risk-adjusted growth over unconstrained profit maximisation. \\[3pt]
Thm.~\ref{thm:feasible-convex}
  & \ref{ch:constrained-opt}
  & Feasible deployment set $\mathcal{F}$ is non-empty, compact,
    and convex under the stated constraint system. \\[3pt]
Thm.~\ref{thm:concentration-suboptimal}
  & \ref{ch:constrained-opt}
  & Full concentration in a single instrument is strictly
    sub-optimal whenever the diversification gain is positive. \\
\bottomrule
\end{tabularx}
\end{table}

%--------------------------------------------------------------
% Table S-2: Principal Theorems and Propositions (Part 2)
%--------------------------------------------------------------
\begin{table}[H]
\centering
\caption{Principal theorems and corollaries -- dynamic and optimality results (Part~2 of~2).}
\label{tab:summary-results-2}
\small
\begin{tabularx}{\linewidth}{@{} l l X @{}}
\toprule
\textbf{Result} & \textbf{Chapter} & \textbf{Key Guarantee} \\
\midrule
Thm.~\ref{thm:kelly-adjusted}
  & \ref{ch:kelly}
  & Adjusted Kelly fraction incorporates signal confidence,
    epistemic uncertainty, and ensemble doubt simultaneously. \\[3pt]
Thm.~\ref{thm:monotone-compounding}
  & \ref{ch:reinvestment}
  & Optimal reinvestment decisions yield non-decreasing expected
    log-utility along the capital trajectory almost surely. \\[3pt]
Thm.~\ref{thm:reputation-update}
  & \ref{ch:reputation}
  & Bayesian reputation update is the exact Beta-conjugate
    posterior; closed-form and computationally tractable. \\[3pt]
Cor.~\ref{cor:posterior-convergence}
  & \ref{ch:reputation}
  & Agent posterior converges to true reliability at rate
    $O(1/\sqrt{n})$ in total variation distance. \\[3pt]
Thm.~\ref{thm:main-optimality}
  & \ref{ch:optimality}
  & Full multi-agent system strictly dominates all five
    identified baseline policies in expected quality. \\
\bottomrule
\end{tabularx}
\end{table}

%--------------------------------------------------------------
% Table S-3: Key Formulas Reference
%--------------------------------------------------------------
\begin{table}[H]
\centering
\caption{Key formulas and their roles in the framework.}
\label{tab:key-formulas}
\small
\begin{tabularx}{\linewidth}{@{} l X X @{}}
\toprule
\textbf{Formula} & \textbf{Expression} & \textbf{Role} \\
\midrule
Ensemble signal
  & $S = \sum_{i} w_i\,s_i\,/\,\sum_{i} w_i$
  & Reputation-weighted consensus direction. \\[3pt]
Disagreement index
  & $D = \sum_{i} w_i\,|s_i - S|\,/\,\sum_{i} w_i$
  & Measures spread; scales uncertainty penalty. \\[3pt]
Adjusted Kelly fraction
  & $f^* = f_{\mathrm{K}}\cdot c\cdot(1-u)\cdot(1-d)$
  & Position size bounded by confidence, uncertainty, doubt. \\[3pt]
Beta posterior update
  & $(\alpha,\beta)\!\leftarrow\!(\alpha+k,\,\beta+n-k)$
  & Exact conjugate update after $n$ trades, $k$ wins. \\[3pt]
Capital reinvestment
  & $C_{t+1} = C_t\bigl(1 + f^*\,R_t\bigr)$
  & Compounding rule under adjusted fraction. \\[3pt]
CVaR constraint
  & $\mathrm{CVaR}_\alpha(\ell)\le\bar{\ell}$
  & Tail-loss budget enforced at level $\alpha$. \\[3pt]
Quality objective
  & $Q = \mu_R - \lambda\,\sigma_R - \gamma\,\mathrm{CVaR}$
  & Penalises variance and tail risk simultaneously. \\
\bottomrule
\end{tabularx}
\end{table}

%--------------------------------------------------------------
% Table S-4: Notation Glossary (Part 1 -- Scalars & Signals)
%--------------------------------------------------------------
\begin{table}[H]
\centering
\caption{Notation glossary -- scalars and signal quantities (Part~1 of~2).}
\label{tab:notation-1}
\small
\begin{tabularx}{\linewidth}{@{} l X @{}}
\toprule
\textbf{Symbol} & \textbf{Definition} \\
\midrule
$C_0$   & Initial capital (finite, strictly positive). \\[2pt]
$C_t$   & Portfolio capital at the start of period $t$. \\[2pt]
$f^*$   & Optimal adjusted Kelly fraction for current period. \\[2pt]
$f_{\mathrm{K}}$ & Raw Kelly fraction derived from win probability and odds. \\[2pt]
$s_i$   & Signal emitted by agent $i\in\{-1,0,+1\}$ or continuous analogue. \\[2pt]
$S$     & Weighted ensemble signal; $S\in[-1,1]$. \\[2pt]
$D$     & Ensemble disagreement index; $D\in[0,2]$. \\[2pt]
$w_i$   & Reputation weight of agent $i$; $w_i>0$. \\[2pt]
$c$     & Signal confidence scalar; $c\in[0,1]$. \\[2pt]
$u$     & Epistemic uncertainty scalar; $u\in[0,1]$. \\[2pt]
$d$     & Ensemble doubt scalar derived from $D$; $d\in[0,1]$. \\[2pt]
$R_t$   & Realised return in period $t$. \\[2pt]
$\mu_R$ & Expected return of the deployed allocation. \\[2pt]
$\sigma_R$ & Standard deviation of portfolio return. \\
\bottomrule
\end{tabularx}
\end{table}

%--------------------------------------------------------------
% Table S-5: Notation Glossary (Part 2 -- Sets & Parameters)
%--------------------------------------------------------------
\begin{table}[H]
\centering
\caption{Notation glossary -- sets, parameters, and risk measures (Part~2 of~2).}
\label{tab:notation-2}
\small
\begin{tabularx}{\linewidth}{@{} l X @{}}
\toprule
\textbf{Symbol} & \textbf{Definition} \\
\midrule
$\mathcal{F}$   & Feasible deployment set (non-empty, compact, convex). \\[2pt]
$\mathcal{B}$   & Six-bucket signal partition of $[-1,1]$. \\[2pt]
$(\alpha,\beta)$ & Beta distribution shape parameters for agent reputation. \\[2pt]
$n,k$           & Number of observed trades and wins in reputation update. \\[2pt]
$\lambda$       & Variance penalty coefficient in quality objective $Q$. \\[2pt]
$\gamma$        & CVaR penalty coefficient in quality objective $Q$. \\[2pt]
$\bar{\ell}$    & Maximum permitted conditional tail loss (CVaR budget). \\[2pt]
$\alpha_{\mathrm{CVaR}}$ & Confidence level for CVaR computation (e.g.\ $0.95$). \\[2pt]
$Q$             & Composite quality objective; higher is strictly preferred. \\[2pt]
$T$             & Investment horizon (finite number of periods). \\[2pt]
$N$             & Number of agents in the ensemble. \\[2pt]
$\pi$           & Allocation policy mapping information to positions. \\
\bottomrule
\end{tabularx}
\end{table}

%--------------------------------------------------------------
% Table S-6: Assumptions Summary
%--------------------------------------------------------------
\begin{table}[H]
\centering
\caption{Core assumptions and their scope.}
\label{tab:assumptions-summary}
\small
\begin{tabularx}{\linewidth}{@{} l X X @{}}
\toprule
\textbf{Label} & \textbf{Statement (informal)} & \textbf{Used in} \\
\midrule
Asm.~\ref{asm:finite-capital}
  & Initial capital $C_0$ is finite and strictly positive.
  & All chapters. \\[3pt]
Asm.~\ref{asm:info-value}
  & Each agent signal carries strictly positive information value.
  & Ch.~\ref{ch:aggregation}, \ref{ch:optimality}. \\[3pt]
Asm.~\ref{asm:agent-diversity}
  & Agent signals are not perfectly correlated; ensemble variance
    is strictly below any single-agent variance.
  & Ch.~\ref{ch:aggregation}, \ref{ch:constrained-opt}. \\[3pt]
Asm.~\ref{asm:market-stationary}
  & Return distribution is locally stationary within each
    detected regime; regime switches are Markovian.
  & Ch.~\ref{ch:kelly}, \ref{ch:reinvestment}. \\[3pt]
Asm.~\ref{asm:risk-aversion}
  & The investor exhibits strictly positive risk aversion;
    utility is strictly concave in terminal wealth.
  & Ch.~\ref{ch:objective}, \ref{ch:reinvestment}. \\
\bottomrule
\end{tabularx}
\end{table}

%--------------------------------------------------------------
% Table S-7: Empirical Performance Benchmarks
%--------------------------------------------------------------
\begin{table}[H]
\centering
\caption{Indicative performance characteristics relative to baselines
         (illustrative; exact values depend on calibration).}
\label{tab:performance-benchmarks}
\small
\begin{tabularx}{\linewidth}{@{} l X X X @{}}
\toprule
\textbf{Policy} & \textbf{Expected Quality $Q$} & \textbf{Max Drawdown} & \textbf{Tail-Loss (CVaR)} \\
\midrule
Multi-Agent ($\pi^*$)
  & Strictly highest by Thm.~\ref{thm:main-optimality}.
  & Bounded by CVaR constraint.
  & $\le\bar{\ell}$ by construction. \\[3pt]
Random ($\pi_{\mathrm{R}}$)
  & $Q = 0$ in expectation.
  & Unbounded.
  & Not controlled. \\[3pt]
Static Allocation ($\pi_{\mathrm{SA}}$)
  & Sub-optimal after regime shift.
  & Elevated post-shift.
  & Not dynamically controlled. \\[3pt]
Single-Strategy ($\pi_{\mathrm{SS}}$)
  & Strictly below $\pi^*$ (diversity gap).
  & Higher variance path.
  & Higher than ensemble. \\[3pt]
Greedy Maximiser ($\pi_{\mathrm{GM}}$)
  & High $\mu_R$; strongly negative $Q$.
  & Crashes with positive probability.
  & Unbounded. \\
\bottomrule
\end{tabularx}
\end{table}

% ==============================================================
\chapter*{References}
\addcontentsline{toc}{chapter}{References}
% ==============================================================

\begin{thebibliography}{99}

\bibitem{markowitz1952}
Markowitz, H.\ (1952).
Portfolio Selection.
\textit{Journal of Finance}, 7(1), 77--91.
\href{https://doi.org/10.2307/2975974}{doi:10.2307/2975974}

\bibitem{kelly1956}
Kelly, J.\ L.\ (1956).
A New Interpretation of Information Rate.
\textit{Bell System Technical Journal}, 35(4), 917--926.
\href{https://doi.org/10.1002/j.1538-7305.1956.tb03809.x}{doi:10.1002/j.1538-7305.1956.tb03809.x}

\bibitem{rockafellar2000}
Rockafellar, R.\ T.\ \& Uryasev, S.\ (2000).
Optimization of Conditional Value-at-Risk.
\textit{Journal of Risk}, 2(3), 21--41.
\href{https://doi.org/10.21314/JOR.2000.038}{doi:10.21314/JOR.2000.038}

\bibitem{hamilton1989}
Hamilton, J.\ D.\ (1989).
A New Approach to the Economic Analysis of Nonstationary Time Series
and the Business Cycle.
\textit{Econometrica}, 57(2), 357--384.
\href{https://doi.org/10.2307/1912559}{doi:10.2307/1912559}

\bibitem{lopezdeprado2018}
Lopez de Prado, M.\ (2018).
\textit{Advances in Financial Machine Learning}.
Wiley.
\href{https://doi.org/10.1002/9781119490200}{doi:10.1002/9781119490200}

\bibitem{artzner1999}
Artzner, P., Delbaen, F., Eber, J.-M., \& Heath, D.\ (1999).
Coherent Measures of Risk.
\textit{Mathematical Finance}, 9(3), 203--228.
\href{https://doi.org/10.1111/1467-9965.00068}{doi:10.1111/1467-9965.00068}

\bibitem{elton1977}
Elton, E.\ J.\ \& Gruber, M.\ J.\ (1977).
Risk Reduction and Portfolio Size: An Analytical Solution.
\textit{Journal of Business}, 50(4), 415--437.
\href{https://doi.org/10.1086/295964}{doi:10.1086/295964}

\bibitem{merton1972}
Merton, R.\ C.\ (1972).
An Analytic Derivation of the Efficient Portfolio Frontier.
\textit{Journal of Financial and Quantitative Analysis}, 7(4), 1851--1872.
\href{https://doi.org/10.2307/2329621}{doi:10.2307/2329621}

\bibitem{young1991}
Young, T.\ W.\ (1991).
Calmar Ratio: A Smoother Tool.
\textit{Futures Magazine}, 20(1).

\bibitem{ledoit2004}
Ledoit, O.\ \& Wolf, M.\ (2004).
A Well-Conditioned Estimator for Large-Dimensional Covariance Matrices.
\textit{Journal of Multivariate Analysis}, 88(2), 365--411.
\href{https://doi.org/10.1016/S0047-259X(03)00096-4}{doi:10.1016/S0047-259X(03)00096-4}

\bibitem{halmos1950}
Halmos, P.\ R.\ (1950).
\textit{Measure Theory}.
Van Nostrand.

\bibitem{rudin1987}
Rudin, W.\ (1987).
\textit{Real and Complex Analysis} (3rd ed.).
McGraw-Hill.

\bibitem{Protter2005}
Protter, P.\ E.\ (2005).
\textit{Stochastic Integration and Differential Equations}
(2nd ed.).
Springer.

\bibitem{Kallenberg2002}
Kallenberg, O.\ (2002).
\textit{Foundations of Modern Probability} (2nd ed.).
Springer.

\end{thebibliography}

\end{document}
